\documentclass[a4paper,12pt,fleqn,openany,twoside,fontset=windows]{ctexbook}

\usepackage[fleqn]{amsmath}
\usepackage{amssymb,amsthm,mathrsfs}
\usepackage{geometry}
\usepackage[dvipsnames]{xcolor}
\usepackage{graphicx}
\usepackage{tikz}
\usetikzlibrary{calc}
\usepackage[most]{tcolorbox}
\usepackage{fancyhdr}
\usepackage{titlesec}
\usepackage{tocloft}
\usepackage{enumitem}
\usepackage{microtype}
\usepackage{pifont}
\usepackage{hyperref}
\usepackage{romannum}
\usepackage{mathtools}

\definecolor{bookblue}{HTML}{264653}
\definecolor{bookteal}{HTML}{4F7C82}
\definecolor{booklight}{HTML}{EDF4F3}
\definecolor{bookgray}{HTML}{5E686B}

\geometry{
    inner=2.7cm,
    outer=2.5cm,
    top=2.7cm,
    bottom=2.7cm,
    headheight=18pt,
    headsep=18pt,
    footskip=25pt
}
\setlength{\mathindent}{0pt}
\setlength{\parindent}{5pt}
\setlength{\parskip}{3pt plus 1pt minus 1pt}
\linespread{1.25}
\allowdisplaybreaks[3]
\AtBeginDocument{%
    \setlength{\parindent}{5pt}
    \setlength{\abovedisplayskip}{8pt plus 2pt minus 2pt}
    \setlength{\belowdisplayskip}{8pt plus 2pt minus 2pt}
    \setlength{\abovedisplayshortskip}{5pt plus 1pt minus 1pt}
    \setlength{\belowdisplayshortskip}{5pt plus 1pt minus 1pt}
}

\hypersetup{
    unicode=true,
    colorlinks=true,
    linkcolor=bookblue,
    citecolor=bookblue,
    urlcolor=bookteal,
    bookmarksopen=true,
    pdfauthor={Chen},
    pdftitle={数学分析习题整理}
}

\renewcommand{\chaptermark}[1]{%
    \edef\@tmp{第\thechapter 章\quad #1}%
    \expandafter\markboth\expandafter{\@tmp}{}%
}
\fancypagestyle{main}{%
    \fancyhf{}
    \fancyhead[LE]{\small\color{bookgray}\nouppercase{\leftmark}}
    \fancyhead[RO]{\small\color{bookgray}数学分析习题整理}
    \fancyfoot[C]{\color{bookblue}\thepage}
    \fancyfoot[R]{\small\bfseries Chen\;\; github.com/mathlover-1/math-analysis}
    \renewcommand{\headrulewidth}{0.4pt}
    \renewcommand{\footrulewidth}{0pt}
}
\fancypagestyle{plain}{%
    \fancyhf{}
    \fancyfoot[C]{\color{bookblue}\thepage}
    \fancyfoot[R]{\small\bfseries Chen\;\; github.com/mathlover-1/math-analysis}
    \renewcommand{\headrulewidth}{0pt}
    \renewcommand{\footrulewidth}{0pt}
}
\pagestyle{main}

\titleformat{\chapter}[display]
  {\normalfont\sffamily\bfseries\color{bookblue}}
  {\raggedleft\fontsize{46}{50}\selectfont\thechapter}
  {-18pt}
  {\titlerule[1.2pt]\vspace{10pt}\filright\Huge}
\titlespacing*{\chapter}{0pt}{-15pt}{36pt}

\renewcommand{\contentsname}{目\quad 录}
\renewcommand{\cfttoctitlefont}{\hfill\sffamily\bfseries\Huge\color{bookblue}}
\renewcommand{\cftaftertoctitle}{\hfill}
\setlength{\cftbeforetoctitleskip}{5pt}
\setlength{\cftaftertoctitleskip}{28pt}
\setlength{\cftbeforechapskip}{12pt}
\renewcommand{\cftchapfont}{\sffamily\bfseries\large\color{bookblue}}
\renewcommand{\cftchappagefont}{\sffamily\bfseries\large\color{bookblue}}
\renewcommand{\cftchapleader}{\color{bookteal}\cftdotfill{2.5}}

\newtcolorbox{ti}{
    enhanced,
    breakable,
    colback=booklight!55!white,
    colframe=bookteal!35!white,
    boxrule=0.45pt,
    boxsep=0pt,
    left=10pt,
    right=10pt,
    top=9pt,
    bottom=9pt,
    before skip=14pt,
    after skip=14pt,
    arc=2pt,
    outer arc=2pt,
    before upper={\parindent=0pt}
}

\newenvironment{booknotebody}{%
    \begin{center}
    \begin{minipage}{0.84\textwidth}
    \songti\zihao{-4}\linespread{1.45}\selectfont
    \setlength{\parindent}{2\ccwd}
    \setlength{\parskip}{0.75em}
}{%
    \end{minipage}
    \end{center}
}

\newenvironment{booknotelongbody}{%
    \begingroup
    \songti\zihao{-4}\linespread{1.45}\selectfont
    \setlength{\parindent}{2\ccwd}
    \setlength{\parskip}{0.75em}
    \begin{list}{}{%
        \setlength{\leftmargin}{0.08\textwidth}
        \setlength{\rightmargin}{0.08\textwidth}
        \setlength{\topsep}{0pt}
        \setlength{\partopsep}{0pt}
        \setlength{\parsep}{0pt}
        \setlength{\itemsep}{0pt}
    }
    \item\relax
}{%
    \end{list}
    \endgroup
}

\fancypagestyle{plainnowm}{%
    \fancyhf{}
    \fancyfoot[C]{\color{bookblue}\thepage}
    \renewcommand{\headrulewidth}{0pt}
    \renewcommand{\footrulewidth}{0pt}
}

\newcommand{\booknotetitle}[2][plain]{%
    \clearpage
    \phantomsection
    \addcontentsline{toc}{chapter}{#2}
    \markboth{#2}{}
    \thispagestyle{#1}
    \vspace*{1.2cm}
    \begin{center}
        {\sffamily\heiti\bfseries\fontsize{26}{34}\selectfont\color{bookblue}#2\par}
        \vspace{0.8em}
        {\color{bookteal}\rule{4.5cm}{0.8pt}\par}
    \end{center}
    \vspace{0.6cm}
}

\renewcommand{\proofname}{证明}
\renewcommand{\qedsymbol}{\color{bookteal}\ensuremath{\blacksquare}}

\title{数学分析习题整理}
\author{Chen}
\date{}

\makeatletter
% ==================== 旧版封面/封底设计备份 ====================
% 2026-08 重设计,旧代码用 \iffalse...\fi 封存,新版见下方。
\iffalse
\renewcommand{\maketitle}{%
    \begin{titlepage}
        \thispagestyle{empty}
        \setcounter{page}{0}
        \noindent\hspace*{-\parindent}\begin{tikzpicture}[remember picture,overlay]
            \fill[bookblue] (current page.north west) rectangle ([yshift=-4.2cm]current page.north east);
            \fill[booklight] ([yshift=-4.2cm]current page.north west) rectangle (current page.south east);
            \draw[bookteal!55,line width=0.7pt]
                ([xshift=2.4cm,yshift=4.3cm]current page.south west)
                .. controls +(2.5cm,1.8cm) and +(-2.8cm,-1.3cm) ..
                ([xshift=-2.4cm,yshift=7.2cm]current page.south east);
            \draw[bookteal!35,line width=0.7pt]
                ([xshift=2.5cm,yshift=6.2cm]current page.south west)
                .. controls +(3.0cm,-1.2cm) and +(-3.0cm,1.6cm) ..
                ([xshift=-2.3cm,yshift=3.8cm]current page.south east);
        \end{tikzpicture}
        \vspace*{5.6cm}
        \noindent\color{bookblue}\rule{\textwidth}{1.4pt}\par
        \vspace{8pt}
        {\centering\sffamily\bfseries\fontsize{32}{42}\selectfont\color{bookblue}\@title\par}
        \vspace{8pt}
        \noindent\color{bookteal}\rule{\textwidth}{0.6pt}\par
        \vspace{2.6cm}
        {\centering\Large\color{bookgray}\@author\par}
        \vfill
        {\centering\color{bookteal}\large
            $\displaystyle \lim_{n\to\infty} a_n$\qquad
            $\displaystyle \int_a^b f(x)\,\mathrm{d}x$\qquad
            $\displaystyle \nabla f(x,y)$\par
        }
        \vspace*{1.5cm}
    \end{titlepage}
}

\newcommand{\makebackcover}{%
    \begingroup
    \let\cleardoublepage\clearpage
    \begin{titlepage}
        \thispagestyle{empty}
        \noindent\hspace*{-\parindent}\begin{tikzpicture}[remember picture,overlay]
            \fill[bookblue] (current page.north west) rectangle (current page.south east);
            \fill[booklight,rounded corners=7pt]
                ([xshift=2.2cm,yshift=-3.0cm]current page.north west)
                rectangle
                ([xshift=-2.2cm,yshift=3.0cm]current page.south east);
            \fill[bookteal!80]
                (current page.north west) --
                ([xshift=5.3cm]current page.north west) --
                ([xshift=2.1cm,yshift=-4.4cm]current page.north west) --
                ([yshift=-6.2cm]current page.north west) -- cycle;
            \fill[bookteal!55]
                (current page.south east) --
                ([xshift=-5.8cm]current page.south east) --
                ([xshift=-2.4cm,yshift=4.2cm]current page.south east) --
                ([yshift=6.0cm]current page.south east) -- cycle;
            \foreach \radius in {0.75,1.15,1.55,1.95} {
                \draw[bookteal!55,line width=0.55pt]
                    ([xshift=-2.9cm,yshift=-4.4cm]current page.north east)
                    circle[radius=\radius cm];
            }
            \foreach \radius in {0.65,1.05,1.45} {
                \draw[booklight!65,line width=0.55pt]
                    ([xshift=2.4cm,yshift=3.6cm]current page.south west)
                    circle[radius=\radius cm];
            }
            \fill[bookteal!16,rounded corners=3pt]
                ([xshift=3.5cm,yshift=-7.1cm]current page.north west)
                rectangle
                ([xshift=-5.0cm,yshift=7.2cm]current page.south east);
            \fill[bookteal!30,rounded corners=3pt]
                ([xshift=4.4cm,yshift=-8.0cm]current page.north west)
                rectangle
                ([xshift=-4.1cm,yshift=8.1cm]current page.south east);
            \begin{scope}[shift={([xshift=-0.75cm]current page.center)}]
                \clip[rounded corners=3pt] (-5.25,-4.75) rectangle (5.25,4.75);
                \foreach \v in {-3.6,-3.0,...,3.6} {
                    \draw[bookteal!48,line width=0.45pt,opacity=0.72]
                        plot[domain=-3.6:3.6,samples=45,smooth,variable=\u]
                        ({\u+0.55*\v},{0.28*\u-0.35*\v+0.11*(\u*\u-\v*\v)});
                }
                \foreach \u in {-3.6,-3.0,...,3.6} {
                    \draw[bookblue!42,line width=0.42pt,opacity=0.62]
                        plot[domain=-3.6:3.6,samples=45,smooth,variable=\v]
                        ({\u+0.55*\v},{0.28*\u-0.35*\v+0.11*(\u*\u-\v*\v)});
                }
                \draw[bookteal!70,line width=0.75pt,opacity=0.82]
                    plot[domain=-3.6:3.6,samples=55,smooth,variable=\u]
                    ({\u},{0.28*\u+0.11*\u*\u});
                \draw[bookblue!60,line width=0.7pt,opacity=0.75]
                    plot[domain=-3.6:3.6,samples=55,smooth,variable=\v]
                    ({0.55*\v},{-0.35*\v-0.11*\v*\v});
            \end{scope}
            \draw[bookteal!45,line width=0.65pt]
                ([xshift=3.0cm,yshift=-8.2cm]current page.north west)
                .. controls +(2.2cm,2.0cm) and +(-2.0cm,-2.1cm) ..
                ([xshift=-3.0cm,yshift=9.0cm]current page.south east);
            \draw[bookteal!25,line width=0.55pt]
                ([xshift=3.1cm,yshift=8.0cm]current page.south west)
                .. controls +(2.7cm,-1.7cm) and +(-2.6cm,1.8cm) ..
                ([xshift=-3.1cm,yshift=-8.7cm]current page.north east);
            \node[anchor=west,text=bookteal!72,scale=0.78] at
                ([xshift=2.8cm,yshift=-4.0cm]current page.north west) {%
                $\displaystyle e^{\mathrm{i}\pi}+1=0$};
            \node[anchor=east,text=bookblue!62,scale=0.67] at
                ([xshift=-2.7cm,yshift=4.0cm]current page.south east) {%
                $\displaystyle
                f(x)=\sum_{n=0}^{\infty}\frac{f^{(n)}(a)}{n!}(x-a)^n$};
        \end{tikzpicture}
        \mbox{}
    \end{titlepage}
    \endgroup
}
\fi

% ==================== 新版封面(2026-08 重设计) ====================
% 设计要点:顶部深绿纵向渐变+波浪下缘+两层半透明过渡波(解决深浅绿生硬过渡);
% 深色区螺旋线/同心圆母题;标题+英文副题+作者居中;
% 中部 Fourier 部分和逼近方波曲线;底部浅色波纹(无文字)。
\renewcommand{\maketitle}{%
    \begin{titlepage}
        \thispagestyle{empty}
        \setcounter{page}{0}
        \noindent\hspace*{-\parindent}\begin{tikzpicture}[remember picture,overlay]
            % ---------- 底色 ----------
            \fill[booklight] (current page.north west) rectangle (current page.south east);

            % ---------- 顶部深色区:纵向渐变 + 波浪下缘 ----------
            \shade[top color=bookblue, bottom color=bookteal!55!bookblue]
                (current page.north west) -- (current page.north east) --
                ([yshift=-8.6cm]current page.north east)
                .. controls ([xshift=-5.2cm,yshift=-10.6cm]current page.north east)
                         and ([xshift=6.0cm,yshift=-8.0cm]current page.north west) ..
                ([yshift=-10.0cm]current page.north west) -- cycle;

            % ---------- 过渡层 1:半透明波浪 ----------
            \fill[bookteal!60!bookblue,opacity=0.42]
                ([yshift=-8.6cm]current page.north east)
                .. controls ([xshift=-5.2cm,yshift=-10.6cm]current page.north east)
                         and ([xshift=6.0cm,yshift=-8.0cm]current page.north west) ..
                ([yshift=-10.0cm]current page.north west) --
                ([yshift=-12.6cm]current page.north west)
                .. controls ([xshift=6.5cm,yshift=-11.0cm]current page.north west)
                         and ([xshift=-5.5cm,yshift=-13.2cm]current page.north east) ..
                ([yshift=-11.4cm]current page.north east) -- cycle;

            % ---------- 过渡层 2:更浅的波浪 ----------
            \fill[bookteal!35,opacity=0.35]
                ([yshift=-10.6cm]current page.north east)
                .. controls ([xshift=-6.0cm,yshift=-12.4cm]current page.north east)
                         and ([xshift=5.0cm,yshift=-9.8cm]current page.north west) ..
                ([yshift=-11.8cm]current page.north west) --
                ([yshift=-13.8cm]current page.north west)
                .. controls ([xshift=5.5cm,yshift=-12.6cm]current page.north west)
                         and ([xshift=-6.5cm,yshift=-14.6cm]current page.north east) ..
                ([yshift=-12.8cm]current page.north east) -- cycle;

            % ---------- 深色区装饰:右上螺旋线 ----------
            \begin{scope}[shift={([xshift=-3.4cm,yshift=-4.3cm]current page.north east)}]
                \draw[booklight!45!bookteal,opacity=0.55,line width=0.6pt]
                    plot[domain=0:1440,samples=220,smooth,variable=\t]
                    ({0.0028*\t*cos(\t)},{0.0028*\t*sin(\t)});
            \end{scope}

            % ---------- 深色区装饰:左下同心圆 ----------
            \foreach \r in {0.7,1.25,1.8,2.35,2.9}{
                \draw[booklight!35!bookteal,opacity=0.5,line width=0.5pt]
                    ([xshift=0.4cm,yshift=-7.9cm]current page.north west) circle[radius=\r cm];
            }

            % ---------- 主标题 ----------
            \node[anchor=center,text=booklight] at
                ([yshift=-4.5cm]current page.north)
                {\sffamily\bfseries\fontsize{40}{46}\selectfont \@title};

            % ---------- 分隔短线 ----------
            \draw[booklight!65,line width=1.2pt]
                ([xshift=-2.6cm,yshift=-5.75cm]current page.north) --
                ([xshift=2.6cm,yshift=-5.75cm]current page.north);

            % ---------- 英文副题 ----------
            \node[anchor=center,text=booklight!70!bookteal] at
                ([yshift=-6.55cm]current page.north)
                {\itshape\large Mathematical Analysis \textperiodcentered\
                 A Collection of Exercises};

            % ---------- 作者(紧随标题下方) ----------
            \node[anchor=center,text=booklight!88] at
                ([yshift=-8.2cm]current page.north)
                {\sffamily\large \@author\quad 整理};

            % ---------- 中部:Fourier 部分和逼近方波 ----------
            \begin{scope}[shift={([xshift=-5.2cm,yshift=-18.1cm]current page.north)}]
                % 坐标轴
                \draw[bookgray!55,line width=0.5pt,->] (-0.5,0) -- (10.9,0);
                \draw[bookgray!55,line width=0.5pt,->] (0,-1.75) -- (0,1.95);
                % 目标方波(虚线)
                \draw[bookgray!75,line width=0.6pt,dashed]
                    (0,0.92) -- (5,0.92) (5,-0.92) -- (10,-0.92);
                % N=1
                \draw[bookteal!45,line width=0.7pt]
                    plot[domain=0:360,samples=90,smooth,variable=\x]
                    ({\x/36},{1.18*sin(\x)});
                % N=3
                \draw[bookteal!80,line width=0.8pt]
                    plot[domain=0:360,samples=120,smooth,variable=\x]
                    ({\x/36},{1.18*(sin(\x)+sin(3*\x)/3)});
                % N=7
                \draw[bookblue,line width=0.95pt]
                    plot[domain=0:360,samples=150,smooth,variable=\x]
                    ({\x/36},{1.18*(sin(\x)+sin(3*\x)/3+sin(5*\x)/5+sin(7*\x)/7)});
                % 刻度标注
                \node[anchor=north east,text=bookgray!80,scale=0.62] at (-0.12,0) {$0$};
                \node[anchor=north,text=bookgray!80,scale=0.62] at (5,0) {$\pi$};
                \node[anchor=north,text=bookgray!80,scale=0.62] at (10,0) {$2\pi$};
            \end{scope}
            % 图注
            \node[anchor=center,text=bookgray,scale=0.8] at
                ([yshift=-20.5cm]current page.north)
                {$S_N(x)=\displaystyle\sum_{k=1}^{N}\frac{\sin\bigl((2k-1)x\bigr)}{2k-1}$
                 \ 逐步逼近方波};

            % ---------- 底部浅色波纹(无文字,与顶部波浪呼应) ----------
            \fill[bookteal!22,opacity=0.5]
                ([yshift=2.4cm]current page.south west)
                .. controls ([xshift=6.0cm,yshift=3.9cm]current page.south west)
                         and ([xshift=-6.5cm,yshift=1.0cm]current page.south east) ..
                ([yshift=2.9cm]current page.south east) --
                (current page.south east) -- (current page.south west) -- cycle;
            \fill[bookteal!40!bookblue,opacity=0.28]
                ([yshift=1.3cm]current page.south west)
                .. controls ([xshift=5.5cm,yshift=2.7cm]current page.south west)
                         and ([xshift=-5.5cm,yshift=0.3cm]current page.south east) ..
                ([yshift=1.9cm]current page.south east) --
                (current page.south east) -- (current page.south west) -- cycle;
        \end{tikzpicture}
        \mbox{}
    \end{titlepage}
}

% ==================== 新版封底(纯背景,2026-08 重设计) ====================
% 设计要点:对角渐变底+角部半透明叠层;中下部外摆线(玫瑰形)为视觉主体;
% 底部两层半透明白色波浪;无任何文字。
\newcommand{\makebackcover}{%
    \begingroup
    \let\cleardoublepage\clearpage
    \begin{titlepage}
        \thispagestyle{empty}
        \noindent\hspace*{-\parindent}\begin{tikzpicture}[remember picture,overlay]
            % ---------- 底色:对角渐变 ----------
            \shade[top color=bookblue, bottom color=bookteal!48!bookblue,
                   shading angle=-30]
                (current page.north west) rectangle (current page.south east);

            % ---------- 角部半透明叠层 ----------
            \fill[white,opacity=0.05]
                ([xshift=2.0cm,yshift=1.0cm]current page.north east) circle[radius=7.2cm];
            \fill[white,opacity=0.06]
                ([xshift=-0.5cm,yshift=-1.5cm]current page.north east) circle[radius=3.6cm];

            % ---------- 左下同心圆 ----------
            \foreach \r in {1.0,1.75,2.5,3.25,4.0,4.75}{
                \draw[booklight!40,opacity=0.55,line width=0.55pt]
                    ([xshift=-0.6cm,yshift=-0.4cm]current page.south west) circle[radius=\r cm];
            }
            % ---------- 右下螺旋线 ----------
            \begin{scope}[shift={([xshift=-2.9cm,yshift=3.2cm]current page.south east)}]
                \draw[booklight!45!bookteal,opacity=0.5,line width=0.55pt]
                    plot[domain=0:1260,samples=200,smooth,variable=\t]
                    ({0.0036*\t*cos(\t)},{0.0036*\t*sin(\t)});
            \end{scope}
            % ---------- 左上斜线组 ----------
            \foreach \i in {0,1,2}{
                \draw[bookteal!60!booklight,opacity=0.5,line width=0.6pt]
                    ([xshift={1.6cm+\i*0.45cm}]current page.north west) --
                    ([yshift={-1.6cm-\i*0.45cm}]current page.north west);
            }

            % ---------- 中下部视觉主体:外摆线(玫瑰形) ----------
            \begin{scope}[shift={([yshift=-16.6cm]current page.north)}]
                % 外圈
                \draw[booklight!35,opacity=0.6,line width=0.6pt]
                    (0,0) circle[radius=4.9cm];
                \draw[booklight!22,opacity=0.5,line width=0.5pt]
                    (0,0) circle[radius=5.25cm];
                % 主玫瑰线
                \draw[booklight!55!bookteal,opacity=0.8,line width=0.85pt]
                    plot[domain=0:1080,samples=400,smooth,variable=\t]
                    ({0.335*(8*cos(\t)-5*cos(2.6667*\t))},
                     {0.335*(8*sin(\t)-5*sin(2.6667*\t))});
                % 内层玫瑰线
                \draw[booklight!38!bookteal,opacity=0.6,line width=0.6pt]
                    plot[domain=0:1080,samples=400,smooth,variable=\t]
                    ({0.21*(8*cos(\t)-5*cos(2.6667*\t))},
                     {0.21*(8*sin(\t)-5*sin(2.6667*\t))});
            \end{scope}

            % ---------- 底部半透明波浪 ----------
            \fill[white,opacity=0.06]
                ([yshift=4.6cm]current page.south west)
                .. controls ([xshift=6.5cm,yshift=6.4cm]current page.south west)
                         and ([xshift=-6.0cm,yshift=2.6cm]current page.south east) ..
                ([yshift=5.2cm]current page.south east) --
                (current page.south east) -- (current page.south west) -- cycle;
            \fill[booklight,opacity=0.08]
                ([yshift=2.6cm]current page.south west)
                .. controls ([xshift=5.5cm,yshift=4.3cm]current page.south west)
                         and ([xshift=-6.5cm,yshift=1.0cm]current page.south east) ..
                ([yshift=3.3cm]current page.south east) --
                (current page.south east) -- (current page.south west) -- cycle;
        \end{tikzpicture}
        \mbox{}
    \end{titlepage}
    \endgroup
}
\makeatother

% 页脚右下方常态化水印(署名 + 仓库地址,加粗黑色),已在 main/plain 分页样式中定义,所有页面均生效.

\begin{document}

\maketitle
\frontmatter

\booknotetitle{前言}
\begin{booknotebody}
本书是笔者在南京大学一年级数学分析学习过程中积攒形成的习题整理,
题目主要来源于谢惠民等老师编著的数学分析习题课讲义以及笔者在大一期间遇到的数学分析课堂例题和南京大学数学分析期中与期末考试卷. \\

本书按照数列,实数系基本定理,函数,一元微分学与一元积分学以及多元微分学与多元积分学的顺序编排. \\

选题侧重较有难度,方法典型或思路新颖的题目,并保留推导过程与解题思路,便于日后复习,补充和查阅. \\

本书仅作为开源共享的笔记式学习资料,
由于个人水平所限,书中难免存在疏漏,后续将结合复习体会和读者反馈持续修订,同时补充更多有价值有思想的数学分析题目.


\end{booknotebody}
\clearpage

\booknotetitle{本书简介}
\begin{booknotebody}
本书依照数学分析课程的学习顺序编排,各章选取具有代表性的习题和证明,并不以完整覆盖全部知识点为目标.

\textbf{第一章\quad 数列}.本章讨论数列收敛,单调性与有界性,递推数列,嵌套根式,组合平均,上下极限及相关极限方法,并包含自然对数底无理性的证明.

\textbf{第二章\quad 实数系基本定理}.本章讨论上确界,单调子列,Bolzano定理,Cauchy收敛准则,闭区间套,有限覆盖以及次可乘数列,着重展示实数完备性的不同表述与应用.

\textbf{第三章\quad 函数}.本章讨论连续性,一致连续性,单侧极限,单调函数,间断点,周期函数及最小正周期,并呈现介值性,连通性与连续性之间的联系.

\textbf{第四章\quad 一元微分学与一元积分学}.本章讨论Rolle定理和Lagrange中值定理的拓展证明,导数与高阶导数估计,积分零点问题以及若干中值定理构造.

\textbf{第五章\quad 多元微分学与多元积分学}.本章讨论齐次函数,隐式零点函数,凹性,二元Riemann可积性,Hölder曲线,调和函数,Gauss公式,曲线积分和球面平均公式.

\end{booknotebody}
\clearpage

\booknotetitle{使用说明}
\begin{booknotebody}
\begin{itemize}[leftmargin=2\ccwd,itemsep=0.55em,topsep=0pt]
    \item 各章题号独立编号,新章从题1开始.
    \item 【题x】表示普通题目,【拓展1】和【拓展2】表示对相关定理或证明方法的扩展.
    \item \textbf{证明}表示主要解答;\textbf{另解}表示可以独立完成论证的另一种方法.若另解只替换原证明中的局部步骤,则由上下文说明.
    \item \textbf{注}和\textbf{注记}用于就近补充条件,记号或结论,有时也用于解释思考动机.
    \item \textbf{一点思考}用于记录思考动机,试探过程或非正式想法,不替代严格证明.
    \item 题目文字未作改动,部分证明和讲解依据课堂笔记转载整理.
    \item 本书目前不提供逐题出处,统一来源见前言及书末的\textbf{参考资料与来源说明}.
\end{itemize}
\end{booknotebody}
\clearpage

\booknotetitle{符号与约定}
\begin{booknotelongbody}
\textbf{数集与数列}

\begin{itemize}[leftmargin=2\ccwd,itemsep=0.5em,topsep=0.4em]
    \item 约定$\displaystyle\mathbb{N}=\{0,1,2,\ldots\}$,$\displaystyle\mathbb{N}^+=\{1,2,3,\ldots\}$, \\
    $\displaystyle\mathbb Z$,$\displaystyle\mathbb Q$,$\displaystyle\mathbb R$分别表示整数集,有理数集和实数集.
    \item 数列写作$\displaystyle\{a_n\}$,$\;\displaystyle\varlimsup a_n$和$\displaystyle\varliminf a_n$分别表示数列的上极限和下极限.
    \item 无穷大量指$\displaystyle|a_n|\to+\infty$;特指趋于正无穷时写作$\displaystyle a_n\to+\infty$.无界数列不一定是无穷大量.
\end{itemize}

\textbf{区间,邻域与拓扑}

\begin{itemize}[leftmargin=2\ccwd,itemsep=0.5em,topsep=0.4em]
    \item 区间均指$\displaystyle\mathbb R$中的区间.端点属于定义域时,连续性和极限按定义域的相对单侧意义理解.
    \item $\displaystyle\overline\Omega$,$\displaystyle\partial\Omega$和$\displaystyle\operatorname{int}\Omega$分别表示$\displaystyle\Omega$的闭包,边界和内部.
    \item $\displaystyle B_r(a)$和$\displaystyle\overline{B_r(a)}$分别表示以$\displaystyle a$为中心,半径为$\displaystyle r$的开球和闭球.
    \item 区域指非空连通开集;涉及闭区域时明确写其闭包.
\end{itemize}

\textbf{凸集,凸函数与凹函数}

\begin{itemize}[leftmargin=2\ccwd,itemsep=0.5em,topsep=0.4em]
    \item 集合$\displaystyle D$称为凸集,是指对任意$\displaystyle x,y\in D$和$\displaystyle t\in[0,1]$,均有$\displaystyle(1-t)x+ty\in D$.
    \item 函数$\displaystyle f$称为凸函数,是指其定义域为凸集,且$\displaystyle f((1-t)x+ty)\leq(1-t)f(x)+tf(y)$;凹函数采用反向不等式.
    \item 严格凸函数和严格凹函数要求当$\displaystyle x\neq y$且$\displaystyle t\in(0,1)$时,相应不等式严格成立.
    \item 凸集$\displaystyle D$称为严格凸集,是指对任意不同的$\displaystyle x,y\in D$和$\displaystyle t\in(0,1)$,均有$\displaystyle(1-t)x+ty\in\operatorname{int}D$.
\end{itemize}

\textbf{范数,微分与函数空间}

\begin{itemize}[leftmargin=2\ccwd,itemsep=0.5em,topsep=0.4em]
    \item $\displaystyle\lVert x\rVert$默认表示欧氏范数,$\displaystyle x\cdot y$表示欧氏内积,$\displaystyle|x|$表示实数的绝对值;在向量语境中,$\displaystyle|x|$亦可表示欧氏范数.
    \item $\displaystyle\nabla u$,$\displaystyle\Delta u$和$\displaystyle\operatorname{div}X$分别表示梯度,$Laplace$算子和散度.
    \item $\displaystyle C(I)$与$\displaystyle C^0(I)$均表示$\displaystyle I$上的连续函数类;$\displaystyle C^k(I)$表示具有连续到$\displaystyle k$阶导数的函数类.
    \item $\displaystyle C^k(\overline\Omega)$表示函数及其不超过$\displaystyle k$阶的偏导数均可连续延拓至$\displaystyle\overline\Omega$.
    \item $\displaystyle C^\alpha[a,b]$表示Hölder指数为$\displaystyle\alpha$的函数类,即存在$\displaystyle M>0$,使任意$\displaystyle s,t\in[a,b]$均满足$\displaystyle|f(t)-f(s)|\leq M|t-s|^\alpha$.
\end{itemize}

\textbf{曲线,曲面与方向}

\begin{itemize}[leftmargin=2\ccwd,itemsep=0.5em,topsep=0.4em]
    \item 平面简单闭曲线未另作说明时,正向指逆时针方向.
    \item $\displaystyle\partial\Omega$上的单位法向量未另作说明时取外法向量.对于挖孔区域,内边界的法向量仍按该区域的外法向量取向.
    \item 曲线积分和曲面积分按题设指定的方向或定向理解;若结果依赖方向,则题设中的方向条件不可省略.
    \item $\displaystyle\operatorname{sgn}(a)$表示符号函数,即当$\displaystyle a>0$,$\displaystyle a=0$和$\displaystyle a<0$时,其值分别为$\displaystyle1$,$\displaystyle0$和$\displaystyle-1$.
\end{itemize}


\end{booknotelongbody}
\clearpage

\tableofcontents

\clearpage
\mainmatter
\pagestyle{main}

\def\BookMaster{}
\ifdefined\BookMaster
\def\BookEnd{}
\else
\documentclass[a4paper,12pt,fleqn,openany,twoside,fontset=windows]{ctexbook}

\usepackage[fleqn]{amsmath}
\usepackage{amssymb,amsthm,mathrsfs}
\usepackage{geometry}
\usepackage[dvipsnames]{xcolor}
\usepackage{graphicx}
\usepackage{tikz}
\usetikzlibrary{calc}
\usepackage[most]{tcolorbox}
\usepackage{fancyhdr}
\usepackage{titlesec}
\usepackage{tocloft}
\usepackage{enumitem}
\usepackage{microtype}
\usepackage{pifont}
\usepackage{hyperref}
\usepackage{romannum}
\usepackage{mathtools}

\definecolor{bookblue}{HTML}{264653}
\definecolor{bookteal}{HTML}{4F7C82}
\definecolor{booklight}{HTML}{EDF4F3}
\definecolor{bookgray}{HTML}{5E686B}

\geometry{
    inner=2.7cm,
    outer=2.5cm,
    top=2.7cm,
    bottom=2.7cm,
    headheight=18pt,
    headsep=18pt,
    footskip=25pt
}
\setlength{\mathindent}{0pt}
\setlength{\parindent}{5pt}
\setlength{\parskip}{3pt plus 1pt minus 1pt}
\linespread{1.25}
\allowdisplaybreaks[3]
\AtBeginDocument{%
    \setlength{\parindent}{5pt}
    \setlength{\abovedisplayskip}{8pt plus 2pt minus 2pt}
    \setlength{\belowdisplayskip}{8pt plus 2pt minus 2pt}
    \setlength{\abovedisplayshortskip}{5pt plus 1pt minus 1pt}
    \setlength{\belowdisplayshortskip}{5pt plus 1pt minus 1pt}
}

\hypersetup{
    unicode=true,
    colorlinks=true,
    linkcolor=bookblue,
    citecolor=bookblue,
    urlcolor=bookteal,
    bookmarksopen=true,
    pdfauthor={Chen},
    pdftitle={数学分析习题整理}
}

\renewcommand{\chaptermark}[1]{%
    \edef\@tmp{第\thechapter 章\quad #1}%
    \expandafter\markboth\expandafter{\@tmp}{}%
}
\fancypagestyle{main}{%
    \fancyhf{}
    \fancyhead[LE]{\small\color{bookgray}\nouppercase{\leftmark}}
    \fancyhead[RO]{\small\color{bookgray}数学分析习题整理}
    \fancyfoot[C]{\color{bookblue}\thepage}
    \fancyfoot[R]{\small\bfseries Chen\;\; github.com/mathlover-1/math-analysis}
    \renewcommand{\headrulewidth}{0.4pt}
    \renewcommand{\footrulewidth}{0pt}
}
\fancypagestyle{plain}{%
    \fancyhf{}
    \fancyfoot[C]{\color{bookblue}\thepage}
    \fancyfoot[R]{\small\bfseries Chen\;\; github.com/mathlover-1/math-analysis}
    \renewcommand{\headrulewidth}{0pt}
    \renewcommand{\footrulewidth}{0pt}
}
\pagestyle{main}

\titleformat{\chapter}[display]
  {\normalfont\sffamily\bfseries\color{bookblue}}
  {\raggedleft\fontsize{46}{50}\selectfont\thechapter}
  {-18pt}
  {\titlerule[1.2pt]\vspace{10pt}\filright\Huge}
\titlespacing*{\chapter}{0pt}{-15pt}{36pt}

\renewcommand{\contentsname}{目\quad 录}
\renewcommand{\cfttoctitlefont}{\hfill\sffamily\bfseries\Huge\color{bookblue}}
\renewcommand{\cftaftertoctitle}{\hfill}
\setlength{\cftbeforetoctitleskip}{5pt}
\setlength{\cftaftertoctitleskip}{28pt}
\setlength{\cftbeforechapskip}{12pt}
\renewcommand{\cftchapfont}{\sffamily\bfseries\large\color{bookblue}}
\renewcommand{\cftchappagefont}{\sffamily\bfseries\large\color{bookblue}}
\renewcommand{\cftchapleader}{\color{bookteal}\cftdotfill{2.5}}

\newtcolorbox{ti}{
    enhanced,
    breakable,
    colback=booklight!55!white,
    colframe=bookteal!35!white,
    boxrule=0.45pt,
    boxsep=0pt,
    left=10pt,
    right=10pt,
    top=9pt,
    bottom=9pt,
    before skip=14pt,
    after skip=14pt,
    arc=2pt,
    outer arc=2pt,
    before upper={\parindent=0pt}
}

\newenvironment{booknotebody}{%
    \begin{center}
    \begin{minipage}{0.84\textwidth}
    \songti\zihao{-4}\linespread{1.45}\selectfont
    \setlength{\parindent}{2\ccwd}
    \setlength{\parskip}{0.75em}
}{%
    \end{minipage}
    \end{center}
}

\newenvironment{booknotelongbody}{%
    \begingroup
    \songti\zihao{-4}\linespread{1.45}\selectfont
    \setlength{\parindent}{2\ccwd}
    \setlength{\parskip}{0.75em}
    \begin{list}{}{%
        \setlength{\leftmargin}{0.08\textwidth}
        \setlength{\rightmargin}{0.08\textwidth}
        \setlength{\topsep}{0pt}
        \setlength{\partopsep}{0pt}
        \setlength{\parsep}{0pt}
        \setlength{\itemsep}{0pt}
    }
    \item\relax
}{%
    \end{list}
    \endgroup
}

\fancypagestyle{plainnowm}{%
    \fancyhf{}
    \fancyfoot[C]{\color{bookblue}\thepage}
    \renewcommand{\headrulewidth}{0pt}
    \renewcommand{\footrulewidth}{0pt}
}

\newcommand{\booknotetitle}[2][plain]{%
    \clearpage
    \phantomsection
    \addcontentsline{toc}{chapter}{#2}
    \markboth{#2}{}
    \thispagestyle{#1}
    \vspace*{1.2cm}
    \begin{center}
        {\sffamily\heiti\bfseries\fontsize{26}{34}\selectfont\color{bookblue}#2\par}
        \vspace{0.8em}
        {\color{bookteal}\rule{4.5cm}{0.8pt}\par}
    \end{center}
    \vspace{0.6cm}
}

\renewcommand{\proofname}{证明}
\renewcommand{\qedsymbol}{\color{bookteal}\ensuremath{\blacksquare}}

\title{数学分析习题整理}
\author{Chen}
\date{}

\makeatletter
% ==================== 旧版封面/封底设计备份 ====================
% 2026-08 重设计,旧代码用 \iffalse...\fi 封存,新版见下方。
\iffalse
\renewcommand{\maketitle}{%
    \begin{titlepage}
        \thispagestyle{empty}
        \setcounter{page}{0}
        \noindent\hspace*{-\parindent}\begin{tikzpicture}[remember picture,overlay]
            \fill[bookblue] (current page.north west) rectangle ([yshift=-4.2cm]current page.north east);
            \fill[booklight] ([yshift=-4.2cm]current page.north west) rectangle (current page.south east);
            \draw[bookteal!55,line width=0.7pt]
                ([xshift=2.4cm,yshift=4.3cm]current page.south west)
                .. controls +(2.5cm,1.8cm) and +(-2.8cm,-1.3cm) ..
                ([xshift=-2.4cm,yshift=7.2cm]current page.south east);
            \draw[bookteal!35,line width=0.7pt]
                ([xshift=2.5cm,yshift=6.2cm]current page.south west)
                .. controls +(3.0cm,-1.2cm) and +(-3.0cm,1.6cm) ..
                ([xshift=-2.3cm,yshift=3.8cm]current page.south east);
        \end{tikzpicture}
        \vspace*{5.6cm}
        \noindent\color{bookblue}\rule{\textwidth}{1.4pt}\par
        \vspace{8pt}
        {\centering\sffamily\bfseries\fontsize{32}{42}\selectfont\color{bookblue}\@title\par}
        \vspace{8pt}
        \noindent\color{bookteal}\rule{\textwidth}{0.6pt}\par
        \vspace{2.6cm}
        {\centering\Large\color{bookgray}\@author\par}
        \vfill
        {\centering\color{bookteal}\large
            $\displaystyle \lim_{n\to\infty} a_n$\qquad
            $\displaystyle \int_a^b f(x)\,\mathrm{d}x$\qquad
            $\displaystyle \nabla f(x,y)$\par
        }
        \vspace*{1.5cm}
    \end{titlepage}
}

\newcommand{\makebackcover}{%
    \begingroup
    \let\cleardoublepage\clearpage
    \begin{titlepage}
        \thispagestyle{empty}
        \noindent\hspace*{-\parindent}\begin{tikzpicture}[remember picture,overlay]
            \fill[bookblue] (current page.north west) rectangle (current page.south east);
            \fill[booklight,rounded corners=7pt]
                ([xshift=2.2cm,yshift=-3.0cm]current page.north west)
                rectangle
                ([xshift=-2.2cm,yshift=3.0cm]current page.south east);
            \fill[bookteal!80]
                (current page.north west) --
                ([xshift=5.3cm]current page.north west) --
                ([xshift=2.1cm,yshift=-4.4cm]current page.north west) --
                ([yshift=-6.2cm]current page.north west) -- cycle;
            \fill[bookteal!55]
                (current page.south east) --
                ([xshift=-5.8cm]current page.south east) --
                ([xshift=-2.4cm,yshift=4.2cm]current page.south east) --
                ([yshift=6.0cm]current page.south east) -- cycle;
            \foreach \radius in {0.75,1.15,1.55,1.95} {
                \draw[bookteal!55,line width=0.55pt]
                    ([xshift=-2.9cm,yshift=-4.4cm]current page.north east)
                    circle[radius=\radius cm];
            }
            \foreach \radius in {0.65,1.05,1.45} {
                \draw[booklight!65,line width=0.55pt]
                    ([xshift=2.4cm,yshift=3.6cm]current page.south west)
                    circle[radius=\radius cm];
            }
            \fill[bookteal!16,rounded corners=3pt]
                ([xshift=3.5cm,yshift=-7.1cm]current page.north west)
                rectangle
                ([xshift=-5.0cm,yshift=7.2cm]current page.south east);
            \fill[bookteal!30,rounded corners=3pt]
                ([xshift=4.4cm,yshift=-8.0cm]current page.north west)
                rectangle
                ([xshift=-4.1cm,yshift=8.1cm]current page.south east);
            \begin{scope}[shift={([xshift=-0.75cm]current page.center)}]
                \clip[rounded corners=3pt] (-5.25,-4.75) rectangle (5.25,4.75);
                \foreach \v in {-3.6,-3.0,...,3.6} {
                    \draw[bookteal!48,line width=0.45pt,opacity=0.72]
                        plot[domain=-3.6:3.6,samples=45,smooth,variable=\u]
                        ({\u+0.55*\v},{0.28*\u-0.35*\v+0.11*(\u*\u-\v*\v)});
                }
                \foreach \u in {-3.6,-3.0,...,3.6} {
                    \draw[bookblue!42,line width=0.42pt,opacity=0.62]
                        plot[domain=-3.6:3.6,samples=45,smooth,variable=\v]
                        ({\u+0.55*\v},{0.28*\u-0.35*\v+0.11*(\u*\u-\v*\v)});
                }
                \draw[bookteal!70,line width=0.75pt,opacity=0.82]
                    plot[domain=-3.6:3.6,samples=55,smooth,variable=\u]
                    ({\u},{0.28*\u+0.11*\u*\u});
                \draw[bookblue!60,line width=0.7pt,opacity=0.75]
                    plot[domain=-3.6:3.6,samples=55,smooth,variable=\v]
                    ({0.55*\v},{-0.35*\v-0.11*\v*\v});
            \end{scope}
            \draw[bookteal!45,line width=0.65pt]
                ([xshift=3.0cm,yshift=-8.2cm]current page.north west)
                .. controls +(2.2cm,2.0cm) and +(-2.0cm,-2.1cm) ..
                ([xshift=-3.0cm,yshift=9.0cm]current page.south east);
            \draw[bookteal!25,line width=0.55pt]
                ([xshift=3.1cm,yshift=8.0cm]current page.south west)
                .. controls +(2.7cm,-1.7cm) and +(-2.6cm,1.8cm) ..
                ([xshift=-3.1cm,yshift=-8.7cm]current page.north east);
            \node[anchor=west,text=bookteal!72,scale=0.78] at
                ([xshift=2.8cm,yshift=-4.0cm]current page.north west) {%
                $\displaystyle e^{\mathrm{i}\pi}+1=0$};
            \node[anchor=east,text=bookblue!62,scale=0.67] at
                ([xshift=-2.7cm,yshift=4.0cm]current page.south east) {%
                $\displaystyle
                f(x)=\sum_{n=0}^{\infty}\frac{f^{(n)}(a)}{n!}(x-a)^n$};
        \end{tikzpicture}
        \mbox{}
    \end{titlepage}
    \endgroup
}
\fi

% ==================== 新版封面(2026-08 重设计) ====================
% 设计要点:顶部深绿纵向渐变+波浪下缘+两层半透明过渡波(解决深浅绿生硬过渡);
% 深色区螺旋线/同心圆母题;标题+英文副题+作者居中;
% 中部 Fourier 部分和逼近方波曲线;底部浅色波纹(无文字)。
\renewcommand{\maketitle}{%
    \begin{titlepage}
        \thispagestyle{empty}
        \setcounter{page}{0}
        \noindent\hspace*{-\parindent}\begin{tikzpicture}[remember picture,overlay]
            % ---------- 底色 ----------
            \fill[booklight] (current page.north west) rectangle (current page.south east);

            % ---------- 顶部深色区:纵向渐变 + 波浪下缘 ----------
            \shade[top color=bookblue, bottom color=bookteal!55!bookblue]
                (current page.north west) -- (current page.north east) --
                ([yshift=-8.6cm]current page.north east)
                .. controls ([xshift=-5.2cm,yshift=-10.6cm]current page.north east)
                         and ([xshift=6.0cm,yshift=-8.0cm]current page.north west) ..
                ([yshift=-10.0cm]current page.north west) -- cycle;

            % ---------- 过渡层 1:半透明波浪 ----------
            \fill[bookteal!60!bookblue,opacity=0.42]
                ([yshift=-8.6cm]current page.north east)
                .. controls ([xshift=-5.2cm,yshift=-10.6cm]current page.north east)
                         and ([xshift=6.0cm,yshift=-8.0cm]current page.north west) ..
                ([yshift=-10.0cm]current page.north west) --
                ([yshift=-12.6cm]current page.north west)
                .. controls ([xshift=6.5cm,yshift=-11.0cm]current page.north west)
                         and ([xshift=-5.5cm,yshift=-13.2cm]current page.north east) ..
                ([yshift=-11.4cm]current page.north east) -- cycle;

            % ---------- 过渡层 2:更浅的波浪 ----------
            \fill[bookteal!35,opacity=0.35]
                ([yshift=-10.6cm]current page.north east)
                .. controls ([xshift=-6.0cm,yshift=-12.4cm]current page.north east)
                         and ([xshift=5.0cm,yshift=-9.8cm]current page.north west) ..
                ([yshift=-11.8cm]current page.north west) --
                ([yshift=-13.8cm]current page.north west)
                .. controls ([xshift=5.5cm,yshift=-12.6cm]current page.north west)
                         and ([xshift=-6.5cm,yshift=-14.6cm]current page.north east) ..
                ([yshift=-12.8cm]current page.north east) -- cycle;

            % ---------- 深色区装饰:右上螺旋线 ----------
            \begin{scope}[shift={([xshift=-3.4cm,yshift=-4.3cm]current page.north east)}]
                \draw[booklight!45!bookteal,opacity=0.55,line width=0.6pt]
                    plot[domain=0:1440,samples=220,smooth,variable=\t]
                    ({0.0028*\t*cos(\t)},{0.0028*\t*sin(\t)});
            \end{scope}

            % ---------- 深色区装饰:左下同心圆 ----------
            \foreach \r in {0.7,1.25,1.8,2.35,2.9}{
                \draw[booklight!35!bookteal,opacity=0.5,line width=0.5pt]
                    ([xshift=0.4cm,yshift=-7.9cm]current page.north west) circle[radius=\r cm];
            }

            % ---------- 主标题 ----------
            \node[anchor=center,text=booklight] at
                ([yshift=-4.5cm]current page.north)
                {\sffamily\bfseries\fontsize{40}{46}\selectfont \@title};

            % ---------- 分隔短线 ----------
            \draw[booklight!65,line width=1.2pt]
                ([xshift=-2.6cm,yshift=-5.75cm]current page.north) --
                ([xshift=2.6cm,yshift=-5.75cm]current page.north);

            % ---------- 英文副题 ----------
            \node[anchor=center,text=booklight!70!bookteal] at
                ([yshift=-6.55cm]current page.north)
                {\itshape\large Mathematical Analysis \textperiodcentered\
                 A Collection of Exercises};

            % ---------- 作者(紧随标题下方) ----------
            \node[anchor=center,text=booklight!88] at
                ([yshift=-8.2cm]current page.north)
                {\sffamily\large \@author\quad 整理};

            % ---------- 中部:Fourier 部分和逼近方波 ----------
            \begin{scope}[shift={([xshift=-5.2cm,yshift=-18.1cm]current page.north)}]
                % 坐标轴
                \draw[bookgray!55,line width=0.5pt,->] (-0.5,0) -- (10.9,0);
                \draw[bookgray!55,line width=0.5pt,->] (0,-1.75) -- (0,1.95);
                % 目标方波(虚线)
                \draw[bookgray!75,line width=0.6pt,dashed]
                    (0,0.92) -- (5,0.92) (5,-0.92) -- (10,-0.92);
                % N=1
                \draw[bookteal!45,line width=0.7pt]
                    plot[domain=0:360,samples=90,smooth,variable=\x]
                    ({\x/36},{1.18*sin(\x)});
                % N=3
                \draw[bookteal!80,line width=0.8pt]
                    plot[domain=0:360,samples=120,smooth,variable=\x]
                    ({\x/36},{1.18*(sin(\x)+sin(3*\x)/3)});
                % N=7
                \draw[bookblue,line width=0.95pt]
                    plot[domain=0:360,samples=150,smooth,variable=\x]
                    ({\x/36},{1.18*(sin(\x)+sin(3*\x)/3+sin(5*\x)/5+sin(7*\x)/7)});
                % 刻度标注
                \node[anchor=north east,text=bookgray!80,scale=0.62] at (-0.12,0) {$0$};
                \node[anchor=north,text=bookgray!80,scale=0.62] at (5,0) {$\pi$};
                \node[anchor=north,text=bookgray!80,scale=0.62] at (10,0) {$2\pi$};
            \end{scope}
            % 图注
            \node[anchor=center,text=bookgray,scale=0.8] at
                ([yshift=-20.5cm]current page.north)
                {$S_N(x)=\displaystyle\sum_{k=1}^{N}\frac{\sin\bigl((2k-1)x\bigr)}{2k-1}$
                 \ 逐步逼近方波};

            % ---------- 底部浅色波纹(无文字,与顶部波浪呼应) ----------
            \fill[bookteal!22,opacity=0.5]
                ([yshift=2.4cm]current page.south west)
                .. controls ([xshift=6.0cm,yshift=3.9cm]current page.south west)
                         and ([xshift=-6.5cm,yshift=1.0cm]current page.south east) ..
                ([yshift=2.9cm]current page.south east) --
                (current page.south east) -- (current page.south west) -- cycle;
            \fill[bookteal!40!bookblue,opacity=0.28]
                ([yshift=1.3cm]current page.south west)
                .. controls ([xshift=5.5cm,yshift=2.7cm]current page.south west)
                         and ([xshift=-5.5cm,yshift=0.3cm]current page.south east) ..
                ([yshift=1.9cm]current page.south east) --
                (current page.south east) -- (current page.south west) -- cycle;
        \end{tikzpicture}
        \mbox{}
    \end{titlepage}
}

% ==================== 新版封底(纯背景,2026-08 重设计) ====================
% 设计要点:对角渐变底+角部半透明叠层;中下部外摆线(玫瑰形)为视觉主体;
% 底部两层半透明白色波浪;无任何文字。
\newcommand{\makebackcover}{%
    \begingroup
    \let\cleardoublepage\clearpage
    \begin{titlepage}
        \thispagestyle{empty}
        \noindent\hspace*{-\parindent}\begin{tikzpicture}[remember picture,overlay]
            % ---------- 底色:对角渐变 ----------
            \shade[top color=bookblue, bottom color=bookteal!48!bookblue,
                   shading angle=-30]
                (current page.north west) rectangle (current page.south east);

            % ---------- 角部半透明叠层 ----------
            \fill[white,opacity=0.05]
                ([xshift=2.0cm,yshift=1.0cm]current page.north east) circle[radius=7.2cm];
            \fill[white,opacity=0.06]
                ([xshift=-0.5cm,yshift=-1.5cm]current page.north east) circle[radius=3.6cm];

            % ---------- 左下同心圆 ----------
            \foreach \r in {1.0,1.75,2.5,3.25,4.0,4.75}{
                \draw[booklight!40,opacity=0.55,line width=0.55pt]
                    ([xshift=-0.6cm,yshift=-0.4cm]current page.south west) circle[radius=\r cm];
            }
            % ---------- 右下螺旋线 ----------
            \begin{scope}[shift={([xshift=-2.9cm,yshift=3.2cm]current page.south east)}]
                \draw[booklight!45!bookteal,opacity=0.5,line width=0.55pt]
                    plot[domain=0:1260,samples=200,smooth,variable=\t]
                    ({0.0036*\t*cos(\t)},{0.0036*\t*sin(\t)});
            \end{scope}
            % ---------- 左上斜线组 ----------
            \foreach \i in {0,1,2}{
                \draw[bookteal!60!booklight,opacity=0.5,line width=0.6pt]
                    ([xshift={1.6cm+\i*0.45cm}]current page.north west) --
                    ([yshift={-1.6cm-\i*0.45cm}]current page.north west);
            }

            % ---------- 中下部视觉主体:外摆线(玫瑰形) ----------
            \begin{scope}[shift={([yshift=-16.6cm]current page.north)}]
                % 外圈
                \draw[booklight!35,opacity=0.6,line width=0.6pt]
                    (0,0) circle[radius=4.9cm];
                \draw[booklight!22,opacity=0.5,line width=0.5pt]
                    (0,0) circle[radius=5.25cm];
                % 主玫瑰线
                \draw[booklight!55!bookteal,opacity=0.8,line width=0.85pt]
                    plot[domain=0:1080,samples=400,smooth,variable=\t]
                    ({0.335*(8*cos(\t)-5*cos(2.6667*\t))},
                     {0.335*(8*sin(\t)-5*sin(2.6667*\t))});
                % 内层玫瑰线
                \draw[booklight!38!bookteal,opacity=0.6,line width=0.6pt]
                    plot[domain=0:1080,samples=400,smooth,variable=\t]
                    ({0.21*(8*cos(\t)-5*cos(2.6667*\t))},
                     {0.21*(8*sin(\t)-5*sin(2.6667*\t))});
            \end{scope}

            % ---------- 底部半透明波浪 ----------
            \fill[white,opacity=0.06]
                ([yshift=4.6cm]current page.south west)
                .. controls ([xshift=6.5cm,yshift=6.4cm]current page.south west)
                         and ([xshift=-6.0cm,yshift=2.6cm]current page.south east) ..
                ([yshift=5.2cm]current page.south east) --
                (current page.south east) -- (current page.south west) -- cycle;
            \fill[booklight,opacity=0.08]
                ([yshift=2.6cm]current page.south west)
                .. controls ([xshift=5.5cm,yshift=4.3cm]current page.south west)
                         and ([xshift=-6.5cm,yshift=1.0cm]current page.south east) ..
                ([yshift=3.3cm]current page.south east) --
                (current page.south east) -- (current page.south west) -- cycle;
        \end{tikzpicture}
        \mbox{}
    \end{titlepage}
    \endgroup
}
\makeatother

% 页脚右下方常态化水印(署名 + 仓库地址,加粗黑色),已在 main/plain 分页样式中定义,所有页面均生效.

\begin{document}
\mainmatter
\pagestyle{main}
\setcounter{chapter}{0}
\def\BookEnd{\end{document}}
\fi

\chapter{数列}

% 在此处使用 ti 环境添加题目,例如:
% \begin{ti}
% 证明:若数列 $\{a_n\}$ 单调有界,则必收敛.
% \end{ti}
%
% 支持行间公式:
% \begin{ti}
% 求极限
% \[
%   \lim_{n \to \infty} \frac{a_n}{b_n}
% \]
% \end{ti}

\begin{ti}
【题1】 设$a_n=\left[\dfrac{(2n-1)!!}{(2n)!!}\right]^2$,证明:$\{a_n\}$收敛于0.
\end{ti}

\textbf{证明}:在不使用斯特林公式的情况下,对任意$\displaystyle n\geq1$,有 \\

$\displaystyle\prod_{k=1}^{n}\dfrac{(2k-1)(2k+1)}{(2k)^2}=\dfrac{1\cdot3^2\cdot5^2\cdots(2n-1)^2(2n+1)}{2^2\cdot4^2\cdots(2n)^2}=(2n+1)a_n,$ \\

又对每个$\displaystyle k\geq1$,有$\displaystyle 0<\dfrac{(2k-1)(2k+1)}{(2k)^2}=1-\dfrac{1}{4k^2}<1,$ \\

故$\displaystyle 0<(2n+1)a_n<1$,从而$\displaystyle 0<a_n<\dfrac{1}{2n+1}\to0,$ \\

由夹逼准则,$\displaystyle\lim_{n\to\infty}a_n=0.$

\begin{ti}
【题2】 $a_1> 0,\quad a_{n+1}=a_n+\dfrac{1}{a_n}\quad \text{证明}:\displaystyle\lim_{n\to\infty}\dfrac{a_n}{\sqrt{2n}}=1$
\end{ti}

\textbf{证明}: 易得$a_n> 0$,且$a_{n+1}> a_n,$ \\

若$\{a_n\}$收敛,设收敛于$a,\;a\in\mathbb{R},\text{则}a=a+\dfrac{1}{a}$,无解, \\

若$\displaystyle\{a_n\}$有上界,则由单调收敛定理必收敛,与上面的结论矛盾. \\

故$\displaystyle\{a_n\}$无上界;又$\displaystyle\{a_n\}$单调递增,所以$\displaystyle a_n\to+\infty$,则$\displaystyle\lim_{n\to\infty}\dfrac{1}{a_n}=0,$ \\

$\because\displaystyle\lim_{n\to\infty}\dfrac{a_{n+1}^2-a_n^2}{2(n+1)-2n}=\lim_{n\to\infty}\dfrac{\frac{1}{a_n^2}+2}{2}=1,$ \\

且$b_n=2n$为严格单调递增且发散到$+\infty$的数列, \\

由\textit{Stolz}定理,$\quad\displaystyle\lim_{n\to\infty}\dfrac{a_n^2}{2n}=1,$ \\

$\therefore\displaystyle\lim_{n\to\infty}\dfrac{a_n}{\sqrt{2n}}=1.$ \\

\textbf{一点思考}:关于为什么要平方,我的思考是:若有$\displaystyle\lim_{n\to\infty}(b_{n+1}-b_n)=0,$ \\

$\text{则由Stolz公式},\displaystyle\lim_{n\to\infty}\dfrac{b_n}{n}=0$, 而这里$a_{n+1}-a_n\rightarrow 0,$ \\

$\text{但是题目要求}\dfrac{a_n}{\sqrt{2n}}\rightarrow 1$,这说明$a_{n+1}-a_n\rightarrow 0$的条件不够强, \\

应该研究$b_n=a_n^2,\text{研究}b_{n+1}-b_n$.

\begin{ti}
【题3】 设$0< x_i\leq \frac{1}{2}\quad i=1,2,...,n,\text{证明}:  $ \\

$\dfrac{\displaystyle\prod_{i=1}^n x_i}{(\displaystyle\sum_{i=1}^n x_i)^n}\leq
\dfrac{\displaystyle\prod_{i=1}^n (1-x_i)}{[\displaystyle\sum_{i=1}^n (1-x_i)]^n}$
\end{ti}

\textbf{证明}:$\;n=1\text{时,上式显然成立},$ \\

$n=2\text{时,即证}\dfrac{x_1x_2}{(x_1+x_2)^2}\leq\dfrac{(1-x_1)(1-x_2)}{(2-x_1-x_2)^2},$ \\

$\text{即证}\quad x_1x_2\left[4+(x_1+x_2)^2-4(x_1+x_2)\right]\leq(x_1+x_2)^2\left[1+x_1x_2-(x_1+x_2)\right],$ \\

$\text{即证}\quad 4x_1x_2\left[1-(x_1+x_2)\right]\leq (x_1+x_2)^2\left[1-(x_1+x_2)\right],$ \\

$\because 1-(x_1+x_2)\geq 0,\text{即证}\quad 4x_1x_2\leq (x_1+x_2)^2\quad\text{易得上式成立},$ \\

$\text{下面使用向前向后证明法},$ \\

$\text{向前:}\quad n\rightarrow 2n,$ \\

$\begin{cases}
    &\dfrac{\displaystyle\prod_{i=1}^n x_i}{(\displaystyle\sum_{i=1}^n x_i)^n}\leq
\dfrac{\displaystyle\prod_{i=1}^n (1-x_i)}{[\displaystyle\sum_{i=1}^n (1-x_i)]^n} \text{\ding{172}} \\
    &\dfrac{\displaystyle\prod_{i=n+1}^{2n} x_i}{(\displaystyle\sum_{i=n+1}^{2n} x_i)^n}\leq
\dfrac{\displaystyle\prod_{i=n+1}^{2n} (1-x_i)}{[\displaystyle\sum_{i=n+1}^{2n} (1-x_i)]^n} \text{\ding{173}} \\
\end{cases}$
\[\resizebox{0.70\textwidth}{!}{$\text{\ding{172}}\times\text{\ding{173}}:\dfrac{\displaystyle\prod_{i=1}^{2n} x_i}{(\displaystyle\sum_{i=1}^n x_i\cdot\sum_{i=n+1}^{2n} x_i)^n}\leq\dfrac{\displaystyle\prod_{i=1}^{2n} (1-x_i)}{[\displaystyle\sum_{i=1}^n (1-x_i)\cdot\sum_{i=n+1}^{2n} (1-x_i)]^n}$}\] \\
$\text{设}A=\displaystyle\sum_{i=1}^n x_i\quad B=\displaystyle\sum_{i=n+1}^{2n} x_i\quad\text{则}A\leq\frac{n}{2},B\leq\frac{n}{2},$ \\

$\therefore\dfrac{\displaystyle\prod_{i=1}^{2n} x_i}{\displaystyle\prod_{i=1}^{2n} (1-x_i)}\leq\dfrac{(AB)^n}{(n-A)^n(n-B)^n},$ $\text{令}A'=\dfrac{A}{n},B'=\dfrac{B}{n},$ \\

$\therefore\dfrac{(AB)^n}{(n-A)^n(n-B)^n}=\left[\dfrac{A'B'}{(1-A')(1-B')}\right]^n\quad \text{由}n=2\text{情形可知},$ \\

$\left[\dfrac{A'B'}{(1-A')(1-B')}\right]^n\leq\dfrac{(A'+B')^{2n}}{(2-A'-B')^{2n}}=\dfrac{(A+B)^{2n}}{(2n-A-B)^{2n}},$ \\

即$\dfrac{\displaystyle\prod_{i=1}^{2n} x_i}{\displaystyle\prod_{i=1}^{2n} (1-x_i)}\leq\dfrac{(\displaystyle\sum_{i=1}^{2n} x_i)^{2n}}{(2n-\displaystyle\sum_{i=1}^{2n} x_i)^{2n}}=\dfrac{(\displaystyle\sum_{i=1}^{2n} x_i)^{2n}}{[\displaystyle\sum_{i=1}^{2n} (1-x_i)]^{2n}},$ \\
 
向后:$n\rightarrow n-1$,  假设对$n$成立, $\because\dfrac{1}{n-1}\displaystyle\sum_{i=1}^{n-1} x_i\leq\dfrac{1}{2}\quad\text{由}x_n\text{的任意性},$ \\

不妨取$x_n=\dfrac{1}{n-1}\displaystyle\sum_{i=1}^{n-1} x_i,$ \\

设$S_n=\displaystyle\sum_{i=1}^{n} x_i,\quad\text{则}x_n=\dfrac{S_{n-1}}{n-1},\quad S_n=S_{n-1}+x_n=S_{n-1}\cdot\dfrac{n}{n-1},$ \\

$\therefore\dfrac{\displaystyle\prod_{i=1}^n x_i}{(S_n)^n}\leq\dfrac{\displaystyle\prod_{i=1}^{n} (1-x_i)}{(n-S_n)^n},$ \\

$\therefore\dfrac{\displaystyle\prod_{i=1}^{n-1} x_i\cdot\dfrac{S_{n-1}}{n-1}}{S_{n-1}^n\cdot\frac{n^n}{(n-1)^n}}\leq\dfrac{\displaystyle\prod_{i=1}^{n-1} (1-x_i)\cdot(1-\dfrac{S_{n-1}}{n-1})}{n^n(1-\dfrac{S_{n-1}}{n-1})^n},$ \\

$\therefore\dfrac{\displaystyle\prod_{i=1}^{n-1} x_i}{(\dfrac{S_{n-1}}{n-1})^{n-1}}\leq\dfrac{\displaystyle\prod_{i=1}^{n-1} (1-x_i)}{(1-\dfrac{S_{n-1}}{n-1})^{n-1}},$ \\

$\therefore\dfrac{\displaystyle\prod_{i=1}^{n-1} x_i}{(\displaystyle\sum_{i=1}^{n-1} x_i)^{n-1}}\leq\dfrac{\displaystyle\prod_{i=1}^{n-1} (1-x_i)}{\left[\displaystyle\sum_{i=1}^{n-1} (1-x_i)\right]^{n-1}},$ \\

又$n=2$时命题成立, \\
由向前-向后归纳法可知 $\dfrac{\displaystyle\prod_{i=1}^{n} x_i}{(\displaystyle\sum_{i=1}^{n} x_i)^n}\leq\dfrac{\displaystyle\prod_{i=1}^{n} (1-x_i)}{\left[\displaystyle\sum_{i=1}^{n} (1-x_i)\right]^n}.$

\begin{ti}
【题4】设$a_n=\sqrt{1+\sqrt{2+\sqrt{...+\sqrt{n}}}}\quad $证明:$\{a_n\}$收敛
\end{ti}

\textbf{证明}: 易得$\{a_n\}$单调递增, \\

对固定的$\displaystyle n$,定义$\displaystyle x_{n,n}=\sqrt n,\quad x_{k,n}=\sqrt{k+x_{k+1,n}}\quad(k=n-1,\ldots,1)$,则$\displaystyle a_n=x_{1,n},$ \\

下面用反向归纳法证明$\displaystyle x_{k,n}\leq\sqrt{k}+1\quad(1\leq k\leq n)$, \\

当$\displaystyle k=n$时,$\displaystyle x_{n,n}=\sqrt n\leq\sqrt n+1$.若$\displaystyle x_{k+1,n}\leq\sqrt{k+1}+1$, \\

则由$\displaystyle\sqrt{k+1}\leq2\sqrt k$可得$\displaystyle x_{k,n}\leq\sqrt{k+\sqrt{k+1}+1}\leq\sqrt{k+2\sqrt k+1}=\sqrt k+1,$ \\

故由反向归纳法,$\displaystyle a_n=x_{1,n}\leq2$,即$\displaystyle\{a_n\}$有上界.\\

又$\displaystyle\{a_n\}$单调递增,所以由单调收敛定理,$\displaystyle\{a_n\}$收敛. \\

\begin{ti}
【题5】设已知$\displaystyle\lim_{n\to\infty}a_n=a\quad$,证明$\displaystyle\lim_{n\to\infty}\frac{1}{2^n}\sum_{k=0}^{n}\binom{n}{k}a_k=a $
\end{ti}

\textbf{证明}:$\;\because 2^n=\displaystyle\sum_{k=0}^{n}\binom{n}{k},$ \\

$\therefore\left\lvert \dfrac{1}{2^n}\displaystyle\sum_{k=0}^{n}\binom{n}{k}a_k-a \right\rvert =\left\lvert \dfrac{1}{2^n}\displaystyle\sum_{k=0}^{n}\binom{n}{k}(a_k-a) \right\rvert$ $\leq\dfrac{1}{2^n}\displaystyle\sum_{k=0}^{n}\binom{n}{k}\left\lvert a_k-a\right\rvert,$ \\

$\because\displaystyle\lim_{n\to\infty}a_n=a\;\therefore\text{对}\forall\varepsilon>0,\;\exists N_1\;\text{s.t.}\;n>N_1\text{时},\;\left\lvert a_n-a\right\rvert<\varepsilon,$ \\

$\therefore\dfrac{1}{2^n}\displaystyle\sum_{k=0}^{n}\binom{n}{k}\left\lvert a_k-a\right\rvert=\dfrac{1}{2^n}\displaystyle\sum_{k=0}^{N_1}\binom{n}{k}\left\lvert a_k-a\right\rvert+\dfrac{1}{2^n}\displaystyle\sum_{k=N_1+1}^{n}\binom{n}{k}\left\lvert a_k-a\right\rvert$ \\

$<\dfrac{1}{2^n}\displaystyle\sum_{k=0}^{N_1}\binom{n}{k}\left\lvert a_k-a\right\rvert+\dfrac{\varepsilon}{2^n}\displaystyle\sum_{k=N_1+1}^{n}\binom{n}{k}$ \\

$<\dfrac{1}{2^n}\displaystyle\sum_{k=0}^{N_1}\binom{n}{k}\left\lvert a_k-a\right\rvert+\varepsilon,$ \\

又$\displaystyle\binom{n}{k} =\dfrac{n!}{k!(n-k)!}\leq\dfrac{n!}{(n-k)!}\leq n^k,\;\text{令}M=\max\{\left\lvert a_0-a\right\rvert,\left\lvert a_1-a\right\rvert,...,\left\lvert a_{N_1}-a\right\rvert  \},$ \\

$\therefore\dfrac{1}{2^n}\displaystyle\sum_{k=0}^{N_1}\binom{n}{k}\left\lvert a_k-a\right\rvert\leq\dfrac{M}{2^n}\displaystyle\sum_{k=0}^{N_1}n^k<\dfrac{M\cdot n^{N_1+1}}{2^n},$ \\

$\because\displaystyle\lim_{n\to\infty}\dfrac{n^{\alpha}}{2^n}=0\;\therefore\forall\varepsilon,\;\exists N_2,\;\text{s.t.}\;n>N_2\text{时},\dfrac{n^{N_1+1}}{2^n}<\varepsilon,$ \\

$\therefore\forall\varepsilon>0,\;\exists N=\max\{N_1,N_2\},\; n>N\text{时},$ $\left\lvert \dfrac{1}{2^n}\displaystyle\sum_{k=0}^{n}\binom{n}{k}a_k-a \right\rvert<(M+1)\varepsilon,$ \\

$\therefore\displaystyle\lim_{n\to\infty}\frac{1}{2^n}\sum_{k=0}^{n}\binom{n}{k}a_k=a.$

\begin{ti}
【题6】设${a_n}$满足条件$0<a_1<1,a_{n+1}=a_n(1-a_n)$,证明:$\displaystyle\lim_{n\to\infty}na_n=1$
\end{ti}

\textbf{证明}:$\;\because a_{n+1}-a_n=-a_n^2\leq 0,\quad\therefore \{a_n\}$单调递减, \\

由数学归纳法可得$0<a_n<1,\quad\therefore \{a_n\}$收敛, \\

设$\displaystyle\lim_{n\to\infty}a_n=\alpha,\text{其中}0\leq\alpha\leq 1,$ $\quad\therefore\alpha=\alpha(1-\alpha)\Rightarrow \alpha=0,$ \\

又$\dfrac{1}{a_{n+1}}=\dfrac{1}{a_n(1-a_n)}=\dfrac{1}{a_n}+\dfrac{1}{1-a_n},\quad\therefore\dfrac{1}{a_{n+1}}-\dfrac{1}{a_n}=\dfrac{1}{1-a_n},$ \\

$\therefore\displaystyle\lim_{n\to\infty}\dfrac{(n+1)-n}{\dfrac{1}{a_{n+1}}-\dfrac{1}{a_n}}=\displaystyle\lim_{n\to\infty}(1-a_n)=1,$ \\

又$\dfrac{1}{a_n}$为严格单调递增的正无穷大量, \\

$\therefore\displaystyle\lim_{n\to\infty}\dfrac{n}{\dfrac{1}{a_n}}=1,\;\text{即}\displaystyle\lim_{n\to\infty}na_n=1.$ \\

\textbf{注}:本题与题2类似,直接用Stolz公式并不好做,取倒数的动机为如果$\displaystyle\lim_{n\to\infty}na_n=1$ 

这就说明$a_n$的衰减速度与$\dfrac{1}{n}$相同,则$b_n=\dfrac{1}{a_n}$就是线性增长,便于研究.

\begin{ti}
【题7】证明:自然对数的底$e$是无理数
\end{ti}

\textbf{证明}:$\;\text{假设}e\in\mathbb{Q},\text{则}\exists\;p,q\in\mathbb{N}^+,\;\text{s.t.}\,e=\dfrac{p}{q},$ \\

设$\varepsilon_n=e-(1+\dfrac{1}{1!}+\dfrac{1}{2!}+...+\dfrac{1}{n!}),$ \\

先证明一个引理:$\;\varepsilon_n<\dfrac{1}{n!n},$ \\

$n!\varepsilon_n=n!\displaystyle\sum_{k=n+1}^{\infty}\dfrac{1}{k!}=\dfrac{1}{n+1}+\dfrac{1}{(n+1)(n+2)}+...,$ \\

$\dfrac{1}{n}-\dfrac{1}{n+1}=\dfrac{1}{n(n+1)},\dfrac{1}{n(n+1)}-\dfrac{1}{(n+1)(n+2)}=\dfrac{2}{n(n+1)(n+2)}...,$  \\

这样归纳下去,可得$\forall m\in\mathbb{N}^+ ,\dfrac{1}{n}>\dfrac{1}{n+1}+\dfrac{1}{(n+1)(n+2)}+...+\dfrac{1}{(n+1)(n+2)...(n+m)}\rlap{,}$ \\

从而$\varepsilon_n<\dfrac{1}{n!n},$\\

$\because q\in\mathbb{N}^+,\,e=(1+\dfrac{1}{1!}+...+\dfrac{1}{q!})+\varepsilon_q=\dfrac{p}{q},$ \\

两边同乘$\displaystyle q!$,有$\displaystyle p(q-1)!-q!\sum_{k=0}^{q}\dfrac{1}{k!}=q!\varepsilon_q,$ \\

等式左边为整数,故$\displaystyle q!\varepsilon_q\in\mathbb{Z}$.另一方面,$\displaystyle 0<q!\varepsilon_q<\dfrac{1}{q}\leq1$,即$\displaystyle 0<q!\varepsilon_q<1$,矛盾!

这就说明$e$是无理数,得证.

\begin{ti}
【题8】设已知存在极限$\displaystyle\lim_{n\to\infty}\dfrac{a_1+a_2+...+a_n}{n},$证明:$\displaystyle\lim_{n \to \infty}\dfrac{a_n}{n}=0  $
\end{ti}

\textbf{证明}: 设$b_n=\dfrac{a_1+a_2+...+a_n}{n},\text{由题意,不妨设}\{b_n\}\text{的极限为}l,$ \\

$\because a_n=nb_n-(n-1)b_{n-1}\,(n\geq 2),$ \\

$\therefore\displaystyle\lim_{n \to \infty}\dfrac{a_n}{n}=\lim_{n\to\infty}\dfrac{nb_n-(n-1)b_{n-1}}{n}=l-l=0.$

\begin{ti}
【题9】 设$x_n=\dfrac{1}{n^2}\displaystyle\sum_{k=0}^{n}ln\binom{n}{k},\text{求}\lim_{n\to\infty}x_n$
\end{ti}

\textbf{证明}: 由于$n^2$是单调递增的正无穷大量, \\

由Stolz公式,$\displaystyle\lim_{n\to\infty}x_n=\lim_{n\to\infty}\dfrac{\displaystyle\sum_{k=0}^{n+1}ln\binom{n+1}{k}-\displaystyle\sum_{k=0}^{n}ln\binom{n}{k}}{2n+1}$ \\

$=\displaystyle\lim_{n\to\infty}\dfrac{\displaystyle\sum_{k = 0}^{n}ln\frac{n+1}{n+1-k}}{2n+1},$又由于$2n+1$也是单调递增的正无穷大量, \\

再由Stolz公式,$\displaystyle\lim_{n\to\infty}\dfrac{\displaystyle\sum_{k = 0}^{n}\ln\frac{n+1}{n+1-k}}{2n+1}=\lim_{n\to\infty}\dfrac{\displaystyle\sum_{k=0}^{n}\ln\frac{n+1}{n+1-k}-\displaystyle\sum_{k=0}^{n-1}\ln\frac{n}{n-k}}{2}$ \\

$=\displaystyle\lim_{n\to\infty}\dfrac{\displaystyle\sum_{k=1}^{n}\ln\frac{n+1}{n+1-k}-\displaystyle\sum_{k=0}^{n-1}\ln\frac{n}{n-k}}{2}$ \\

$=\displaystyle\lim_{n\to\infty}\dfrac{\displaystyle\sum_{k=0}^{n-1}\ln\frac{n+1}{n-k}-\displaystyle\sum_{k=0}^{n-1}\ln\frac{n}{n-k}}{2}$ \\

$=\displaystyle\lim_{n\to\infty}\dfrac{n\ln\dfrac{n+1}{n}}{2},$ \\

$\because\dfrac{1}{n+1}\leq\ln(1+\dfrac{1}{n})\leq\dfrac{1}{n},\quad\therefore\dfrac{n}{n+1}\leq n\ln(1+\dfrac{1}{n})\leq 1,$ \\

由夹逼准则,$\displaystyle\lim_{n\to\infty}n\ln(1+\dfrac{1}{n})=1,$ \\

$\therefore\displaystyle\lim_{n\to\infty}x_n=\dfrac{1}{2}.$

\begin{ti}
【题10】设$a_n>0,\{\dfrac{a_{n+1}+a_{n+2}}{a_n}\}\text{收敛于}\alpha,\alpha>2$,证明:数列$\{a_n\}$没有上界
\end{ti}

\textbf{证明}:用反证法,假设$\{a_n\}$有上界,又$a_n>0\Rightarrow\{a_n\}$为有界序列, \\

所以$\{a_n\}$上下极限存在,记$a=\displaystyle\varlimsup_{n\to\infty}a_n,b=\displaystyle\varliminf_{n\to\infty}a_n,\text{那么}a\geq b\geq 0,$ \\

\ding{172}$\;a>0$:$\;\exists\,\beta>2,N_1,\;\text{s.t.}\,n>N_1\text{时},a_{n+1}+a_{n+2}>\beta a_n,$ \\

由上极限的保不等式性质,$2a\geq\beta a\Rightarrow\beta\leq 2,$矛盾!

\ding{173}$\;a=0$:同上有$a_{n+1}+a_{n+2}>\beta a_n(\beta>2,n>N_1),$ \\

记$M_n=\displaystyle\sup_{k\geq n}a_k,\text{则}a_{n+1}\leq M_n,a_{n+2}\leq M_n\Rightarrow M_n>\dfrac{\beta}{2}a_n>a_n,$ \\

又$M_n=\max\{a_n,M_{n+1}\},\;\therefore\,M_n=M_{n+1},\forall\,n>N_1,$ \\

又$a=0\Rightarrow\displaystyle\lim_{n\to\infty}M_n=0\;\therefore\forall\,n>N_1,\text{都有}M_n=0$,这与$a_n>0$矛盾! \\

综上,$\{a_n\}$没有上界. \\

\textbf{注意}:上下极限没有乘除法的四则运算性质,所以即使是在$a>b>0$的情况下,\\
$\alpha\neq\dfrac{2a}{b},\alpha\neq\dfrac{2b}{a}$ 



\BookEnd

\ifdefined\BookMaster
\def\BookEnd{}
\else
\documentclass[a4paper,12pt,fleqn,openany,twoside,fontset=windows]{ctexbook}

\usepackage[fleqn]{amsmath}
\usepackage{amssymb,amsthm,mathrsfs}
\usepackage{geometry}
\usepackage[dvipsnames]{xcolor}
\usepackage{graphicx}
\usepackage{tikz}
\usetikzlibrary{calc}
\usepackage[most]{tcolorbox}
\usepackage{fancyhdr}
\usepackage{titlesec}
\usepackage{tocloft}
\usepackage{enumitem}
\usepackage{microtype}
\usepackage{pifont}
\usepackage{hyperref}
\usepackage{romannum}
\usepackage{mathtools}

\definecolor{bookblue}{HTML}{264653}
\definecolor{bookteal}{HTML}{4F7C82}
\definecolor{booklight}{HTML}{EDF4F3}
\definecolor{bookgray}{HTML}{5E686B}

\geometry{
    inner=2.7cm,
    outer=2.5cm,
    top=2.7cm,
    bottom=2.7cm,
    headheight=18pt,
    headsep=18pt,
    footskip=25pt
}
\setlength{\mathindent}{0pt}
\setlength{\parindent}{5pt}
\setlength{\parskip}{3pt plus 1pt minus 1pt}
\linespread{1.25}
\allowdisplaybreaks[3]
\AtBeginDocument{%
    \setlength{\parindent}{5pt}
    \setlength{\abovedisplayskip}{8pt plus 2pt minus 2pt}
    \setlength{\belowdisplayskip}{8pt plus 2pt minus 2pt}
    \setlength{\abovedisplayshortskip}{5pt plus 1pt minus 1pt}
    \setlength{\belowdisplayshortskip}{5pt plus 1pt minus 1pt}
}

\hypersetup{
    unicode=true,
    colorlinks=true,
    linkcolor=bookblue,
    citecolor=bookblue,
    urlcolor=bookteal,
    bookmarksopen=true,
    pdfauthor={Chen},
    pdftitle={数学分析习题整理}
}

\renewcommand{\chaptermark}[1]{%
    \edef\@tmp{第\thechapter 章\quad #1}%
    \expandafter\markboth\expandafter{\@tmp}{}%
}
\fancypagestyle{main}{%
    \fancyhf{}
    \fancyhead[LE]{\small\color{bookgray}\nouppercase{\leftmark}}
    \fancyhead[RO]{\small\color{bookgray}数学分析习题整理}
    \fancyfoot[C]{\color{bookblue}\thepage}
    \fancyfoot[R]{\small\bfseries Chen\;\; github.com/mathlover-1/math-analysis}
    \renewcommand{\headrulewidth}{0.4pt}
    \renewcommand{\footrulewidth}{0pt}
}
\fancypagestyle{plain}{%
    \fancyhf{}
    \fancyfoot[C]{\color{bookblue}\thepage}
    \fancyfoot[R]{\small\bfseries Chen\;\; github.com/mathlover-1/math-analysis}
    \renewcommand{\headrulewidth}{0pt}
    \renewcommand{\footrulewidth}{0pt}
}
\pagestyle{main}

\titleformat{\chapter}[display]
  {\normalfont\sffamily\bfseries\color{bookblue}}
  {\raggedleft\fontsize{46}{50}\selectfont\thechapter}
  {-18pt}
  {\titlerule[1.2pt]\vspace{10pt}\filright\Huge}
\titlespacing*{\chapter}{0pt}{-15pt}{36pt}

\renewcommand{\contentsname}{目\quad 录}
\renewcommand{\cfttoctitlefont}{\hfill\sffamily\bfseries\Huge\color{bookblue}}
\renewcommand{\cftaftertoctitle}{\hfill}
\setlength{\cftbeforetoctitleskip}{5pt}
\setlength{\cftaftertoctitleskip}{28pt}
\setlength{\cftbeforechapskip}{12pt}
\renewcommand{\cftchapfont}{\sffamily\bfseries\large\color{bookblue}}
\renewcommand{\cftchappagefont}{\sffamily\bfseries\large\color{bookblue}}
\renewcommand{\cftchapleader}{\color{bookteal}\cftdotfill{2.5}}

\newtcolorbox{ti}{
    enhanced,
    breakable,
    colback=booklight!55!white,
    colframe=bookteal!35!white,
    boxrule=0.45pt,
    boxsep=0pt,
    left=10pt,
    right=10pt,
    top=9pt,
    bottom=9pt,
    before skip=14pt,
    after skip=14pt,
    arc=2pt,
    outer arc=2pt,
    before upper={\parindent=0pt}
}

\newenvironment{booknotebody}{%
    \begin{center}
    \begin{minipage}{0.84\textwidth}
    \songti\zihao{-4}\linespread{1.45}\selectfont
    \setlength{\parindent}{2\ccwd}
    \setlength{\parskip}{0.75em}
}{%
    \end{minipage}
    \end{center}
}

\newenvironment{booknotelongbody}{%
    \begingroup
    \songti\zihao{-4}\linespread{1.45}\selectfont
    \setlength{\parindent}{2\ccwd}
    \setlength{\parskip}{0.75em}
    \begin{list}{}{%
        \setlength{\leftmargin}{0.08\textwidth}
        \setlength{\rightmargin}{0.08\textwidth}
        \setlength{\topsep}{0pt}
        \setlength{\partopsep}{0pt}
        \setlength{\parsep}{0pt}
        \setlength{\itemsep}{0pt}
    }
    \item\relax
}{%
    \end{list}
    \endgroup
}

\fancypagestyle{plainnowm}{%
    \fancyhf{}
    \fancyfoot[C]{\color{bookblue}\thepage}
    \renewcommand{\headrulewidth}{0pt}
    \renewcommand{\footrulewidth}{0pt}
}

\newcommand{\booknotetitle}[2][plain]{%
    \clearpage
    \phantomsection
    \addcontentsline{toc}{chapter}{#2}
    \markboth{#2}{}
    \thispagestyle{#1}
    \vspace*{1.2cm}
    \begin{center}
        {\sffamily\heiti\bfseries\fontsize{26}{34}\selectfont\color{bookblue}#2\par}
        \vspace{0.8em}
        {\color{bookteal}\rule{4.5cm}{0.8pt}\par}
    \end{center}
    \vspace{0.6cm}
}

\renewcommand{\proofname}{证明}
\renewcommand{\qedsymbol}{\color{bookteal}\ensuremath{\blacksquare}}

\title{数学分析习题整理}
\author{Chen}
\date{}

\makeatletter
% ==================== 旧版封面/封底设计备份 ====================
% 2026-08 重设计,旧代码用 \iffalse...\fi 封存,新版见下方。
\iffalse
\renewcommand{\maketitle}{%
    \begin{titlepage}
        \thispagestyle{empty}
        \setcounter{page}{0}
        \noindent\hspace*{-\parindent}\begin{tikzpicture}[remember picture,overlay]
            \fill[bookblue] (current page.north west) rectangle ([yshift=-4.2cm]current page.north east);
            \fill[booklight] ([yshift=-4.2cm]current page.north west) rectangle (current page.south east);
            \draw[bookteal!55,line width=0.7pt]
                ([xshift=2.4cm,yshift=4.3cm]current page.south west)
                .. controls +(2.5cm,1.8cm) and +(-2.8cm,-1.3cm) ..
                ([xshift=-2.4cm,yshift=7.2cm]current page.south east);
            \draw[bookteal!35,line width=0.7pt]
                ([xshift=2.5cm,yshift=6.2cm]current page.south west)
                .. controls +(3.0cm,-1.2cm) and +(-3.0cm,1.6cm) ..
                ([xshift=-2.3cm,yshift=3.8cm]current page.south east);
        \end{tikzpicture}
        \vspace*{5.6cm}
        \noindent\color{bookblue}\rule{\textwidth}{1.4pt}\par
        \vspace{8pt}
        {\centering\sffamily\bfseries\fontsize{32}{42}\selectfont\color{bookblue}\@title\par}
        \vspace{8pt}
        \noindent\color{bookteal}\rule{\textwidth}{0.6pt}\par
        \vspace{2.6cm}
        {\centering\Large\color{bookgray}\@author\par}
        \vfill
        {\centering\color{bookteal}\large
            $\displaystyle \lim_{n\to\infty} a_n$\qquad
            $\displaystyle \int_a^b f(x)\,\mathrm{d}x$\qquad
            $\displaystyle \nabla f(x,y)$\par
        }
        \vspace*{1.5cm}
    \end{titlepage}
}

\newcommand{\makebackcover}{%
    \begingroup
    \let\cleardoublepage\clearpage
    \begin{titlepage}
        \thispagestyle{empty}
        \noindent\hspace*{-\parindent}\begin{tikzpicture}[remember picture,overlay]
            \fill[bookblue] (current page.north west) rectangle (current page.south east);
            \fill[booklight,rounded corners=7pt]
                ([xshift=2.2cm,yshift=-3.0cm]current page.north west)
                rectangle
                ([xshift=-2.2cm,yshift=3.0cm]current page.south east);
            \fill[bookteal!80]
                (current page.north west) --
                ([xshift=5.3cm]current page.north west) --
                ([xshift=2.1cm,yshift=-4.4cm]current page.north west) --
                ([yshift=-6.2cm]current page.north west) -- cycle;
            \fill[bookteal!55]
                (current page.south east) --
                ([xshift=-5.8cm]current page.south east) --
                ([xshift=-2.4cm,yshift=4.2cm]current page.south east) --
                ([yshift=6.0cm]current page.south east) -- cycle;
            \foreach \radius in {0.75,1.15,1.55,1.95} {
                \draw[bookteal!55,line width=0.55pt]
                    ([xshift=-2.9cm,yshift=-4.4cm]current page.north east)
                    circle[radius=\radius cm];
            }
            \foreach \radius in {0.65,1.05,1.45} {
                \draw[booklight!65,line width=0.55pt]
                    ([xshift=2.4cm,yshift=3.6cm]current page.south west)
                    circle[radius=\radius cm];
            }
            \fill[bookteal!16,rounded corners=3pt]
                ([xshift=3.5cm,yshift=-7.1cm]current page.north west)
                rectangle
                ([xshift=-5.0cm,yshift=7.2cm]current page.south east);
            \fill[bookteal!30,rounded corners=3pt]
                ([xshift=4.4cm,yshift=-8.0cm]current page.north west)
                rectangle
                ([xshift=-4.1cm,yshift=8.1cm]current page.south east);
            \begin{scope}[shift={([xshift=-0.75cm]current page.center)}]
                \clip[rounded corners=3pt] (-5.25,-4.75) rectangle (5.25,4.75);
                \foreach \v in {-3.6,-3.0,...,3.6} {
                    \draw[bookteal!48,line width=0.45pt,opacity=0.72]
                        plot[domain=-3.6:3.6,samples=45,smooth,variable=\u]
                        ({\u+0.55*\v},{0.28*\u-0.35*\v+0.11*(\u*\u-\v*\v)});
                }
                \foreach \u in {-3.6,-3.0,...,3.6} {
                    \draw[bookblue!42,line width=0.42pt,opacity=0.62]
                        plot[domain=-3.6:3.6,samples=45,smooth,variable=\v]
                        ({\u+0.55*\v},{0.28*\u-0.35*\v+0.11*(\u*\u-\v*\v)});
                }
                \draw[bookteal!70,line width=0.75pt,opacity=0.82]
                    plot[domain=-3.6:3.6,samples=55,smooth,variable=\u]
                    ({\u},{0.28*\u+0.11*\u*\u});
                \draw[bookblue!60,line width=0.7pt,opacity=0.75]
                    plot[domain=-3.6:3.6,samples=55,smooth,variable=\v]
                    ({0.55*\v},{-0.35*\v-0.11*\v*\v});
            \end{scope}
            \draw[bookteal!45,line width=0.65pt]
                ([xshift=3.0cm,yshift=-8.2cm]current page.north west)
                .. controls +(2.2cm,2.0cm) and +(-2.0cm,-2.1cm) ..
                ([xshift=-3.0cm,yshift=9.0cm]current page.south east);
            \draw[bookteal!25,line width=0.55pt]
                ([xshift=3.1cm,yshift=8.0cm]current page.south west)
                .. controls +(2.7cm,-1.7cm) and +(-2.6cm,1.8cm) ..
                ([xshift=-3.1cm,yshift=-8.7cm]current page.north east);
            \node[anchor=west,text=bookteal!72,scale=0.78] at
                ([xshift=2.8cm,yshift=-4.0cm]current page.north west) {%
                $\displaystyle e^{\mathrm{i}\pi}+1=0$};
            \node[anchor=east,text=bookblue!62,scale=0.67] at
                ([xshift=-2.7cm,yshift=4.0cm]current page.south east) {%
                $\displaystyle
                f(x)=\sum_{n=0}^{\infty}\frac{f^{(n)}(a)}{n!}(x-a)^n$};
        \end{tikzpicture}
        \mbox{}
    \end{titlepage}
    \endgroup
}
\fi

% ==================== 新版封面(2026-08 重设计) ====================
% 设计要点:顶部深绿纵向渐变+波浪下缘+两层半透明过渡波(解决深浅绿生硬过渡);
% 深色区螺旋线/同心圆母题;标题+英文副题+作者居中;
% 中部 Fourier 部分和逼近方波曲线;底部浅色波纹(无文字)。
\renewcommand{\maketitle}{%
    \begin{titlepage}
        \thispagestyle{empty}
        \setcounter{page}{0}
        \noindent\hspace*{-\parindent}\begin{tikzpicture}[remember picture,overlay]
            % ---------- 底色 ----------
            \fill[booklight] (current page.north west) rectangle (current page.south east);

            % ---------- 顶部深色区:纵向渐变 + 波浪下缘 ----------
            \shade[top color=bookblue, bottom color=bookteal!55!bookblue]
                (current page.north west) -- (current page.north east) --
                ([yshift=-8.6cm]current page.north east)
                .. controls ([xshift=-5.2cm,yshift=-10.6cm]current page.north east)
                         and ([xshift=6.0cm,yshift=-8.0cm]current page.north west) ..
                ([yshift=-10.0cm]current page.north west) -- cycle;

            % ---------- 过渡层 1:半透明波浪 ----------
            \fill[bookteal!60!bookblue,opacity=0.42]
                ([yshift=-8.6cm]current page.north east)
                .. controls ([xshift=-5.2cm,yshift=-10.6cm]current page.north east)
                         and ([xshift=6.0cm,yshift=-8.0cm]current page.north west) ..
                ([yshift=-10.0cm]current page.north west) --
                ([yshift=-12.6cm]current page.north west)
                .. controls ([xshift=6.5cm,yshift=-11.0cm]current page.north west)
                         and ([xshift=-5.5cm,yshift=-13.2cm]current page.north east) ..
                ([yshift=-11.4cm]current page.north east) -- cycle;

            % ---------- 过渡层 2:更浅的波浪 ----------
            \fill[bookteal!35,opacity=0.35]
                ([yshift=-10.6cm]current page.north east)
                .. controls ([xshift=-6.0cm,yshift=-12.4cm]current page.north east)
                         and ([xshift=5.0cm,yshift=-9.8cm]current page.north west) ..
                ([yshift=-11.8cm]current page.north west) --
                ([yshift=-13.8cm]current page.north west)
                .. controls ([xshift=5.5cm,yshift=-12.6cm]current page.north west)
                         and ([xshift=-6.5cm,yshift=-14.6cm]current page.north east) ..
                ([yshift=-12.8cm]current page.north east) -- cycle;

            % ---------- 深色区装饰:右上螺旋线 ----------
            \begin{scope}[shift={([xshift=-3.4cm,yshift=-4.3cm]current page.north east)}]
                \draw[booklight!45!bookteal,opacity=0.55,line width=0.6pt]
                    plot[domain=0:1440,samples=220,smooth,variable=\t]
                    ({0.0028*\t*cos(\t)},{0.0028*\t*sin(\t)});
            \end{scope}

            % ---------- 深色区装饰:左下同心圆 ----------
            \foreach \r in {0.7,1.25,1.8,2.35,2.9}{
                \draw[booklight!35!bookteal,opacity=0.5,line width=0.5pt]
                    ([xshift=0.4cm,yshift=-7.9cm]current page.north west) circle[radius=\r cm];
            }

            % ---------- 主标题 ----------
            \node[anchor=center,text=booklight] at
                ([yshift=-4.5cm]current page.north)
                {\sffamily\bfseries\fontsize{40}{46}\selectfont \@title};

            % ---------- 分隔短线 ----------
            \draw[booklight!65,line width=1.2pt]
                ([xshift=-2.6cm,yshift=-5.75cm]current page.north) --
                ([xshift=2.6cm,yshift=-5.75cm]current page.north);

            % ---------- 英文副题 ----------
            \node[anchor=center,text=booklight!70!bookteal] at
                ([yshift=-6.55cm]current page.north)
                {\itshape\large Mathematical Analysis \textperiodcentered\
                 A Collection of Exercises};

            % ---------- 作者(紧随标题下方) ----------
            \node[anchor=center,text=booklight!88] at
                ([yshift=-8.2cm]current page.north)
                {\sffamily\large \@author\quad 整理};

            % ---------- 中部:Fourier 部分和逼近方波 ----------
            \begin{scope}[shift={([xshift=-5.2cm,yshift=-18.1cm]current page.north)}]
                % 坐标轴
                \draw[bookgray!55,line width=0.5pt,->] (-0.5,0) -- (10.9,0);
                \draw[bookgray!55,line width=0.5pt,->] (0,-1.75) -- (0,1.95);
                % 目标方波(虚线)
                \draw[bookgray!75,line width=0.6pt,dashed]
                    (0,0.92) -- (5,0.92) (5,-0.92) -- (10,-0.92);
                % N=1
                \draw[bookteal!45,line width=0.7pt]
                    plot[domain=0:360,samples=90,smooth,variable=\x]
                    ({\x/36},{1.18*sin(\x)});
                % N=3
                \draw[bookteal!80,line width=0.8pt]
                    plot[domain=0:360,samples=120,smooth,variable=\x]
                    ({\x/36},{1.18*(sin(\x)+sin(3*\x)/3)});
                % N=7
                \draw[bookblue,line width=0.95pt]
                    plot[domain=0:360,samples=150,smooth,variable=\x]
                    ({\x/36},{1.18*(sin(\x)+sin(3*\x)/3+sin(5*\x)/5+sin(7*\x)/7)});
                % 刻度标注
                \node[anchor=north east,text=bookgray!80,scale=0.62] at (-0.12,0) {$0$};
                \node[anchor=north,text=bookgray!80,scale=0.62] at (5,0) {$\pi$};
                \node[anchor=north,text=bookgray!80,scale=0.62] at (10,0) {$2\pi$};
            \end{scope}
            % 图注
            \node[anchor=center,text=bookgray,scale=0.8] at
                ([yshift=-20.5cm]current page.north)
                {$S_N(x)=\displaystyle\sum_{k=1}^{N}\frac{\sin\bigl((2k-1)x\bigr)}{2k-1}$
                 \ 逐步逼近方波};

            % ---------- 底部浅色波纹(无文字,与顶部波浪呼应) ----------
            \fill[bookteal!22,opacity=0.5]
                ([yshift=2.4cm]current page.south west)
                .. controls ([xshift=6.0cm,yshift=3.9cm]current page.south west)
                         and ([xshift=-6.5cm,yshift=1.0cm]current page.south east) ..
                ([yshift=2.9cm]current page.south east) --
                (current page.south east) -- (current page.south west) -- cycle;
            \fill[bookteal!40!bookblue,opacity=0.28]
                ([yshift=1.3cm]current page.south west)
                .. controls ([xshift=5.5cm,yshift=2.7cm]current page.south west)
                         and ([xshift=-5.5cm,yshift=0.3cm]current page.south east) ..
                ([yshift=1.9cm]current page.south east) --
                (current page.south east) -- (current page.south west) -- cycle;
        \end{tikzpicture}
        \mbox{}
    \end{titlepage}
}

% ==================== 新版封底(纯背景,2026-08 重设计) ====================
% 设计要点:对角渐变底+角部半透明叠层;中下部外摆线(玫瑰形)为视觉主体;
% 底部两层半透明白色波浪;无任何文字。
\newcommand{\makebackcover}{%
    \begingroup
    \let\cleardoublepage\clearpage
    \begin{titlepage}
        \thispagestyle{empty}
        \noindent\hspace*{-\parindent}\begin{tikzpicture}[remember picture,overlay]
            % ---------- 底色:对角渐变 ----------
            \shade[top color=bookblue, bottom color=bookteal!48!bookblue,
                   shading angle=-30]
                (current page.north west) rectangle (current page.south east);

            % ---------- 角部半透明叠层 ----------
            \fill[white,opacity=0.05]
                ([xshift=2.0cm,yshift=1.0cm]current page.north east) circle[radius=7.2cm];
            \fill[white,opacity=0.06]
                ([xshift=-0.5cm,yshift=-1.5cm]current page.north east) circle[radius=3.6cm];

            % ---------- 左下同心圆 ----------
            \foreach \r in {1.0,1.75,2.5,3.25,4.0,4.75}{
                \draw[booklight!40,opacity=0.55,line width=0.55pt]
                    ([xshift=-0.6cm,yshift=-0.4cm]current page.south west) circle[radius=\r cm];
            }
            % ---------- 右下螺旋线 ----------
            \begin{scope}[shift={([xshift=-2.9cm,yshift=3.2cm]current page.south east)}]
                \draw[booklight!45!bookteal,opacity=0.5,line width=0.55pt]
                    plot[domain=0:1260,samples=200,smooth,variable=\t]
                    ({0.0036*\t*cos(\t)},{0.0036*\t*sin(\t)});
            \end{scope}
            % ---------- 左上斜线组 ----------
            \foreach \i in {0,1,2}{
                \draw[bookteal!60!booklight,opacity=0.5,line width=0.6pt]
                    ([xshift={1.6cm+\i*0.45cm}]current page.north west) --
                    ([yshift={-1.6cm-\i*0.45cm}]current page.north west);
            }

            % ---------- 中下部视觉主体:外摆线(玫瑰形) ----------
            \begin{scope}[shift={([yshift=-16.6cm]current page.north)}]
                % 外圈
                \draw[booklight!35,opacity=0.6,line width=0.6pt]
                    (0,0) circle[radius=4.9cm];
                \draw[booklight!22,opacity=0.5,line width=0.5pt]
                    (0,0) circle[radius=5.25cm];
                % 主玫瑰线
                \draw[booklight!55!bookteal,opacity=0.8,line width=0.85pt]
                    plot[domain=0:1080,samples=400,smooth,variable=\t]
                    ({0.335*(8*cos(\t)-5*cos(2.6667*\t))},
                     {0.335*(8*sin(\t)-5*sin(2.6667*\t))});
                % 内层玫瑰线
                \draw[booklight!38!bookteal,opacity=0.6,line width=0.6pt]
                    plot[domain=0:1080,samples=400,smooth,variable=\t]
                    ({0.21*(8*cos(\t)-5*cos(2.6667*\t))},
                     {0.21*(8*sin(\t)-5*sin(2.6667*\t))});
            \end{scope}

            % ---------- 底部半透明波浪 ----------
            \fill[white,opacity=0.06]
                ([yshift=4.6cm]current page.south west)
                .. controls ([xshift=6.5cm,yshift=6.4cm]current page.south west)
                         and ([xshift=-6.0cm,yshift=2.6cm]current page.south east) ..
                ([yshift=5.2cm]current page.south east) --
                (current page.south east) -- (current page.south west) -- cycle;
            \fill[booklight,opacity=0.08]
                ([yshift=2.6cm]current page.south west)
                .. controls ([xshift=5.5cm,yshift=4.3cm]current page.south west)
                         and ([xshift=-6.5cm,yshift=1.0cm]current page.south east) ..
                ([yshift=3.3cm]current page.south east) --
                (current page.south east) -- (current page.south west) -- cycle;
        \end{tikzpicture}
        \mbox{}
    \end{titlepage}
    \endgroup
}
\makeatother

% 页脚右下方常态化水印(署名 + 仓库地址,加粗黑色),已在 main/plain 分页样式中定义,所有页面均生效.

\begin{document}
\mainmatter
\pagestyle{main}
\setcounter{chapter}{1}
\def\BookEnd{\end{document}}
\fi

\chapter{实数系基本定理}
\begin{ti}
【题1】 设$A$是有上界的数集,$\beta=\sup A$  证明:$\exists\{a_n\},\text{s.t.}\;a_n\in A\text{且}\displaystyle\lim_{n\to\infty}a_n=\beta;$ \\

又若$\beta\notin A,$则$\{a_n\}$可以是严格单调递增的
\end{ti}

\textbf{证明}:$\;\because\beta=\sup A,\quad\therefore\forall\,n\geq 1,\text{都存在}a_n\in A,\;\text{s.t.}\;\beta-\dfrac{1}{n}<a_n\leq\beta,$ \\

由此生成的数列$\{a_n\}\text{就满足}\displaystyle\lim_{n\to\infty}a_n=\beta,$ \\

若$\beta\notin A,\text{则}a_n<\beta,\text{则}\beta-1<a_1<\beta,$ \\

$\because\max\{a_1,\beta-\dfrac{1}{2}\}<\beta\quad\therefore\exists a_2\in A,\;\text{s.t.}\;\max\{a_1,\beta-\dfrac{1}{2}\}<a_2<\beta,$ \\

由此归纳地做下去,可以找到数列$\{a_n\}$,满足$\beta-\dfrac{1}{n}<a_n<\beta\text{且}a_n>a_{n-1},$ \\

$\therefore\{a_n\}$可以是严格单调递增的.

\begin{ti}
【题2】 证明:每个数列都有单调子列
\end{ti}

\textbf{证明}:称$\{a_n\}$中某项$a_n$满足以下条件时具有性质$(M):\forall k>n,\text{都有}a_k\leq a_n,$ \\

分情况讨论, \\

\ding{172}$\quad\text{若}\{a_n\}\text{中有无限多项具有性质}(M),\text{则把这些下标拿出来,记为}n_1,n_2,...,n_k,$ \\

则有$a_{n_1}\geq a_{n_2}\geq...\geq a_{n_k}\geq...\quad\text{那么}\{a_{n_k}\}$为单调递减的子列, \\

\ding{173}$\quad\text{若}\{a_n\}\text{只有有限多项满足性质}(M),\text{则}\,\exists N,\;\text{s.t.}\;n>N\text{时},a_n\text{均不满足性质}(M),$ \\

取$n_1>N,\text{则}\exists\,n_2>n_1,\;\text{s.t.}\,a_{n_2}>a_{n_1},$ \\

同理继续操作下去,可得到单调递增的子列$\{a_{n_k}\},$ \\

综上,每个数列都有单调子列.

\begin{ti}
【题3】证明:若$\{x_n\}$无界,但不是无穷大量,则存在两个子列,其中一个子列收敛,另一个子列是无穷大量
\end{ti}

\textbf{证明}:先证明前半部分:若$\{x_n\}$无界,则有无穷大量子列, \\

$\because\{x_n\}\text{无界},\quad\therefore\forall\,M\in\mathbb{R},\exists\,n\in\mathbb{N}\;\text{s.t.}\left\lvert x_n\right\rvert>M,$ \\

事实上,这样的$n$总是有无限个, \\

否则,令$M'=\max\{\left\lvert x_n\right\rvert\,|\;\left\lvert x_n\right\rvert>M\},\text{因为只有有限项满足}\left\lvert x_n\right\rvert>M,\text{所以}M'\in\mathbb{R},$ \\

则$\forall\,n\in\mathbb{N}\;\text{总有}\left\lvert x_n\right\rvert\leq M',\text{与}\{x_n\}\text{无界矛盾},$ \\

$\therefore\text{取}M=1,\;\exists\,p_1\in\mathbb{N},\;\text{s.t.}\,\left\lvert x_{p_1}\right\rvert>1,$ \\

$\phantom{\therefore}\text{取}M=2,\;\exists\,p_2>p_1,\;\text{s.t.}\,\left\lvert x_{p_2}\right\rvert>2,$ \\

从而归纳地做下去,可以得到子列$\{x_{p_m}\}$,$\text{其中}p_1<p_2<...<p_m<...,\,\left\lvert x_{p_m}\right\rvert>m,$ \\

$\therefore\{x_{p_m}\}$是无穷大量的子列. \\

再证下半部分:若$\{x_n\}$不是无穷大量,则$\{x_n\}$有收敛子列, \\

$\because\{x_n\}\text{不是无穷大量},\therefore\exists\,G_0>0,\,\forall\,N,\,\exists\,n>N,\,\text{s.t.}\left\lvert x_n\right\rvert \leq G_0,$ \\

取$N=1,\;\exists\,n_1>1,\,\text{s.t.}\,\left\lvert x_{n_1}\right\rvert\leq G_0,$ \\

再取$N=n_1,\;\exists\,n_2>n_1,\,\text{s.t.}\,\left\lvert x_{n_2}\right\rvert\leq G_0,$ \\

由此归纳地做下去,我们可以得到子列$\{x_{n_k}\}$,$\text{其中}n_1<n_2<...<n_k<...,\,\left\lvert x_{n_k}\right\rvert\leq G_0,$ \\

$\therefore\{x_{n_k}\}$是有界数列, 运用$Bolzano$定理 ,$\{x_{n_k}\}$有收敛子列,  则$\{x_n\}$也有收敛子列, \\

综上,命题得证. \\

\textbf{另解}:对上半部分的另一种证法, \\

取$M=1\quad\exists\,p_1\in\mathbb{N},\;\text{s.t.}\,\left\lvert x_{p_1}\right\rvert>1,$ \\

再取$M=\max\{\left\lvert x_1\right\rvert,\left\lvert x_2\right\rvert,...,\left\lvert x_{p_1}\right\rvert,2\},\;\exists\,p_2,\;\text{s.t.}\,\left\lvert x_{p_2}\right\rvert>M,$ \\

则$p_2>p_1,\text{且}\left\lvert x_{p_2}\right\rvert>2,$ \\

这样归纳地做下去,同样可得到所需的无穷大量子列. \\

\textbf{注记}:$\,Bolzano$定理告诉我们有界数列必有收敛子列,而下半部分的证明事实上提供了$Bolzano$定理的加强版: $\{x_n\}\text{有收敛子列}\Leftrightarrow\{x_n\}\text{不是无穷大量}$

\begin{ti}
【题4】用闭区间套原理直接证明$Cauchy$收敛准则
\end{ti}

\textbf{证明}:设$\{x_n\}\text{是}Cauchy$列,则易得$\{x_n\}$有界, \\

$\text{设}x_n\in[a_1,b_1],\;\text{对}[a_1,b_1]\text{三等分,令}c_1=\dfrac{2a_1+b_1}{3},\;c_2=\dfrac{a_1+2b_1}{3},$ \\

设$J_1=[a_1,c_1],J_2=[c_1,c_2],J_3=[c_2,b_1],$ \\

$Claim:J_1\text{和}J_3$中至少有一个区间只包含$\{x_n\}$的有限多项, \\

事实上,如果$J_1\text{和}J_3$均包含$\{x_n\}$中的无限多项, \\

则$\exists\,\varepsilon=\dfrac{b_1-a_1}{3},\;\forall N,\;\exists\,m,n>N\;\text{s.t.}\,x_n\in J_1,x_m\in J_3,\text{且}\left\lvert x_m-x_n\right\rvert\geq\varepsilon,$ \\

这与$\{x_n\}\text{是}Cauchy$列矛盾! \\

从而可以去掉$J_1,J_3$中只包含有限项的那个区间(如果$J_1,J_3$都只包含有限多项,则任意去掉一个区间),与未去掉的$J_2$组成新的闭区间$[a_2,b_2],$ \\

如此继续操作下去,可以得到闭区间套$\{[a_k,b_k]\},\text{且}b_n-a_n=\dfrac{2}{3}(b_{n-1}-a_{n-1}),$ \\

$\therefore\displaystyle\lim_{n\to\infty}(b_n-a_n)=0,$ \\

由闭区间套原理,可知闭区间套有一个公共点$\xi$,且根据闭区间套的构造, \\

每个区间$[a_k,b_k]$均包含$\{x_n\}$中从某项开始往后的所有项, \\

$\because\displaystyle\lim_{n\to\infty}b_n=\lim_{n\to\infty}a_n=\xi,\;\therefore\forall\varepsilon>0,\;\exists N_1,\;\text{s.t.}\,[a_{N_1},b_{N_1}]\subseteq(\xi-\varepsilon,\xi+\varepsilon),$ \\

又$\exists\,N_2,\;n>N_2\text{时},x_n\in[a_{N_1},b_{N_1}],$ \\

$\therefore\,\forall\,\varepsilon>0,\;\exists\,N,\,n>N\text{时},\left\lvert x_n-\xi\right\rvert<\varepsilon\Rightarrow\{x_n\}\text{收敛}$,得证.

\begin{ti}
【题5】证明:有限覆盖定理在$\mathbb{Q}$中不成立
\end{ti}

\textbf{证明}:考虑$J=[0,2]\cap\mathbb{Q},$ \\

$\because\sqrt{2}\notin\mathbb{Q},\;\therefore\forall\,x\in J,\;\exists r_x\in\mathbb{Q}^+,\;\text{s.t.}\,\sqrt{2}\notin(x-r_x,x+r_x),$ \\

又$J\subseteq\displaystyle\bigcup_{x\in J}[(x-r_x,x+r_x)\cap\mathbb{Q}]$,如果有限覆盖定理对$J$成立, \\

即$J\subseteq\displaystyle\bigcup_{k=1}^{n}(x_k-r_k,x_k+r_k)\cap\mathbb{Q},$ \\

令$\delta_k=\min\{\left\lvert \sqrt{2}-x_k+r_k\right\rvert,\left\lvert \sqrt{2}-x_k-r_k\right\rvert\}$,再令$\delta=\min\{\delta_1,\delta_2,...,\delta_n\},$ \\

则$(\sqrt{2}-\delta,\sqrt{2}+\delta)\bigcap\left(\displaystyle\bigcup_{k=1}^{n}(x_k-r_k,x_k+r_k)\right)=\varnothing,$ \\

又$\exists\,q\in J,\text{且}q\in(\sqrt{2}-\delta,\sqrt{2}+\delta),$ \\

$\therefore\displaystyle\bigcup_{k=1}^{n}(x_k-r_k,x_k+r_k)$不能覆盖整个$J$,自然$\left(\displaystyle\bigcup_{k=1}^{n}(x_k-r_k,x_k+r_k)\right)\bigcap\mathbb{Q}$不能覆盖整个$J,$ \\

所以有限覆盖定理在$\mathbb{Q}$中不成立.

\begin{ti}
【题6】用有限覆盖定理证明闭区间套原理
\end{ti}

\textbf{证明}:设存在闭区间套$[a_1,b_1]\supseteq[a_2,b_2]\supseteq...\supseteq[a_n,b_n]\supseteq...,\text{且}\displaystyle\lim_{n\to\infty}(b_n-a_n)=0,$ \\

假设这个闭区间套无公共点,则$\forall\,x\in[a_1,b_1],\;\exists\,n_x,\;\text{s.t.}\,x\notin[a_{n_x},b_{n_x}],$ \\

构造集合$O_x=\mathbb{R}\backslash[a_{n_x},b_{n_x}],\text{则}x\in O_x,\text{且}O_x\text{是开集},$ \\

则$\displaystyle\bigcup_{x\in[a_1,b_1]}O_x$是$[a_1,b_1]$的一个开覆盖,则存在有限个$O_x$覆盖整个$[a_1,b_1],$ \\

设$[a_1,b_1]\subseteq\displaystyle\bigcup_{i=1}^{k}O_{x_i},N=\max\{n_{x_1},n_{x_2},...n_{x_k}\},$ \\

则$[a_N,b_N]\subseteq[a_{n_{x_i}},b_{n_{x_i}}],\text{从而}[a_N,b_N]\text{不与}O_{x_1},O_{x_2},...,O_{x_k}$中任意一个相交, \\

但$[a_N,b_N]\subseteq[a_1,b_1],\text{而}O_{x_1},O_{x_2},...,O_{x_k}\text{覆盖了整个}[a_1,b_1]$,矛盾! \\

$\therefore$闭区间套必有公共点, 设$\xi,\xi'$都是公共点,则$\left\lvert \xi-\xi'\right\rvert\leq\left\lvert b_n-a_n\right\rvert,$ \\

令$n\rightarrow\infty,\text{则}\left\lvert \xi-\xi'\right\rvert\leq 0\Rightarrow\xi=\xi',$ \\

$\therefore$闭区间套$\{[a_n,b_n]\}$有唯一公共点.

\begin{ti}
【题7】设正数列$\{a_n\}$满足$a_{m+n}\leq a_ma_n,\,\forall\,n,m\in\mathbb{N}^+$,证明: \\
若$\{\dfrac{\ln{a_n}}{n}\}$有界,则$\displaystyle\lim_{n\to\infty}\dfrac{\ln{a_n}}{n}=\displaystyle\inf_{n\geq 1}\{\dfrac{\ln{a_n}}{n}\}$
\end{ti}

\textbf{证明}:设$\alpha=\displaystyle\inf_{n\geq 1}\{\dfrac{\ln{a_n}}{n}\},\text{则}\displaystyle\varliminf_{n\to\infty}\dfrac{\ln{a_n}}{n}\geq\alpha,$ \\

又$\alpha$是下确界,由确界的刻画,$\forall\,\varepsilon>0,\;\exists\,N,\;\text{s.t.}\,\dfrac{\ln{a_N}}{N}<\alpha+\varepsilon,$ \\

固定这个$N,\text{对}\forall\,n\in\mathbb{N}^+,\;\exists\,m,k\in\mathbb{N},\;\text{s.t.}\,n=mN+k,0<k\leq N,n$足够大时$m\geq 1$, \\

则$a_n=a_{mN+k}\leq a_{mN}a_k\leq a_N^ma_k,$ \\

$\therefore\dfrac{\ln{a_n}}{n}\leq\dfrac{m\ln{a_N}}{n}+\dfrac{\ln{a_k}}{n}\leq\dfrac{mN(\alpha+\varepsilon)}{n}+\dfrac{\ln{a_k}}{n}$,\ding{172} \\

令$\displaystyle C=\max_{1\leq j\leq N}|\ln a_j|$,因为$\displaystyle 1\leq k\leq N$,所以$\displaystyle\dfrac{|\ln a_k|}{n}\leq\dfrac{C}{n}\to0,$ \\

又$\displaystyle\dfrac{mN}{n}=1-\dfrac{k}{n}$,故$\displaystyle\lim_{n\to\infty}\dfrac{mN}{n}=1,\quad\lim_{n\to\infty}\dfrac{\ln a_k}{n}=0,$ \\

对\ding{172}式两边同时关于$n$取上极限:$\displaystyle\varlimsup_{n\to\infty}\dfrac{\ln{a_n}}{n}\leq\alpha+\varepsilon,$ \\

由$\varepsilon$的任意性,$\displaystyle\varlimsup_{n\to\infty}\dfrac{\ln{a_n}}{n}\leq\alpha,\text{又}\displaystyle\varlimsup_{n\to\infty}\dfrac{\ln{a_n}}{n}\geq\displaystyle\varliminf_{n\to\infty}\dfrac{\ln{a_n}}{n},$ \\

$\therefore\displaystyle\varlimsup_{n\to\infty}\dfrac{\ln{a_n}}{n}=\displaystyle\varliminf_{n\to\infty}\dfrac{\ln{a_n}}{n}=\alpha,$ \\

$\therefore\displaystyle\lim_{n\to\infty}\dfrac{\ln{a_n}}{n}\text{有意义,且}\displaystyle\lim_{n\to\infty}\dfrac{\ln{a_n}}{n}=\displaystyle\inf_{n\geq 1}\{\dfrac{\ln{a_n}}{n}\}.$ \\

\textbf{注}:由上可得到正数列$\{a_n\}$满足条件$a_{m+n}\leq a_na_m$时,$\{\sqrt[n]{a_n} \}$一定收敛于$\displaystyle\inf_{n\geq 1}\{\sqrt[n]{a_n} \}$



\BookEnd

\ifdefined\BookMaster
\def\BookEnd{}
\else
\documentclass[a4paper,12pt,fleqn,openany,twoside,fontset=windows]{ctexbook}

\usepackage[fleqn]{amsmath}
\usepackage{amssymb,amsthm,mathrsfs}
\usepackage{geometry}
\usepackage[dvipsnames]{xcolor}
\usepackage{graphicx}
\usepackage{tikz}
\usetikzlibrary{calc}
\usepackage[most]{tcolorbox}
\usepackage{fancyhdr}
\usepackage{titlesec}
\usepackage{tocloft}
\usepackage{enumitem}
\usepackage{microtype}
\usepackage{pifont}
\usepackage{hyperref}
\usepackage{romannum}
\usepackage{mathtools}

\definecolor{bookblue}{HTML}{264653}
\definecolor{bookteal}{HTML}{4F7C82}
\definecolor{booklight}{HTML}{EDF4F3}
\definecolor{bookgray}{HTML}{5E686B}

\geometry{
    inner=2.7cm,
    outer=2.5cm,
    top=2.7cm,
    bottom=2.7cm,
    headheight=18pt,
    headsep=18pt,
    footskip=25pt
}
\setlength{\mathindent}{0pt}
\setlength{\parindent}{5pt}
\setlength{\parskip}{3pt plus 1pt minus 1pt}
\linespread{1.25}
\allowdisplaybreaks[3]
\AtBeginDocument{%
    \setlength{\parindent}{5pt}
    \setlength{\abovedisplayskip}{8pt plus 2pt minus 2pt}
    \setlength{\belowdisplayskip}{8pt plus 2pt minus 2pt}
    \setlength{\abovedisplayshortskip}{5pt plus 1pt minus 1pt}
    \setlength{\belowdisplayshortskip}{5pt plus 1pt minus 1pt}
}

\hypersetup{
    unicode=true,
    colorlinks=true,
    linkcolor=bookblue,
    citecolor=bookblue,
    urlcolor=bookteal,
    bookmarksopen=true,
    pdfauthor={Chen},
    pdftitle={数学分析习题整理}
}

\renewcommand{\chaptermark}[1]{%
    \edef\@tmp{第\thechapter 章\quad #1}%
    \expandafter\markboth\expandafter{\@tmp}{}%
}
\fancypagestyle{main}{%
    \fancyhf{}
    \fancyhead[LE]{\small\color{bookgray}\nouppercase{\leftmark}}
    \fancyhead[RO]{\small\color{bookgray}数学分析习题整理}
    \fancyfoot[C]{\color{bookblue}\thepage}
    \fancyfoot[R]{\small\bfseries Chen\;\; github.com/mathlover-1/math-analysis}
    \renewcommand{\headrulewidth}{0.4pt}
    \renewcommand{\footrulewidth}{0pt}
}
\fancypagestyle{plain}{%
    \fancyhf{}
    \fancyfoot[C]{\color{bookblue}\thepage}
    \fancyfoot[R]{\small\bfseries Chen\;\; github.com/mathlover-1/math-analysis}
    \renewcommand{\headrulewidth}{0pt}
    \renewcommand{\footrulewidth}{0pt}
}
\pagestyle{main}

\titleformat{\chapter}[display]
  {\normalfont\sffamily\bfseries\color{bookblue}}
  {\raggedleft\fontsize{46}{50}\selectfont\thechapter}
  {-18pt}
  {\titlerule[1.2pt]\vspace{10pt}\filright\Huge}
\titlespacing*{\chapter}{0pt}{-15pt}{36pt}

\renewcommand{\contentsname}{目\quad 录}
\renewcommand{\cfttoctitlefont}{\hfill\sffamily\bfseries\Huge\color{bookblue}}
\renewcommand{\cftaftertoctitle}{\hfill}
\setlength{\cftbeforetoctitleskip}{5pt}
\setlength{\cftaftertoctitleskip}{28pt}
\setlength{\cftbeforechapskip}{12pt}
\renewcommand{\cftchapfont}{\sffamily\bfseries\large\color{bookblue}}
\renewcommand{\cftchappagefont}{\sffamily\bfseries\large\color{bookblue}}
\renewcommand{\cftchapleader}{\color{bookteal}\cftdotfill{2.5}}

\newtcolorbox{ti}{
    enhanced,
    breakable,
    colback=booklight!55!white,
    colframe=bookteal!35!white,
    boxrule=0.45pt,
    boxsep=0pt,
    left=10pt,
    right=10pt,
    top=9pt,
    bottom=9pt,
    before skip=14pt,
    after skip=14pt,
    arc=2pt,
    outer arc=2pt,
    before upper={\parindent=0pt}
}

\newenvironment{booknotebody}{%
    \begin{center}
    \begin{minipage}{0.84\textwidth}
    \songti\zihao{-4}\linespread{1.45}\selectfont
    \setlength{\parindent}{2\ccwd}
    \setlength{\parskip}{0.75em}
}{%
    \end{minipage}
    \end{center}
}

\newenvironment{booknotelongbody}{%
    \begingroup
    \songti\zihao{-4}\linespread{1.45}\selectfont
    \setlength{\parindent}{2\ccwd}
    \setlength{\parskip}{0.75em}
    \begin{list}{}{%
        \setlength{\leftmargin}{0.08\textwidth}
        \setlength{\rightmargin}{0.08\textwidth}
        \setlength{\topsep}{0pt}
        \setlength{\partopsep}{0pt}
        \setlength{\parsep}{0pt}
        \setlength{\itemsep}{0pt}
    }
    \item\relax
}{%
    \end{list}
    \endgroup
}

\fancypagestyle{plainnowm}{%
    \fancyhf{}
    \fancyfoot[C]{\color{bookblue}\thepage}
    \renewcommand{\headrulewidth}{0pt}
    \renewcommand{\footrulewidth}{0pt}
}

\newcommand{\booknotetitle}[2][plain]{%
    \clearpage
    \phantomsection
    \addcontentsline{toc}{chapter}{#2}
    \markboth{#2}{}
    \thispagestyle{#1}
    \vspace*{1.2cm}
    \begin{center}
        {\sffamily\heiti\bfseries\fontsize{26}{34}\selectfont\color{bookblue}#2\par}
        \vspace{0.8em}
        {\color{bookteal}\rule{4.5cm}{0.8pt}\par}
    \end{center}
    \vspace{0.6cm}
}

\renewcommand{\proofname}{证明}
\renewcommand{\qedsymbol}{\color{bookteal}\ensuremath{\blacksquare}}

\title{数学分析习题整理}
\author{Chen}
\date{}

\makeatletter
% ==================== 旧版封面/封底设计备份 ====================
% 2026-08 重设计,旧代码用 \iffalse...\fi 封存,新版见下方。
\iffalse
\renewcommand{\maketitle}{%
    \begin{titlepage}
        \thispagestyle{empty}
        \setcounter{page}{0}
        \noindent\hspace*{-\parindent}\begin{tikzpicture}[remember picture,overlay]
            \fill[bookblue] (current page.north west) rectangle ([yshift=-4.2cm]current page.north east);
            \fill[booklight] ([yshift=-4.2cm]current page.north west) rectangle (current page.south east);
            \draw[bookteal!55,line width=0.7pt]
                ([xshift=2.4cm,yshift=4.3cm]current page.south west)
                .. controls +(2.5cm,1.8cm) and +(-2.8cm,-1.3cm) ..
                ([xshift=-2.4cm,yshift=7.2cm]current page.south east);
            \draw[bookteal!35,line width=0.7pt]
                ([xshift=2.5cm,yshift=6.2cm]current page.south west)
                .. controls +(3.0cm,-1.2cm) and +(-3.0cm,1.6cm) ..
                ([xshift=-2.3cm,yshift=3.8cm]current page.south east);
        \end{tikzpicture}
        \vspace*{5.6cm}
        \noindent\color{bookblue}\rule{\textwidth}{1.4pt}\par
        \vspace{8pt}
        {\centering\sffamily\bfseries\fontsize{32}{42}\selectfont\color{bookblue}\@title\par}
        \vspace{8pt}
        \noindent\color{bookteal}\rule{\textwidth}{0.6pt}\par
        \vspace{2.6cm}
        {\centering\Large\color{bookgray}\@author\par}
        \vfill
        {\centering\color{bookteal}\large
            $\displaystyle \lim_{n\to\infty} a_n$\qquad
            $\displaystyle \int_a^b f(x)\,\mathrm{d}x$\qquad
            $\displaystyle \nabla f(x,y)$\par
        }
        \vspace*{1.5cm}
    \end{titlepage}
}

\newcommand{\makebackcover}{%
    \begingroup
    \let\cleardoublepage\clearpage
    \begin{titlepage}
        \thispagestyle{empty}
        \noindent\hspace*{-\parindent}\begin{tikzpicture}[remember picture,overlay]
            \fill[bookblue] (current page.north west) rectangle (current page.south east);
            \fill[booklight,rounded corners=7pt]
                ([xshift=2.2cm,yshift=-3.0cm]current page.north west)
                rectangle
                ([xshift=-2.2cm,yshift=3.0cm]current page.south east);
            \fill[bookteal!80]
                (current page.north west) --
                ([xshift=5.3cm]current page.north west) --
                ([xshift=2.1cm,yshift=-4.4cm]current page.north west) --
                ([yshift=-6.2cm]current page.north west) -- cycle;
            \fill[bookteal!55]
                (current page.south east) --
                ([xshift=-5.8cm]current page.south east) --
                ([xshift=-2.4cm,yshift=4.2cm]current page.south east) --
                ([yshift=6.0cm]current page.south east) -- cycle;
            \foreach \radius in {0.75,1.15,1.55,1.95} {
                \draw[bookteal!55,line width=0.55pt]
                    ([xshift=-2.9cm,yshift=-4.4cm]current page.north east)
                    circle[radius=\radius cm];
            }
            \foreach \radius in {0.65,1.05,1.45} {
                \draw[booklight!65,line width=0.55pt]
                    ([xshift=2.4cm,yshift=3.6cm]current page.south west)
                    circle[radius=\radius cm];
            }
            \fill[bookteal!16,rounded corners=3pt]
                ([xshift=3.5cm,yshift=-7.1cm]current page.north west)
                rectangle
                ([xshift=-5.0cm,yshift=7.2cm]current page.south east);
            \fill[bookteal!30,rounded corners=3pt]
                ([xshift=4.4cm,yshift=-8.0cm]current page.north west)
                rectangle
                ([xshift=-4.1cm,yshift=8.1cm]current page.south east);
            \begin{scope}[shift={([xshift=-0.75cm]current page.center)}]
                \clip[rounded corners=3pt] (-5.25,-4.75) rectangle (5.25,4.75);
                \foreach \v in {-3.6,-3.0,...,3.6} {
                    \draw[bookteal!48,line width=0.45pt,opacity=0.72]
                        plot[domain=-3.6:3.6,samples=45,smooth,variable=\u]
                        ({\u+0.55*\v},{0.28*\u-0.35*\v+0.11*(\u*\u-\v*\v)});
                }
                \foreach \u in {-3.6,-3.0,...,3.6} {
                    \draw[bookblue!42,line width=0.42pt,opacity=0.62]
                        plot[domain=-3.6:3.6,samples=45,smooth,variable=\v]
                        ({\u+0.55*\v},{0.28*\u-0.35*\v+0.11*(\u*\u-\v*\v)});
                }
                \draw[bookteal!70,line width=0.75pt,opacity=0.82]
                    plot[domain=-3.6:3.6,samples=55,smooth,variable=\u]
                    ({\u},{0.28*\u+0.11*\u*\u});
                \draw[bookblue!60,line width=0.7pt,opacity=0.75]
                    plot[domain=-3.6:3.6,samples=55,smooth,variable=\v]
                    ({0.55*\v},{-0.35*\v-0.11*\v*\v});
            \end{scope}
            \draw[bookteal!45,line width=0.65pt]
                ([xshift=3.0cm,yshift=-8.2cm]current page.north west)
                .. controls +(2.2cm,2.0cm) and +(-2.0cm,-2.1cm) ..
                ([xshift=-3.0cm,yshift=9.0cm]current page.south east);
            \draw[bookteal!25,line width=0.55pt]
                ([xshift=3.1cm,yshift=8.0cm]current page.south west)
                .. controls +(2.7cm,-1.7cm) and +(-2.6cm,1.8cm) ..
                ([xshift=-3.1cm,yshift=-8.7cm]current page.north east);
            \node[anchor=west,text=bookteal!72,scale=0.78] at
                ([xshift=2.8cm,yshift=-4.0cm]current page.north west) {%
                $\displaystyle e^{\mathrm{i}\pi}+1=0$};
            \node[anchor=east,text=bookblue!62,scale=0.67] at
                ([xshift=-2.7cm,yshift=4.0cm]current page.south east) {%
                $\displaystyle
                f(x)=\sum_{n=0}^{\infty}\frac{f^{(n)}(a)}{n!}(x-a)^n$};
        \end{tikzpicture}
        \mbox{}
    \end{titlepage}
    \endgroup
}
\fi

% ==================== 新版封面(2026-08 重设计) ====================
% 设计要点:顶部深绿纵向渐变+波浪下缘+两层半透明过渡波(解决深浅绿生硬过渡);
% 深色区螺旋线/同心圆母题;标题+英文副题+作者居中;
% 中部 Fourier 部分和逼近方波曲线;底部浅色波纹(无文字)。
\renewcommand{\maketitle}{%
    \begin{titlepage}
        \thispagestyle{empty}
        \setcounter{page}{0}
        \noindent\hspace*{-\parindent}\begin{tikzpicture}[remember picture,overlay]
            % ---------- 底色 ----------
            \fill[booklight] (current page.north west) rectangle (current page.south east);

            % ---------- 顶部深色区:纵向渐变 + 波浪下缘 ----------
            \shade[top color=bookblue, bottom color=bookteal!55!bookblue]
                (current page.north west) -- (current page.north east) --
                ([yshift=-8.6cm]current page.north east)
                .. controls ([xshift=-5.2cm,yshift=-10.6cm]current page.north east)
                         and ([xshift=6.0cm,yshift=-8.0cm]current page.north west) ..
                ([yshift=-10.0cm]current page.north west) -- cycle;

            % ---------- 过渡层 1:半透明波浪 ----------
            \fill[bookteal!60!bookblue,opacity=0.42]
                ([yshift=-8.6cm]current page.north east)
                .. controls ([xshift=-5.2cm,yshift=-10.6cm]current page.north east)
                         and ([xshift=6.0cm,yshift=-8.0cm]current page.north west) ..
                ([yshift=-10.0cm]current page.north west) --
                ([yshift=-12.6cm]current page.north west)
                .. controls ([xshift=6.5cm,yshift=-11.0cm]current page.north west)
                         and ([xshift=-5.5cm,yshift=-13.2cm]current page.north east) ..
                ([yshift=-11.4cm]current page.north east) -- cycle;

            % ---------- 过渡层 2:更浅的波浪 ----------
            \fill[bookteal!35,opacity=0.35]
                ([yshift=-10.6cm]current page.north east)
                .. controls ([xshift=-6.0cm,yshift=-12.4cm]current page.north east)
                         and ([xshift=5.0cm,yshift=-9.8cm]current page.north west) ..
                ([yshift=-11.8cm]current page.north west) --
                ([yshift=-13.8cm]current page.north west)
                .. controls ([xshift=5.5cm,yshift=-12.6cm]current page.north west)
                         and ([xshift=-6.5cm,yshift=-14.6cm]current page.north east) ..
                ([yshift=-12.8cm]current page.north east) -- cycle;

            % ---------- 深色区装饰:右上螺旋线 ----------
            \begin{scope}[shift={([xshift=-3.4cm,yshift=-4.3cm]current page.north east)}]
                \draw[booklight!45!bookteal,opacity=0.55,line width=0.6pt]
                    plot[domain=0:1440,samples=220,smooth,variable=\t]
                    ({0.0028*\t*cos(\t)},{0.0028*\t*sin(\t)});
            \end{scope}

            % ---------- 深色区装饰:左下同心圆 ----------
            \foreach \r in {0.7,1.25,1.8,2.35,2.9}{
                \draw[booklight!35!bookteal,opacity=0.5,line width=0.5pt]
                    ([xshift=0.4cm,yshift=-7.9cm]current page.north west) circle[radius=\r cm];
            }

            % ---------- 主标题 ----------
            \node[anchor=center,text=booklight] at
                ([yshift=-4.5cm]current page.north)
                {\sffamily\bfseries\fontsize{40}{46}\selectfont \@title};

            % ---------- 分隔短线 ----------
            \draw[booklight!65,line width=1.2pt]
                ([xshift=-2.6cm,yshift=-5.75cm]current page.north) --
                ([xshift=2.6cm,yshift=-5.75cm]current page.north);

            % ---------- 英文副题 ----------
            \node[anchor=center,text=booklight!70!bookteal] at
                ([yshift=-6.55cm]current page.north)
                {\itshape\large Mathematical Analysis \textperiodcentered\
                 A Collection of Exercises};

            % ---------- 作者(紧随标题下方) ----------
            \node[anchor=center,text=booklight!88] at
                ([yshift=-8.2cm]current page.north)
                {\sffamily\large \@author\quad 整理};

            % ---------- 中部:Fourier 部分和逼近方波 ----------
            \begin{scope}[shift={([xshift=-5.2cm,yshift=-18.1cm]current page.north)}]
                % 坐标轴
                \draw[bookgray!55,line width=0.5pt,->] (-0.5,0) -- (10.9,0);
                \draw[bookgray!55,line width=0.5pt,->] (0,-1.75) -- (0,1.95);
                % 目标方波(虚线)
                \draw[bookgray!75,line width=0.6pt,dashed]
                    (0,0.92) -- (5,0.92) (5,-0.92) -- (10,-0.92);
                % N=1
                \draw[bookteal!45,line width=0.7pt]
                    plot[domain=0:360,samples=90,smooth,variable=\x]
                    ({\x/36},{1.18*sin(\x)});
                % N=3
                \draw[bookteal!80,line width=0.8pt]
                    plot[domain=0:360,samples=120,smooth,variable=\x]
                    ({\x/36},{1.18*(sin(\x)+sin(3*\x)/3)});
                % N=7
                \draw[bookblue,line width=0.95pt]
                    plot[domain=0:360,samples=150,smooth,variable=\x]
                    ({\x/36},{1.18*(sin(\x)+sin(3*\x)/3+sin(5*\x)/5+sin(7*\x)/7)});
                % 刻度标注
                \node[anchor=north east,text=bookgray!80,scale=0.62] at (-0.12,0) {$0$};
                \node[anchor=north,text=bookgray!80,scale=0.62] at (5,0) {$\pi$};
                \node[anchor=north,text=bookgray!80,scale=0.62] at (10,0) {$2\pi$};
            \end{scope}
            % 图注
            \node[anchor=center,text=bookgray,scale=0.8] at
                ([yshift=-20.5cm]current page.north)
                {$S_N(x)=\displaystyle\sum_{k=1}^{N}\frac{\sin\bigl((2k-1)x\bigr)}{2k-1}$
                 \ 逐步逼近方波};

            % ---------- 底部浅色波纹(无文字,与顶部波浪呼应) ----------
            \fill[bookteal!22,opacity=0.5]
                ([yshift=2.4cm]current page.south west)
                .. controls ([xshift=6.0cm,yshift=3.9cm]current page.south west)
                         and ([xshift=-6.5cm,yshift=1.0cm]current page.south east) ..
                ([yshift=2.9cm]current page.south east) --
                (current page.south east) -- (current page.south west) -- cycle;
            \fill[bookteal!40!bookblue,opacity=0.28]
                ([yshift=1.3cm]current page.south west)
                .. controls ([xshift=5.5cm,yshift=2.7cm]current page.south west)
                         and ([xshift=-5.5cm,yshift=0.3cm]current page.south east) ..
                ([yshift=1.9cm]current page.south east) --
                (current page.south east) -- (current page.south west) -- cycle;
        \end{tikzpicture}
        \mbox{}
    \end{titlepage}
}

% ==================== 新版封底(纯背景,2026-08 重设计) ====================
% 设计要点:对角渐变底+角部半透明叠层;中下部外摆线(玫瑰形)为视觉主体;
% 底部两层半透明白色波浪;无任何文字。
\newcommand{\makebackcover}{%
    \begingroup
    \let\cleardoublepage\clearpage
    \begin{titlepage}
        \thispagestyle{empty}
        \noindent\hspace*{-\parindent}\begin{tikzpicture}[remember picture,overlay]
            % ---------- 底色:对角渐变 ----------
            \shade[top color=bookblue, bottom color=bookteal!48!bookblue,
                   shading angle=-30]
                (current page.north west) rectangle (current page.south east);

            % ---------- 角部半透明叠层 ----------
            \fill[white,opacity=0.05]
                ([xshift=2.0cm,yshift=1.0cm]current page.north east) circle[radius=7.2cm];
            \fill[white,opacity=0.06]
                ([xshift=-0.5cm,yshift=-1.5cm]current page.north east) circle[radius=3.6cm];

            % ---------- 左下同心圆 ----------
            \foreach \r in {1.0,1.75,2.5,3.25,4.0,4.75}{
                \draw[booklight!40,opacity=0.55,line width=0.55pt]
                    ([xshift=-0.6cm,yshift=-0.4cm]current page.south west) circle[radius=\r cm];
            }
            % ---------- 右下螺旋线 ----------
            \begin{scope}[shift={([xshift=-2.9cm,yshift=3.2cm]current page.south east)}]
                \draw[booklight!45!bookteal,opacity=0.5,line width=0.55pt]
                    plot[domain=0:1260,samples=200,smooth,variable=\t]
                    ({0.0036*\t*cos(\t)},{0.0036*\t*sin(\t)});
            \end{scope}
            % ---------- 左上斜线组 ----------
            \foreach \i in {0,1,2}{
                \draw[bookteal!60!booklight,opacity=0.5,line width=0.6pt]
                    ([xshift={1.6cm+\i*0.45cm}]current page.north west) --
                    ([yshift={-1.6cm-\i*0.45cm}]current page.north west);
            }

            % ---------- 中下部视觉主体:外摆线(玫瑰形) ----------
            \begin{scope}[shift={([yshift=-16.6cm]current page.north)}]
                % 外圈
                \draw[booklight!35,opacity=0.6,line width=0.6pt]
                    (0,0) circle[radius=4.9cm];
                \draw[booklight!22,opacity=0.5,line width=0.5pt]
                    (0,0) circle[radius=5.25cm];
                % 主玫瑰线
                \draw[booklight!55!bookteal,opacity=0.8,line width=0.85pt]
                    plot[domain=0:1080,samples=400,smooth,variable=\t]
                    ({0.335*(8*cos(\t)-5*cos(2.6667*\t))},
                     {0.335*(8*sin(\t)-5*sin(2.6667*\t))});
                % 内层玫瑰线
                \draw[booklight!38!bookteal,opacity=0.6,line width=0.6pt]
                    plot[domain=0:1080,samples=400,smooth,variable=\t]
                    ({0.21*(8*cos(\t)-5*cos(2.6667*\t))},
                     {0.21*(8*sin(\t)-5*sin(2.6667*\t))});
            \end{scope}

            % ---------- 底部半透明波浪 ----------
            \fill[white,opacity=0.06]
                ([yshift=4.6cm]current page.south west)
                .. controls ([xshift=6.5cm,yshift=6.4cm]current page.south west)
                         and ([xshift=-6.0cm,yshift=2.6cm]current page.south east) ..
                ([yshift=5.2cm]current page.south east) --
                (current page.south east) -- (current page.south west) -- cycle;
            \fill[booklight,opacity=0.08]
                ([yshift=2.6cm]current page.south west)
                .. controls ([xshift=5.5cm,yshift=4.3cm]current page.south west)
                         and ([xshift=-6.5cm,yshift=1.0cm]current page.south east) ..
                ([yshift=3.3cm]current page.south east) --
                (current page.south east) -- (current page.south west) -- cycle;
        \end{tikzpicture}
        \mbox{}
    \end{titlepage}
    \endgroup
}
\makeatother

% 页脚右下方常态化水印(署名 + 仓库地址,加粗黑色),已在 main/plain 分页样式中定义,所有页面均生效.

\begin{document}
\mainmatter
\pagestyle{main}
\setcounter{chapter}{2}
\def\BookEnd{\end{document}}
\fi

\chapter{函数}

\begin{ti}
【题1】用覆盖定理证明零点存在性定理
\end{ti}

\textbf{证明}:给定连续函数$f,f(a)<0,f(b)>0,$ \\

假设函数$f\text{在}[a,b]$上无零点,则对$\forall\,x_0\in[a,b],$由连续函数保号性, \\

$\exists\,\delta_x>0,\;\text{s.t.}\,\forall\,x\in O_{\delta_x}(x_0)\cap[a,b],\text{都有}f(x)f(x_0)>0,$ \\

又$\displaystyle\bigcup_{x\in[a,b]}O_{\delta_x}(x)\text{构成}[a,b]$的一个开覆盖,那么,由$Lebesgue$定理, \\

$\exists\,\zeta>0,\;\text{s.t.}\,x',x''\in[a,b]\text{且}\left\lvert x'-x''\right\rvert<\zeta\text{时},x',x''\text{可被包在同一个}O_{\delta}(x)$里, \\

在$[a,b]\text{中插入一系列点,连同端点一起,记为:}a=x_0<x_1<x_2<...<x_n=b,$ \\

$\text{且满足}\left\lvert x_i-x_{i-1}\right\rvert<\zeta,\,i=1,2,...,n,$ \\

由插入点的特点可知,$f(x_i)f(x_{i-1})>0\Longrightarrow f(a)f(b)>0,$矛盾.

\begin{ti}
【题2】设$f\in C[a,b),\{x_n\},\{y_n\}\subset[a,b),\,\,\text{满足}:$\\

$\displaystyle\lim_{n\to\infty}x_n=\lim_{n\to\infty}y_n=b,\displaystyle\lim_{n\to\infty}f(x_n)=A<B=\lim_{n\to\infty}f(y_n)$ \\

证明:$\forall\,\eta\in(A,B),\;\exists\,\{z_n\}\subset[a,b),\;\text{s.t.}\,\displaystyle\lim_{n\to\infty}z_n=b,\lim_{n\to\infty}f(z_n)=\eta$
\end{ti}

\textbf{证明}:取$\varepsilon=\min\{\eta-A,B-\eta\},\text{则}\exists\,N,$ \\

$n>N\text{时},\,f(x_n)<A+\varepsilon=\eta,f(y_n)>B-\varepsilon=\eta,$ \\

又$x_n,y_n\in[a,b),f\in C[a,b)$,由零点存在性定理, \\

$\exists\,z_n\text{介于}x_n,y_n\text{之间},\text{且}f(z_n)=\eta,$ \\

又由夹逼定理,可得$\displaystyle\lim_{n\to\infty}z_n=b.$

\begin{ti}
【题3】设$f$在$(a,b)$上一致连续,证明:存在两个有限的单侧极限$f(a^+),f(b^-)$
\end{ti}

\textbf{证明}:$\forall\,\varepsilon>0,\;\exists\,\delta>0,\;\text{s.t.}\,\left\lvert x'-x''\right\rvert<\delta\text{且}x',x''\in(a,b)\text{时},\left\lvert f(x')-f(x'')\right\rvert<\varepsilon,$ \\

$\therefore\,x_1,x_2\in(a,a+\delta)\text{时},\left\lvert f(x_1)-f(x_2)\right\rvert<\varepsilon,$ \\

由$Cauchy$收敛准则,$\displaystyle\lim_{x\to a^+}f(x)$存在且有限.同理$f(b^-)$有限,得证.

\begin{ti}
【题4】证明:$f(x)=x\sin{\dfrac{1}{x}}$在$(0,+\infty)$上一致连续
\end{ti}

\textbf{证明}:$\left\lvert f(x)-f(y)\right\rvert=\left\lvert x\sin{\dfrac{1}{x}}-y\sin{\dfrac{1}{y}}\right\rvert=\left\lvert x\sin{\dfrac{1}{x}}-y\sin{\dfrac{1}{x}}+y\sin{\dfrac{1}{x}}-y\sin{\dfrac{1}{y}}\right\rvert$ \\

$<\left\lvert \sin{\dfrac{1}{x}}\right\rvert\left\lvert x-y\right\rvert+y\left\lvert \sin{\dfrac{1}{x}-\dfrac{1}{y}}\right\rvert<\left\lvert x-y\right\rvert+y\left\lvert \sin{\dfrac{1}{x}-\dfrac{1}{y}}\right\rvert,$ \\

利用$\sin{x}\text{是}Lipschitz\text{函数},\,\left\lvert \sin{\dfrac{1}{x}-\dfrac{1}{y}}\right\rvert<\left\lvert \dfrac{1}{x}-\dfrac{1}{y}\right\rvert,$ \\

$\therefore\left\lvert f(x)-f(y)\right\rvert<\left\lvert x-y\right\rvert+y\left\lvert \dfrac{1}{x}-\dfrac{1}{y}\right\rvert=\left\lvert x-y\right\rvert+\dfrac{1}{x}\left\lvert x-y\right\rvert,$ \\

同理$\left\lvert f(x)-f(y)\right\rvert<\left\lvert x-y\right\rvert+\dfrac{1}{y}\left\lvert x-y\right\rvert,$ \\

又$\left\lvert f(x)-f(y)\right\rvert<\left\lvert x\sin{\dfrac{1}{x}}\right\rvert+\left\lvert y\sin{\dfrac{1}{y}}\right\rvert<x+y,$ \\

对$\forall\,\varepsilon>0,\;\text{取}\delta=\varepsilon(1+\dfrac{2}{\varepsilon})^{-1},$我们断言:$\left\lvert x-y\right\rvert<\delta\text{时},\left\lvert f(x)-f(y)\right\rvert<\varepsilon,$ \\

分情况讨论, \\

\romannum{1} 若$x<\dfrac{\varepsilon}{2},\,y<\dfrac{\varepsilon}{2},\text{则}\left\lvert f(x)-f(y)\right\rvert<x+y<\varepsilon,$ \\

\romannum{2} 若$x\geq\dfrac{\varepsilon}{2}(y\geq\dfrac{\varepsilon}{2}\text{时类似可这样处理}),\text{则}\left\lvert f(x)-f(y)\right\rvert<\left\lvert x-y\right\rvert+\dfrac{1}{x}\left\lvert x-y\right\rvert$ \\

$\leq\left\lvert x-y\right\rvert+\dfrac{2}{\varepsilon}\left\lvert x-y\right\rvert<(1+\dfrac{2}{\varepsilon})\delta=\varepsilon,$ \\

即$\forall\,\varepsilon>0,\;\exists\,\delta=\varepsilon(1+\dfrac{2}{\varepsilon})^{-1},\left\lvert x-y\right\rvert<\delta\text{时},\left\lvert f(x)-f(y)\right\rvert<\varepsilon,$ \\

即$f(x)=x\sin{\dfrac{1}{x}}\text{在}(0,+\infty)$上一致连续. \\

\textbf{注}:本题中$\left\lvert f(x)-f(y)\right\rvert=\left\lvert x\sin{\dfrac{1}{x}}-y\sin{\dfrac{1}{x}}+y\sin{\dfrac{1}{x}}-y\sin{\dfrac{1}{y}}\right\rvert$的处理方法在解决这种差值估计当中十分常见,
但是本题中这种方法在0附近失效了,不能直接解决问题,所以要分段去用不同的放缩,在0附近利用$f$趋近于0来放缩.

\begin{ti}
【题5】设$g\text{是区间}I\text{上的单调函数}$,证明:值域$g(I)$为区间的充要条件为$g\in C^0(I)$
\end{ti}

\textbf{证明}:充分性易得,下证必要性,先设$\displaystyle g$是单调递增的, \\

若$\displaystyle g$在内点$\displaystyle x_0\in I$处不连续,则$\displaystyle g(x_0^-)<g(x_0^+)$,取$\displaystyle y\in(g(x_0^-),g(x_0^+))$且$\displaystyle y\neq g(x_0)$, \\

由单调性可知$\displaystyle y\notin g(I)$,与$\displaystyle g(I)$是区间矛盾, \\

若$\displaystyle I$含左端点$\displaystyle a$,且$\displaystyle g$在$\displaystyle a$处相对右不连续,则$\displaystyle g(a)<g(a^+)$, \\

任取$\displaystyle y\in(g(a),g(a^+))$,对$\displaystyle x=a$有$\displaystyle g(x)<y$,对$\displaystyle x>a$有$\displaystyle g(x)\geq g(a^+)>y$, \\

故$\displaystyle y\notin g(I)$,矛盾, \\

若$\displaystyle I$含右端点$\displaystyle b$,且$\displaystyle g$在$\displaystyle b$处相对左不连续,则$\displaystyle g(b^-)<g(b)$,\\

任取$\displaystyle y\in(g(b^-),g(b))$,对$\displaystyle x<b$有$\displaystyle g(x)\leq g(b^-)<y$,而$\displaystyle g(b)>y$,故$\displaystyle y\notin g(I)$,矛盾, \\

所以单调递增时$\displaystyle g\in C^0(I)$.若$\displaystyle g$单调递减,对$\displaystyle-g$应用上述结论即可,故$\displaystyle g\in C^0(I).$

\begin{ti}
【题6】设$f$是$(-\infty,+\infty)$上的单调函数,在每一点定义$g(x)=f(x^+)$,证明:$g\text{是}(-\infty,+\infty)$上处处右连续的函数
\end{ti}

\textbf{证明}:即证$\forall\,x_0\in\mathbb{R},\;\displaystyle\lim_{x\to x_0^+}f(x^+)=f(x_0^+),$ \\

不妨设$f$是单调递增函数,则$x>x_0\text{时},f(x_0^+)\leq f(x)\leq f(x^+),$ \\

又$f(x_0^+)=\displaystyle\inf_{x>x_0}f(x),$由确界的刻画可知, \\

$\forall\,\varepsilon>0,\;\exists\,u_0>x_0,\;\text{s.t.}\,f(x_0^+)+\varepsilon>f(u_0),$ \\

令$\delta=u_0-x_0>0,\,x\in(x_0,x_0+\delta)\text{时},f(x^+)\leq f(u_0)<f(x_0^+)+\varepsilon,$ \\

$\therefore\,\forall\varepsilon>0,\;\exists\,\delta>0,x\in(x_0,x_0+\delta)\text{时},\left\lvert f(x^+)-f(x_0^+)\right\rvert<\varepsilon,$ \\

$\therefore\displaystyle\lim_{x\to x_0^+}f(x^+)=f(x_0^+)$对$\forall\,x_0\in\mathbb{R}$成立,即$g(x)$在$(-\infty,+\infty)$上处处右连续.

\begin{ti}
【题7】设函数$f$在区间$I$内只有可去间断点,定义$g(x)=\displaystyle\lim_{t\to x}f(t),x\in I,$证明:$g\in C(I)$
\end{ti}

\textbf{证明}:因为$\displaystyle f$只有可去间断点,所以$\displaystyle g$在$\displaystyle I$中每一点都有定义, \\

给定$\displaystyle\varepsilon>0$,由极限定义,存在$\displaystyle\delta>0$,当$\displaystyle0<|t-x_0|<\delta$时,$\displaystyle\left\lvert f(t)-g(x_0)\right\rvert<\dfrac{\varepsilon}{2},$ \\

若$\displaystyle|x-x_0|<\dfrac{\delta}{2}$,令$\displaystyle y\to x$且$\displaystyle0<|y-x|<\dfrac{\delta}{2}-|x-x_0|$,则$\displaystyle|y-x_0|<\delta$,从而 \\

$\displaystyle g(x_0)-\dfrac{\varepsilon}{2}\leq\lim_{y\to x}f(y)=g(x)\leq g(x_0)+\dfrac{\varepsilon}{2},$ 故$\displaystyle\left\lvert g(x)-g(x_0)\right\rvert\leq\dfrac{\varepsilon}{2}<\varepsilon$,\\

因此$\displaystyle\lim_{x\to x_0}g(x)=g(x_0)$对任意$\displaystyle x_0\in I$成立,即$\displaystyle g\in C(I).$

\begin{ti}
【题8】证明:若$f$在$[0,+\infty)$上连续且有界,则对每个给定的$\lambda$,存在一个数列$\{x_n\}$满足要求: \\
$\displaystyle(1)\lim_{n\to\infty}x_n=+\infty;\quad(2)\lim_{n\to\infty}[f(x_n+\lambda)-f(x_n)]=0$
\end{ti}

\textbf{证明}:若$\lambda=0,\text{取}x_n=n$即可,下设$\lambda\neq 0$, \\

考虑函数$\displaystyle g(x)=f(x+\lambda)-f(x)$,则$\displaystyle g$在$\displaystyle[\max\{0,-\lambda\},+\infty)$上连续且有界, \\

下面证明:$\displaystyle\forall\,\varepsilon>0,\forall\,M>0,\;\exists\,x\geq M,\;\text{s.t.}\,\left\lvert g(x)\right\rvert<\varepsilon,$ \\

用反证法,假设$\displaystyle\exists\,\varepsilon_0>0,M_0>0,\;\text{s.t.}\,\forall\,x\geq M_0,\left\lvert g(x)\right\rvert\geq\varepsilon_0$, \\

令$\displaystyle L=\max\{0,-\lambda,M_0\}$,因为$\displaystyle g$在连通区间$\displaystyle[L,+\infty)$上连续且不取零值, \\

所以由介值定理,$\displaystyle g$在该区间上恒取同一符号, \\

即$\displaystyle g(x)\geq\varepsilon_0$恒成立或$\displaystyle g(x)\leq-\varepsilon_0$恒成立, \\

下设前一种情形,后一种情形同理, \\

当$\displaystyle\lambda>0$时,$\displaystyle f(L+\lambda)-f(L)\geq\varepsilon_0,$ \\

\phantom{当$\displaystyle\lambda>0$时,$\displaystyle\;$}$\displaystyle f(L+2\lambda)-f(L+\lambda)\geq\varepsilon_0,$ \\
\phantom{当$\displaystyle\lambda>0$时,$\displaystyle\;f(L+2\lambda)$}\vdots \\
\phantom{当$\displaystyle\lambda>0$时,$\displaystyle\;$}$\displaystyle f(L+n\lambda)-f(L+n\lambda-\lambda)\geq\varepsilon_0,$ \\

\phantom{当$\displaystyle\lambda>0\;$}$\displaystyle\therefore\,f(L+n\lambda)-f(L)\geq n\varepsilon_0,$ \\

\phantom{当$\displaystyle\lambda>0\;$}$\displaystyle\therefore\lim_{n\to\infty}f(L+n\lambda)=+\infty$,与$\displaystyle f(x)$在$\displaystyle[0,+\infty)$上有界矛盾, \\

当$\displaystyle\lambda<0$时,$\displaystyle f(L)-f(L-\lambda)\geq\varepsilon_0,$ \\

\phantom{当$\displaystyle\lambda<0$时,$\displaystyle\;$}$\displaystyle f(L-\lambda)-f(L-2\lambda)\geq\varepsilon_0,$ \\
\phantom{当$\displaystyle\lambda<0$时,$\displaystyle\;f(L-\lambda)$}\vdots \\
\phantom{当$\displaystyle\lambda<0$时,$\displaystyle\;$}$\displaystyle f(L-n\lambda+\lambda)-f(L-n\lambda)\geq\varepsilon_0,$ \\

\phantom{当$\displaystyle\lambda<0\;$}$\displaystyle\therefore\,f(L)-f(L-n\lambda)\geq n\varepsilon_0,$ \\

\phantom{当$\displaystyle\lambda<0\;$}$\displaystyle\therefore\lim_{n\to\infty}f(L-n\lambda)=-\infty$,与$\displaystyle f(x)$在$\displaystyle[0,+\infty)$上有界矛盾, \\

综上,$\forall\,\varepsilon>0,\forall\,M>0,\;\exists\,x\geq M,\;\text{s.t.}\,\left\lvert g(x)\right\rvert<\varepsilon,$ \\

取$\varepsilon=1,M=1,\;\exists\,x_1\geq 1,\text{且}\left\lvert g(x_1)\right\rvert<1,$ \\

取$\varepsilon=\dfrac{1}{2},M=2,\;\exists\,x_2\geq 2,\text{且}\left\lvert g(x_2)\right\rvert<\dfrac{1}{2},$ \\
\phantom{取$\varepsilon=1,M=1,$}\vdots \\
取$\varepsilon=\dfrac{1}{n},M=n,\;\exists\,x_n\geq n,\text{且}\left\lvert g(x_n)\right\rvert<\dfrac{1}{n},$ \\

$\therefore\displaystyle\lim_{n\to\infty}x_n=+\infty,\lim_{n\to\infty}\left\lvert f(x_n+\lambda)-f(x_n)\right\rvert=0,$ \\

综上,$\forall\,\lambda\in\mathbb{R},\;\exists\{x_n\},\;\text{s.t.}\,\displaystyle\lim_{n\to\infty}x_n=+\infty,\lim_{n\to\infty}[f(x_n+\lambda)-f(x_n)]=0.$

\begin{ti}
【题9】证明:不等于常数的连续周期函数一定有最小正周期
\end{ti}

\textbf{证明}:记集合$T=\left\{t>0\,|\;f(x+t)=f(x)\right\}$,函数$f(x)$为不等于常数的连续周期函数, \\

$\because\,f(x)$是周期函数,$\quad\therefore\,T$非空,又$t>0,\quad\therefore\,T$有下界, \\

由确界原理,$\exists\,T_0=\inf{T},T_0\geq 0$,下证:$T_0>0\text{且}T_0\in T,$ \\

\ding{172} 若$\displaystyle T_0=0$,则对每个$\displaystyle n\in\mathbb{N}^+$,可取$\displaystyle t_n\in T$使$\displaystyle0<t_n<\dfrac{1}{n}$, \\

任取$\displaystyle x_1<x_2$,令$\displaystyle k_n$为满足$\displaystyle k_nt_n\geq x_2-x_1$的最小正整数,并令$\displaystyle r_n=x_2-x_1-(k_n-1)t_n$, \\

则$\displaystyle0<r_n\leq t_n<\dfrac{1}{n}$,故$\displaystyle r_n\to0^+$,又$\displaystyle(k_n-1)t_n$为$\displaystyle f$的周期, \\

$\therefore\displaystyle f(x_2)=f(x_1+(k_n-1)t_n+r_n)=f(x_1+r_n)\to f(x_1),$ \\

故$\displaystyle f(x_2)=f(x_1)$.由$\displaystyle x_1,x_2$的任意性,$\displaystyle f$为常值函数,与题目条件矛盾!$\displaystyle\quad\therefore\,T_0>0,$ \\

\ding{173} $\because\,T_0=\inf{T},\;\therefore\exists\,\{t_n\}\subseteq T,\;\text{s.t.}\,\displaystyle\lim_{n\to\infty}t_n=T_0,$ \\

又$f(x)=f(x+t_n)\text{对}\forall\,n\in\mathbb{N}^+$成立,$\quad\therefore\,f(x)=\displaystyle\lim_{n\to\infty}f(x+t_n),$ \\

由$Heine$定理,$\displaystyle\lim_{n\to\infty}f(x+t_n)=f(x+T_0),\quad\therefore\,f(x)=f(x+T_0),\text{即}T_0\in T,$ \\

综上,$T_0$为$f(x)$的最小正周期,即$f(x)$有最小正周期. \\

\textbf{注记}:关键在于,当周期很小的时候,可以用整数倍周期很好的逼近任意实数.

\begin{ti}
【题10】设$f$在区间$[0,n]$连续,$f(0)=f(n),n\in\mathbb{N}^+$,证明:至少有$n$对不同的$(x,y)$使得$f(x)=f(y)$,同时$x-y$为非$0$整数
\end{ti}

\textbf{证明}:对$n$进行归纳,$n=1$时显然成立, \\

假设$n=k-1$时结论成立,$\;n=k$时,设$g(x)=f(x+1)-f(x),$ \\

则$g(0)+g(1)+...+g(n-1)=0,\;\therefore\,g(0),g(1),...g(n-1)$不能恒正或恒负, \\

$\therefore\exists\,k_1,k_2\in\{0,1,...,n-1\},\;\text{s.t.}\,g(k_1)\geq 0,g(k_2)\leq 0,k_1\neq k_2,$ \\

$\therefore\exists\,x_0\text{介于}k_1,k_2\text{之间},\;\text{s.t.}\,g(x_0)=0,$ \\

记$\displaystyle\phi(x)=\begin{cases}
    \raisebox{0.45ex}[0pt][0pt]{$\displaystyle x$}&\raisebox{0.45ex}[0pt][0pt]{$\displaystyle\quad x\in[0,x_0]$} \\
    \displaystyle x+1&\displaystyle\quad x\in(x_0,n-1]
\end{cases}$,并令$\displaystyle h(x)=f(\phi(x))$,
\\

因为$\displaystyle f(x_0)=f(x_0+1)$,所以$\displaystyle h\in C[0,n-1]$, \\

由归纳假设,$\displaystyle h$至少有$\displaystyle n-1$对不同的$\displaystyle(x_i,y_i)$,满足$\displaystyle h(x_i)=h(y_i)$且$\displaystyle x_i-y_i\in\mathbb{Z}\setminus\{0\}$, \\

易得函数$\displaystyle\phi$在$\displaystyle[0,n-1]$上严格递增,故为单射, \\

对每一对$\displaystyle(x_i,y_i)$,有$\displaystyle f(\phi(x_i))=f(\phi(y_i))$, \\

且$\displaystyle\phi(x_i)-\phi(y_i)$等于$\displaystyle x_i-y_i$加上$\displaystyle-1,0,1$中的一个,因而仍为整数; \\

又由单射性可知该整数非零, 单射性还保证不同点对的像仍不相同, \\

又$\displaystyle\phi([0,n-1])=[0,x_0]\cup(x_0+1,n]$,所以$\displaystyle x_0+1$不在像集中, \\

从而新增点对$\displaystyle(x_0,x_0+1)$不与上述$\displaystyle n-1$对重复, \\

因此$\displaystyle f$至少有$\displaystyle n$对不同的$\displaystyle(x,y)$满足$\displaystyle f(x)=f(y)$且$\displaystyle x-y$为非零整数,\\

由数学归纳法可知,至少有$\displaystyle n$对不同的$\displaystyle(x,y)$使得$\displaystyle f(x)=f(y)$,同时$\displaystyle x-y$为非$\displaystyle0$整数.



\BookEnd

\ifdefined\BookMaster
\def\BookEnd{}
\else
\documentclass[a4paper,12pt,fleqn,openany,twoside,fontset=windows]{ctexbook}

\usepackage[fleqn]{amsmath}
\usepackage{amssymb,amsthm,mathrsfs}
\usepackage{geometry}
\usepackage[dvipsnames]{xcolor}
\usepackage{graphicx}
\usepackage{tikz}
\usetikzlibrary{calc}
\usepackage[most]{tcolorbox}
\usepackage{fancyhdr}
\usepackage{titlesec}
\usepackage{tocloft}
\usepackage{enumitem}
\usepackage{microtype}
\usepackage{pifont}
\usepackage{hyperref}
\usepackage{romannum}
\usepackage{mathtools}

\definecolor{bookblue}{HTML}{264653}
\definecolor{bookteal}{HTML}{4F7C82}
\definecolor{booklight}{HTML}{EDF4F3}
\definecolor{bookgray}{HTML}{5E686B}

\geometry{
    inner=2.7cm,
    outer=2.5cm,
    top=2.7cm,
    bottom=2.7cm,
    headheight=18pt,
    headsep=18pt,
    footskip=25pt
}
\setlength{\mathindent}{0pt}
\setlength{\parindent}{5pt}
\setlength{\parskip}{3pt plus 1pt minus 1pt}
\linespread{1.25}
\allowdisplaybreaks[3]
\AtBeginDocument{%
    \setlength{\parindent}{5pt}
    \setlength{\abovedisplayskip}{8pt plus 2pt minus 2pt}
    \setlength{\belowdisplayskip}{8pt plus 2pt minus 2pt}
    \setlength{\abovedisplayshortskip}{5pt plus 1pt minus 1pt}
    \setlength{\belowdisplayshortskip}{5pt plus 1pt minus 1pt}
}

\hypersetup{
    unicode=true,
    colorlinks=true,
    linkcolor=bookblue,
    citecolor=bookblue,
    urlcolor=bookteal,
    bookmarksopen=true,
    pdfauthor={Chen},
    pdftitle={数学分析习题整理}
}

\renewcommand{\chaptermark}[1]{%
    \edef\@tmp{第\thechapter 章\quad #1}%
    \expandafter\markboth\expandafter{\@tmp}{}%
}
\fancypagestyle{main}{%
    \fancyhf{}
    \fancyhead[LE]{\small\color{bookgray}\nouppercase{\leftmark}}
    \fancyhead[RO]{\small\color{bookgray}数学分析习题整理}
    \fancyfoot[C]{\color{bookblue}\thepage}
    \fancyfoot[R]{\small\bfseries Chen\;\; github.com/mathlover-1/math-analysis}
    \renewcommand{\headrulewidth}{0.4pt}
    \renewcommand{\footrulewidth}{0pt}
}
\fancypagestyle{plain}{%
    \fancyhf{}
    \fancyfoot[C]{\color{bookblue}\thepage}
    \fancyfoot[R]{\small\bfseries Chen\;\; github.com/mathlover-1/math-analysis}
    \renewcommand{\headrulewidth}{0pt}
    \renewcommand{\footrulewidth}{0pt}
}
\pagestyle{main}

\titleformat{\chapter}[display]
  {\normalfont\sffamily\bfseries\color{bookblue}}
  {\raggedleft\fontsize{46}{50}\selectfont\thechapter}
  {-18pt}
  {\titlerule[1.2pt]\vspace{10pt}\filright\Huge}
\titlespacing*{\chapter}{0pt}{-15pt}{36pt}

\renewcommand{\contentsname}{目\quad 录}
\renewcommand{\cfttoctitlefont}{\hfill\sffamily\bfseries\Huge\color{bookblue}}
\renewcommand{\cftaftertoctitle}{\hfill}
\setlength{\cftbeforetoctitleskip}{5pt}
\setlength{\cftaftertoctitleskip}{28pt}
\setlength{\cftbeforechapskip}{12pt}
\renewcommand{\cftchapfont}{\sffamily\bfseries\large\color{bookblue}}
\renewcommand{\cftchappagefont}{\sffamily\bfseries\large\color{bookblue}}
\renewcommand{\cftchapleader}{\color{bookteal}\cftdotfill{2.5}}

\newtcolorbox{ti}{
    enhanced,
    breakable,
    colback=booklight!55!white,
    colframe=bookteal!35!white,
    boxrule=0.45pt,
    boxsep=0pt,
    left=10pt,
    right=10pt,
    top=9pt,
    bottom=9pt,
    before skip=14pt,
    after skip=14pt,
    arc=2pt,
    outer arc=2pt,
    before upper={\parindent=0pt}
}

\newenvironment{booknotebody}{%
    \begin{center}
    \begin{minipage}{0.84\textwidth}
    \songti\zihao{-4}\linespread{1.45}\selectfont
    \setlength{\parindent}{2\ccwd}
    \setlength{\parskip}{0.75em}
}{%
    \end{minipage}
    \end{center}
}

\newenvironment{booknotelongbody}{%
    \begingroup
    \songti\zihao{-4}\linespread{1.45}\selectfont
    \setlength{\parindent}{2\ccwd}
    \setlength{\parskip}{0.75em}
    \begin{list}{}{%
        \setlength{\leftmargin}{0.08\textwidth}
        \setlength{\rightmargin}{0.08\textwidth}
        \setlength{\topsep}{0pt}
        \setlength{\partopsep}{0pt}
        \setlength{\parsep}{0pt}
        \setlength{\itemsep}{0pt}
    }
    \item\relax
}{%
    \end{list}
    \endgroup
}

\fancypagestyle{plainnowm}{%
    \fancyhf{}
    \fancyfoot[C]{\color{bookblue}\thepage}
    \renewcommand{\headrulewidth}{0pt}
    \renewcommand{\footrulewidth}{0pt}
}

\newcommand{\booknotetitle}[2][plain]{%
    \clearpage
    \phantomsection
    \addcontentsline{toc}{chapter}{#2}
    \markboth{#2}{}
    \thispagestyle{#1}
    \vspace*{1.2cm}
    \begin{center}
        {\sffamily\heiti\bfseries\fontsize{26}{34}\selectfont\color{bookblue}#2\par}
        \vspace{0.8em}
        {\color{bookteal}\rule{4.5cm}{0.8pt}\par}
    \end{center}
    \vspace{0.6cm}
}

\renewcommand{\proofname}{证明}
\renewcommand{\qedsymbol}{\color{bookteal}\ensuremath{\blacksquare}}

\title{数学分析习题整理}
\author{Chen}
\date{}

\makeatletter
% ==================== 旧版封面/封底设计备份 ====================
% 2026-08 重设计,旧代码用 \iffalse...\fi 封存,新版见下方。
\iffalse
\renewcommand{\maketitle}{%
    \begin{titlepage}
        \thispagestyle{empty}
        \setcounter{page}{0}
        \noindent\hspace*{-\parindent}\begin{tikzpicture}[remember picture,overlay]
            \fill[bookblue] (current page.north west) rectangle ([yshift=-4.2cm]current page.north east);
            \fill[booklight] ([yshift=-4.2cm]current page.north west) rectangle (current page.south east);
            \draw[bookteal!55,line width=0.7pt]
                ([xshift=2.4cm,yshift=4.3cm]current page.south west)
                .. controls +(2.5cm,1.8cm) and +(-2.8cm,-1.3cm) ..
                ([xshift=-2.4cm,yshift=7.2cm]current page.south east);
            \draw[bookteal!35,line width=0.7pt]
                ([xshift=2.5cm,yshift=6.2cm]current page.south west)
                .. controls +(3.0cm,-1.2cm) and +(-3.0cm,1.6cm) ..
                ([xshift=-2.3cm,yshift=3.8cm]current page.south east);
        \end{tikzpicture}
        \vspace*{5.6cm}
        \noindent\color{bookblue}\rule{\textwidth}{1.4pt}\par
        \vspace{8pt}
        {\centering\sffamily\bfseries\fontsize{32}{42}\selectfont\color{bookblue}\@title\par}
        \vspace{8pt}
        \noindent\color{bookteal}\rule{\textwidth}{0.6pt}\par
        \vspace{2.6cm}
        {\centering\Large\color{bookgray}\@author\par}
        \vfill
        {\centering\color{bookteal}\large
            $\displaystyle \lim_{n\to\infty} a_n$\qquad
            $\displaystyle \int_a^b f(x)\,\mathrm{d}x$\qquad
            $\displaystyle \nabla f(x,y)$\par
        }
        \vspace*{1.5cm}
    \end{titlepage}
}

\newcommand{\makebackcover}{%
    \begingroup
    \let\cleardoublepage\clearpage
    \begin{titlepage}
        \thispagestyle{empty}
        \noindent\hspace*{-\parindent}\begin{tikzpicture}[remember picture,overlay]
            \fill[bookblue] (current page.north west) rectangle (current page.south east);
            \fill[booklight,rounded corners=7pt]
                ([xshift=2.2cm,yshift=-3.0cm]current page.north west)
                rectangle
                ([xshift=-2.2cm,yshift=3.0cm]current page.south east);
            \fill[bookteal!80]
                (current page.north west) --
                ([xshift=5.3cm]current page.north west) --
                ([xshift=2.1cm,yshift=-4.4cm]current page.north west) --
                ([yshift=-6.2cm]current page.north west) -- cycle;
            \fill[bookteal!55]
                (current page.south east) --
                ([xshift=-5.8cm]current page.south east) --
                ([xshift=-2.4cm,yshift=4.2cm]current page.south east) --
                ([yshift=6.0cm]current page.south east) -- cycle;
            \foreach \radius in {0.75,1.15,1.55,1.95} {
                \draw[bookteal!55,line width=0.55pt]
                    ([xshift=-2.9cm,yshift=-4.4cm]current page.north east)
                    circle[radius=\radius cm];
            }
            \foreach \radius in {0.65,1.05,1.45} {
                \draw[booklight!65,line width=0.55pt]
                    ([xshift=2.4cm,yshift=3.6cm]current page.south west)
                    circle[radius=\radius cm];
            }
            \fill[bookteal!16,rounded corners=3pt]
                ([xshift=3.5cm,yshift=-7.1cm]current page.north west)
                rectangle
                ([xshift=-5.0cm,yshift=7.2cm]current page.south east);
            \fill[bookteal!30,rounded corners=3pt]
                ([xshift=4.4cm,yshift=-8.0cm]current page.north west)
                rectangle
                ([xshift=-4.1cm,yshift=8.1cm]current page.south east);
            \begin{scope}[shift={([xshift=-0.75cm]current page.center)}]
                \clip[rounded corners=3pt] (-5.25,-4.75) rectangle (5.25,4.75);
                \foreach \v in {-3.6,-3.0,...,3.6} {
                    \draw[bookteal!48,line width=0.45pt,opacity=0.72]
                        plot[domain=-3.6:3.6,samples=45,smooth,variable=\u]
                        ({\u+0.55*\v},{0.28*\u-0.35*\v+0.11*(\u*\u-\v*\v)});
                }
                \foreach \u in {-3.6,-3.0,...,3.6} {
                    \draw[bookblue!42,line width=0.42pt,opacity=0.62]
                        plot[domain=-3.6:3.6,samples=45,smooth,variable=\v]
                        ({\u+0.55*\v},{0.28*\u-0.35*\v+0.11*(\u*\u-\v*\v)});
                }
                \draw[bookteal!70,line width=0.75pt,opacity=0.82]
                    plot[domain=-3.6:3.6,samples=55,smooth,variable=\u]
                    ({\u},{0.28*\u+0.11*\u*\u});
                \draw[bookblue!60,line width=0.7pt,opacity=0.75]
                    plot[domain=-3.6:3.6,samples=55,smooth,variable=\v]
                    ({0.55*\v},{-0.35*\v-0.11*\v*\v});
            \end{scope}
            \draw[bookteal!45,line width=0.65pt]
                ([xshift=3.0cm,yshift=-8.2cm]current page.north west)
                .. controls +(2.2cm,2.0cm) and +(-2.0cm,-2.1cm) ..
                ([xshift=-3.0cm,yshift=9.0cm]current page.south east);
            \draw[bookteal!25,line width=0.55pt]
                ([xshift=3.1cm,yshift=8.0cm]current page.south west)
                .. controls +(2.7cm,-1.7cm) and +(-2.6cm,1.8cm) ..
                ([xshift=-3.1cm,yshift=-8.7cm]current page.north east);
            \node[anchor=west,text=bookteal!72,scale=0.78] at
                ([xshift=2.8cm,yshift=-4.0cm]current page.north west) {%
                $\displaystyle e^{\mathrm{i}\pi}+1=0$};
            \node[anchor=east,text=bookblue!62,scale=0.67] at
                ([xshift=-2.7cm,yshift=4.0cm]current page.south east) {%
                $\displaystyle
                f(x)=\sum_{n=0}^{\infty}\frac{f^{(n)}(a)}{n!}(x-a)^n$};
        \end{tikzpicture}
        \mbox{}
    \end{titlepage}
    \endgroup
}
\fi

% ==================== 新版封面(2026-08 重设计) ====================
% 设计要点:顶部深绿纵向渐变+波浪下缘+两层半透明过渡波(解决深浅绿生硬过渡);
% 深色区螺旋线/同心圆母题;标题+英文副题+作者居中;
% 中部 Fourier 部分和逼近方波曲线;底部浅色波纹(无文字)。
\renewcommand{\maketitle}{%
    \begin{titlepage}
        \thispagestyle{empty}
        \setcounter{page}{0}
        \noindent\hspace*{-\parindent}\begin{tikzpicture}[remember picture,overlay]
            % ---------- 底色 ----------
            \fill[booklight] (current page.north west) rectangle (current page.south east);

            % ---------- 顶部深色区:纵向渐变 + 波浪下缘 ----------
            \shade[top color=bookblue, bottom color=bookteal!55!bookblue]
                (current page.north west) -- (current page.north east) --
                ([yshift=-8.6cm]current page.north east)
                .. controls ([xshift=-5.2cm,yshift=-10.6cm]current page.north east)
                         and ([xshift=6.0cm,yshift=-8.0cm]current page.north west) ..
                ([yshift=-10.0cm]current page.north west) -- cycle;

            % ---------- 过渡层 1:半透明波浪 ----------
            \fill[bookteal!60!bookblue,opacity=0.42]
                ([yshift=-8.6cm]current page.north east)
                .. controls ([xshift=-5.2cm,yshift=-10.6cm]current page.north east)
                         and ([xshift=6.0cm,yshift=-8.0cm]current page.north west) ..
                ([yshift=-10.0cm]current page.north west) --
                ([yshift=-12.6cm]current page.north west)
                .. controls ([xshift=6.5cm,yshift=-11.0cm]current page.north west)
                         and ([xshift=-5.5cm,yshift=-13.2cm]current page.north east) ..
                ([yshift=-11.4cm]current page.north east) -- cycle;

            % ---------- 过渡层 2:更浅的波浪 ----------
            \fill[bookteal!35,opacity=0.35]
                ([yshift=-10.6cm]current page.north east)
                .. controls ([xshift=-6.0cm,yshift=-12.4cm]current page.north east)
                         and ([xshift=5.0cm,yshift=-9.8cm]current page.north west) ..
                ([yshift=-11.8cm]current page.north west) --
                ([yshift=-13.8cm]current page.north west)
                .. controls ([xshift=5.5cm,yshift=-12.6cm]current page.north west)
                         and ([xshift=-6.5cm,yshift=-14.6cm]current page.north east) ..
                ([yshift=-12.8cm]current page.north east) -- cycle;

            % ---------- 深色区装饰:右上螺旋线 ----------
            \begin{scope}[shift={([xshift=-3.4cm,yshift=-4.3cm]current page.north east)}]
                \draw[booklight!45!bookteal,opacity=0.55,line width=0.6pt]
                    plot[domain=0:1440,samples=220,smooth,variable=\t]
                    ({0.0028*\t*cos(\t)},{0.0028*\t*sin(\t)});
            \end{scope}

            % ---------- 深色区装饰:左下同心圆 ----------
            \foreach \r in {0.7,1.25,1.8,2.35,2.9}{
                \draw[booklight!35!bookteal,opacity=0.5,line width=0.5pt]
                    ([xshift=0.4cm,yshift=-7.9cm]current page.north west) circle[radius=\r cm];
            }

            % ---------- 主标题 ----------
            \node[anchor=center,text=booklight] at
                ([yshift=-4.5cm]current page.north)
                {\sffamily\bfseries\fontsize{40}{46}\selectfont \@title};

            % ---------- 分隔短线 ----------
            \draw[booklight!65,line width=1.2pt]
                ([xshift=-2.6cm,yshift=-5.75cm]current page.north) --
                ([xshift=2.6cm,yshift=-5.75cm]current page.north);

            % ---------- 英文副题 ----------
            \node[anchor=center,text=booklight!70!bookteal] at
                ([yshift=-6.55cm]current page.north)
                {\itshape\large Mathematical Analysis \textperiodcentered\
                 A Collection of Exercises};

            % ---------- 作者(紧随标题下方) ----------
            \node[anchor=center,text=booklight!88] at
                ([yshift=-8.2cm]current page.north)
                {\sffamily\large \@author\quad 整理};

            % ---------- 中部:Fourier 部分和逼近方波 ----------
            \begin{scope}[shift={([xshift=-5.2cm,yshift=-18.1cm]current page.north)}]
                % 坐标轴
                \draw[bookgray!55,line width=0.5pt,->] (-0.5,0) -- (10.9,0);
                \draw[bookgray!55,line width=0.5pt,->] (0,-1.75) -- (0,1.95);
                % 目标方波(虚线)
                \draw[bookgray!75,line width=0.6pt,dashed]
                    (0,0.92) -- (5,0.92) (5,-0.92) -- (10,-0.92);
                % N=1
                \draw[bookteal!45,line width=0.7pt]
                    plot[domain=0:360,samples=90,smooth,variable=\x]
                    ({\x/36},{1.18*sin(\x)});
                % N=3
                \draw[bookteal!80,line width=0.8pt]
                    plot[domain=0:360,samples=120,smooth,variable=\x]
                    ({\x/36},{1.18*(sin(\x)+sin(3*\x)/3)});
                % N=7
                \draw[bookblue,line width=0.95pt]
                    plot[domain=0:360,samples=150,smooth,variable=\x]
                    ({\x/36},{1.18*(sin(\x)+sin(3*\x)/3+sin(5*\x)/5+sin(7*\x)/7)});
                % 刻度标注
                \node[anchor=north east,text=bookgray!80,scale=0.62] at (-0.12,0) {$0$};
                \node[anchor=north,text=bookgray!80,scale=0.62] at (5,0) {$\pi$};
                \node[anchor=north,text=bookgray!80,scale=0.62] at (10,0) {$2\pi$};
            \end{scope}
            % 图注
            \node[anchor=center,text=bookgray,scale=0.8] at
                ([yshift=-20.5cm]current page.north)
                {$S_N(x)=\displaystyle\sum_{k=1}^{N}\frac{\sin\bigl((2k-1)x\bigr)}{2k-1}$
                 \ 逐步逼近方波};

            % ---------- 底部浅色波纹(无文字,与顶部波浪呼应) ----------
            \fill[bookteal!22,opacity=0.5]
                ([yshift=2.4cm]current page.south west)
                .. controls ([xshift=6.0cm,yshift=3.9cm]current page.south west)
                         and ([xshift=-6.5cm,yshift=1.0cm]current page.south east) ..
                ([yshift=2.9cm]current page.south east) --
                (current page.south east) -- (current page.south west) -- cycle;
            \fill[bookteal!40!bookblue,opacity=0.28]
                ([yshift=1.3cm]current page.south west)
                .. controls ([xshift=5.5cm,yshift=2.7cm]current page.south west)
                         and ([xshift=-5.5cm,yshift=0.3cm]current page.south east) ..
                ([yshift=1.9cm]current page.south east) --
                (current page.south east) -- (current page.south west) -- cycle;
        \end{tikzpicture}
        \mbox{}
    \end{titlepage}
}

% ==================== 新版封底(纯背景,2026-08 重设计) ====================
% 设计要点:对角渐变底+角部半透明叠层;中下部外摆线(玫瑰形)为视觉主体;
% 底部两层半透明白色波浪;无任何文字。
\newcommand{\makebackcover}{%
    \begingroup
    \let\cleardoublepage\clearpage
    \begin{titlepage}
        \thispagestyle{empty}
        \noindent\hspace*{-\parindent}\begin{tikzpicture}[remember picture,overlay]
            % ---------- 底色:对角渐变 ----------
            \shade[top color=bookblue, bottom color=bookteal!48!bookblue,
                   shading angle=-30]
                (current page.north west) rectangle (current page.south east);

            % ---------- 角部半透明叠层 ----------
            \fill[white,opacity=0.05]
                ([xshift=2.0cm,yshift=1.0cm]current page.north east) circle[radius=7.2cm];
            \fill[white,opacity=0.06]
                ([xshift=-0.5cm,yshift=-1.5cm]current page.north east) circle[radius=3.6cm];

            % ---------- 左下同心圆 ----------
            \foreach \r in {1.0,1.75,2.5,3.25,4.0,4.75}{
                \draw[booklight!40,opacity=0.55,line width=0.55pt]
                    ([xshift=-0.6cm,yshift=-0.4cm]current page.south west) circle[radius=\r cm];
            }
            % ---------- 右下螺旋线 ----------
            \begin{scope}[shift={([xshift=-2.9cm,yshift=3.2cm]current page.south east)}]
                \draw[booklight!45!bookteal,opacity=0.5,line width=0.55pt]
                    plot[domain=0:1260,samples=200,smooth,variable=\t]
                    ({0.0036*\t*cos(\t)},{0.0036*\t*sin(\t)});
            \end{scope}
            % ---------- 左上斜线组 ----------
            \foreach \i in {0,1,2}{
                \draw[bookteal!60!booklight,opacity=0.5,line width=0.6pt]
                    ([xshift={1.6cm+\i*0.45cm}]current page.north west) --
                    ([yshift={-1.6cm-\i*0.45cm}]current page.north west);
            }

            % ---------- 中下部视觉主体:外摆线(玫瑰形) ----------
            \begin{scope}[shift={([yshift=-16.6cm]current page.north)}]
                % 外圈
                \draw[booklight!35,opacity=0.6,line width=0.6pt]
                    (0,0) circle[radius=4.9cm];
                \draw[booklight!22,opacity=0.5,line width=0.5pt]
                    (0,0) circle[radius=5.25cm];
                % 主玫瑰线
                \draw[booklight!55!bookteal,opacity=0.8,line width=0.85pt]
                    plot[domain=0:1080,samples=400,smooth,variable=\t]
                    ({0.335*(8*cos(\t)-5*cos(2.6667*\t))},
                     {0.335*(8*sin(\t)-5*sin(2.6667*\t))});
                % 内层玫瑰线
                \draw[booklight!38!bookteal,opacity=0.6,line width=0.6pt]
                    plot[domain=0:1080,samples=400,smooth,variable=\t]
                    ({0.21*(8*cos(\t)-5*cos(2.6667*\t))},
                     {0.21*(8*sin(\t)-5*sin(2.6667*\t))});
            \end{scope}

            % ---------- 底部半透明波浪 ----------
            \fill[white,opacity=0.06]
                ([yshift=4.6cm]current page.south west)
                .. controls ([xshift=6.5cm,yshift=6.4cm]current page.south west)
                         and ([xshift=-6.0cm,yshift=2.6cm]current page.south east) ..
                ([yshift=5.2cm]current page.south east) --
                (current page.south east) -- (current page.south west) -- cycle;
            \fill[booklight,opacity=0.08]
                ([yshift=2.6cm]current page.south west)
                .. controls ([xshift=5.5cm,yshift=4.3cm]current page.south west)
                         and ([xshift=-6.5cm,yshift=1.0cm]current page.south east) ..
                ([yshift=3.3cm]current page.south east) --
                (current page.south east) -- (current page.south west) -- cycle;
        \end{tikzpicture}
        \mbox{}
    \end{titlepage}
    \endgroup
}
\makeatother

% 页脚右下方常态化水印(署名 + 仓库地址,加粗黑色),已在 main/plain 分页样式中定义,所有页面均生效.

\begin{document}
\mainmatter
\pagestyle{main}
\setcounter{chapter}{3}
\def\BookEnd{\end{document}}
\fi

\chapter{一元微分学与一元积分学}
\begin{ti}
【拓展1】$Rolle$定理的$Samuleson$证明
\end{ti}

\textbf{证明}:设$f\in C^0[a,b],\text{且在开区间}(a,b)\text{可导},f(a)=f(b),\text{则要证的就是}\;\exists\xi,\;\text{s.t.}\,f'(\xi)=0,$ \\

构造$F(x)=f(x)-f(x+\dfrac{b-a}{2}),F\in C[a,\dfrac{a+b}{2}],F(a)=-F(\dfrac{a+b}{2}),$ \\

$\therefore\exists\,a_1\in[a,\dfrac{a+b}{2}],\;\text{s.t.}\,F(a_1)=0,$ \\

记$b_1=a_1+\dfrac{b-a}{2},\,\,\text{则}f(a_1)=f(b_1),\,\,\text{且}b_1-a_1=\dfrac{b-a}{2},$ \\

以此类推地做下去,可以得到闭区间套$\{[a_n,b_n]\},\text{且}f(a_n)=f(b_n),b_n-a_n=\dfrac{b-a}{2^n},$ \\

由闭区间套原理,这些闭区间有唯一公共点$\xi,\text{且}\displaystyle\lim_{n\to\infty}a_n=\lim_{n\to\infty}b_n=\xi,a_n\leq\xi\leq b_n,$ \\

事实上,总可以使得$\displaystyle\xi\in(a,b)$.\\

若$\displaystyle\xi=a$,则由$\displaystyle a\leq a_n\leq\xi=a$可知$\displaystyle a_n=a$对每个$\displaystyle n$成立, \\

从而$\displaystyle f(b_1)=f(a)=f(b_2)$, \\

又$\displaystyle b_2=a+\dfrac{b-a}{4}<b_1=a+\dfrac{b-a}{2}$,故$\displaystyle[b_2,b_1]\subset(a,b)$, \\

以$\displaystyle[b_2,b_1]$为新的初始区间,对$\displaystyle f$重新进行上述构造, \\

所得闭区间套的公共点属于$\displaystyle[b_2,b_1]\subset(a,b)$, \\

若$\displaystyle\xi=b$,则同理由$\displaystyle\xi=b\leq b_n\leq b$可知$\displaystyle b_n=b$对每个$\displaystyle n$成立, \\

从而$\displaystyle f(a_1)=f(b)=f(a_2)$,且$\displaystyle[a_1,a_2]\subset(a,b)$, \\

以$\displaystyle[a_1,a_2]$为新的初始区间重新进行上述构造即可, \\

因此总能得到公共点$\displaystyle\xi\in(a,b)$的闭区间套, \\

下面证明一个引理:若有$x_n\leq c\leq y_n,\displaystyle\lim_{n\to\infty}x_n=\lim_{n\to\infty}y_n=c,f\text{在}c\text{处可导},$ \\
$\displaystyle\text{则}f'(c)=\lim_{n\to\infty}\dfrac{f(x_n)-f(y_n)}{x_n-y_n}.$ \\

不妨设$f'(c)=0,$否则令$g(x)=f(x)-f'(c)x,\;\text{则}g'(c)=0,$ \\

$\displaystyle\lim_{n\to\infty}\dfrac{f(x_n)-f(y_n)}{x_n-y_n}=\lim_{n\to\infty}\dfrac{g(x_n)-g(y_n)}{x_n-y_n}+f'(c),$ \\

$\because f'(c)=0,\quad\therefore\forall\,\varepsilon>0,\;\exists\,N,\;n>N\text{时},$ \\

$\left\lvert f(c)-f(x_n)\right\rvert<\varepsilon(c-x_n),\left\lvert f(y_n)-f(c)\right\rvert<\varepsilon(y_n-c),$ \\

$\therefore\left\lvert f(x_n)-f(y_n)\right\rvert\leq\left\lvert f(c)-f(x_n)\right\rvert+\left\lvert f(y_n)-f(c)\right\rvert<\varepsilon(y_n-x_n),$ \\

即$\forall\,\varepsilon>0,\;\exists\,N,\;n>N\text{时},\;\left\lvert \dfrac{f(x_n)-f(y_n)}{x_n-y_n}\right\rvert<\varepsilon,$ \\

$\therefore\displaystyle\lim_{n\to\infty}\dfrac{f(x_n)-f(y_n)}{x_n-y_n}=f'(c)=0,$得证, \\
 
由引理,$f'(\xi)=\displaystyle\lim_{n\to\infty}\dfrac{f(a_n)-f(b_n)}{a_n-b_n}=0$,命题得证. \\

\textbf{注}:以上证明并未用到最值定理与极值的思想,事实上,由$Samuleson$证明得到的$\xi$不一定是极值点, \\

如$f(x)=\begin{cases}
    x^3\sin{\dfrac{1}{x}}&\quad x\in[-\dfrac{1}{\pi},0)\cup(0,\dfrac{1}{\pi}] \\
    0&\quad x=0 \\
\end{cases}$

\begin{ti}
【拓展2】$Lagrange$定理的面积证法
\end{ti}

\textbf{证明}:令$S(x)=\dfrac{1}{2}\begin{vmatrix}
    1 &1 &1 \\
    a &b &x \\
    f(a) &f(b) &f(x)
\end{vmatrix},$ \\

则$S(x)=\dfrac{1}{2}\begin{vmatrix}
    1 &0 &0 \\
    a &b-a &x-a \\
    f(a) &f(b)-f(a) &f(x)-f(a)
\end{vmatrix}=\dfrac{1}{2}\begin{vmatrix}
     b-a &x-a \\
     f(b)-f(a) &f(x)-f(a)
\end{vmatrix},$ \\

$\therefore\,S(a)=S(b)=0,$由$Rolle$定理,$\exists\,\xi\in(a,b),\;\text{s.t.}\,S'(\xi)=0,$ \\

由行列式求导法则,$2S'(x)=\begin{vmatrix}
    0 &x-a \\
    0 &f(x)-f(a)
\end{vmatrix}+\begin{vmatrix}
    b-a &1 \\
    f(b)-f(a) &f'(x)
\end{vmatrix},$ \\

$\therefore\,2S'(\xi)=\begin{vmatrix}
    b-a &1 \\
    f(b)-f(a) &f'(\xi)
\end{vmatrix}=0\Rightarrow f'(\xi)=\dfrac{f(b)-f(a)}{b-a}.$ \\

\begin{minipage}[t]{0.59\textwidth}
\textbf{注}:这个证法来自于几何直观,我们要找函数上切线与给定连线平行的点 ,$S(x)$是$(a,f(a)),(b,f(b)),(x,f(x))$三点连成的三角形的有向面积,当$(x,f(x))$沿着平行于$(a,f(a)),(b,f(b))$的连线的直线移动时,$S(x)$不变,
猜想:$S(x)$的驻点处,切线与该连线平行
\end{minipage}\hfill
\begin{minipage}[t]{0.37\textwidth}
\begin{tikzpicture}[x=0.72cm,y=0.72cm,baseline=(current bounding box.north)]
    \draw[bookteal!55,line width=1pt] (0.35,0.35) -- (2.90,2.05) -- (4.35,0.75) -- cycle;
    \draw[bookteal!75,line width=0.8pt]
        (0.35,0.35) .. controls (0.95,1.35) and (2.30,1.99) .. (2.90,2.05)
        .. controls (3.12,2.072) and (3.25,2.08) .. (3.35,2.02)
        .. controls (3.72,1.80) and (4.08,1.02) .. (4.35,0.75);
    \draw[bookgray!80,dashed,line width=0.55pt]
        (1.80,1.94) -- (3.80,2.14);
    \fill[bookblue] (0.35,0.35) circle (1.4pt);
    \fill[bookblue] (2.90,2.05) circle (1.4pt);
    \fill[bookblue] (4.35,0.75) circle (1.4pt);
    \node[below left,font=\scriptsize] at (0.35,0.35) {$(a,f(a))$};
    \node[above left,font=\scriptsize] at (2.90,2.05) {$(x,f(x))$};
    \node[below right,font=\scriptsize] at (4.35,0.75) {$(b,f(b))$};
\end{tikzpicture}

\vspace{2pt}
{\centering\small\color{bookgray}图4-1\par}
\end{minipage}
\par

\begin{ti}
【题3】设$f\in C[0,1],\text{在}(0,1)\text{可微},f(0)=0,f(1)=1,\text{设}k_1,k_2,...,k_n\text{是满足}$ \\
$k_1+k_2+...+k_n=1$的$n$个正数,证明:在$(0,1)$中存在$n$个互不相同的数$t_1,t_2,...,t_n,\text{使得}\dfrac{k_1}{f'(t_1)}+...+\dfrac{k_n}{f'(t_n)}=1$
\end{ti}

\textbf{证明}:由介值定理,$\exists\,x_1\in(0,1),\;\text{s.t.}\,f(x_1)=k_1,$ \\

同理$\exists\,x_2\in(x_1,1),\;\text{s.t.}\,f(x_2)=k_1+k_2,$ \\

以此类推,可以在$[0,1]$中插入一系列数:$0=x_0<x_1<x_2<...<x_{n-1}<x_n=1,$ \\

同时$f(x_m)=\displaystyle\sum_{i=1}^{m}k_i$,在$[x_{i-1},x_i]\text{上,由}Lagrange$中值定理, \\

$\exists\,t_i\in(x_{i-1},x_i),\;\text{s.t.}\,f(x_i)-f(x_{i-1})=f'(t_i)(x_i-x_{i-1}),i=1,2,...,n,$ \\

$\therefore\displaystyle\sum_{i=1}^{n}\dfrac{k_i}{f'(t_i)}=\sum_{i=1}^{n}(x_i-x_{i-1})=1,\text{且}t_1,t_2,...,t_n$互不相同,得证.

\begin{ti}
【题4】设$f(x)=\displaystyle\sum_{k=1}^{n}c_ke^{\lambda_k x},\text{其中}\lambda_1,\lambda_2,...,\lambda_n$为互异实数,$c_1,c_2,...,c_n$不同时为$0$, \\

证明:$f$的零点个数小于$n$
\end{ti}

\textbf{证明}:$n=1\text{时},f(x)=c_1e^{\lambda_1 x}$无零点,符合要求, \\

假设$n=m$时结论成立,$n=m+1$时,令$g(x)=e^{-\lambda_{m+1} x}f(x)=\displaystyle\sum_{k=1}^{m+1}c_ke^{(\lambda_k-\lambda_{m+1})x},$ \\

则$g'(x)=\displaystyle\sum_{k=1}^{m+1}c_k(\lambda_k-\lambda_{m+1})e^{(\lambda_k-\lambda_{m+1})x}=\sum_{k=1}^{m}c_k(\lambda_k-\lambda_{m+1})e^{(\lambda_k-\lambda_{m+1})x},$ \\

记$d_k=c_k(\lambda_k-\lambda_{m+1}),\;\lambda_k'=\lambda_k-\lambda_{m+1},$ \\

那么$d_1,...,d_m$不同时为$0$,$\lambda_1',...,\lambda_m'$为互异实数,又$g'(x)=\displaystyle\sum_{k=1}^{m}c_k(\lambda_k-\lambda_{m+1})e^{(\lambda_k-\lambda_{m+1})x}\rlap{,}$ \\

由归纳假设可知,$g'(x)$零点个数小于$m$,\\

但如果$f$有大于等于$m+1$个零点,则$g$也有大于等于$m+1$个零点, \\

从而$g'$有至少$m$个零点,矛盾!$\quad\therefore f$的零点个数小于$m+1,$ \\

由数学归纳法可知,$f$的零点个数小于$n.$

\begin{ti}
【题5】$\;f\text{在}[a,b]\text{上连续},(a,b)\text{上二阶可导},f(a)=f(b)=0,\int_{a}^{b}f(x)\,dx=0,$ \\
证明:$\exists\,\zeta\in[a,b],\;\text{s.t.}f(\zeta)=f''(\zeta)$
\end{ti}

\textbf{证明}:若$f(x)_{min}=0,\text{则由}\int_{a}^{b}f(x)\,dx=0\Rightarrow f\equiv 0,$ \\

则$\forall\zeta\in(a,b),\text{都有}f(\zeta)=f''(\zeta)$,同理$f(x)_{max}=0$时也有$f(\zeta)=f''(\zeta),$ \\

下设$f(x_1)=f(x)_{min}<0,f(x_2)=f(x)_{max}>0,x_1,x_2\in(a,b),$ \\

由零值定理,$\exists\,x_0\in(a,b),\;\text{s.t.}f(x_0)=0,$ \\

令$G(x)=e^xf(x),\text{则}G(a)=G(x_0)=G(b)=0,G'(x)=e^x[f(x)+f'(x)],$ \\

由$Rolle$定理,$\;\exists\,\xi_1\in(a,x_0),\xi_2\in(x_0,b),\;\text{s.t.}\,G'(\xi_1)=G'(\xi_2)=0,$ \\

令$F(x)=e^{-x}[f(x)+f'(x)],\text{则}F(\xi_1)=F(\xi_2)=0,$ \\

$\because f\text{在}[a,b]\text{上连续},\text{在}(a,b)\text{上二阶可导},\quad\therefore F\text{在}[\xi_1,\xi_2]\text{上连续},\text{在}(\xi_1,\xi_2)\text{上可导},$ \\

由$Rolle$定理,$\;\exists\,\zeta\in(\xi_1,\xi_2),\;\text{s.t.}\,F'(\zeta)=0$, \\

又$F'(x)=e^{-x}[f''(x)-f(x)],\quad\therefore f''(\zeta)=f(\zeta)$, \\

综上,$\exists\,\zeta\in[a,b],\;\text{s.t.}f(\zeta)=f''(\zeta).$ 

\begin{ti}
【题6】设$f\text{在}[a,b]\text{上可微,在}(a,b)\text{上二阶可微},$ 证明:$\exists\,\xi\in(a,b),\;\text{s.t.}\,f'(b)-f'(a)=f''(\xi)(b-a)$
\end{ti}

\textbf{证明}:假设不存在这样的$\xi$,由$Darboux$定理, \\

$f''(x)>\dfrac{f'(b)-f'(a)}{b-a}\text{恒成立或}f''(x)<\dfrac{f'(b)-f'(a)}{b-a}\text{恒成立}$, \\

不妨设$f''(x)>\dfrac{f'(b)-f'(a)}{b-a},\;\forall\,x\in(a,b)$, \\

设$\lambda=\dfrac{f'(b)-f'(a)}{b-a},\;g(x)=f'(x)-\lambda x,\text{则}x\in(a,b)\text{时},g'(x)=f''(x)-\lambda>0$, \\

$\therefore\,g(x)\text{在}(a,b)$上严格单调递增,且$g(a)=g(b)$, \\

$\therefore\,g(a^+),g(b^-)$均存在,即$\displaystyle\lim_{x\to a^+}f'(x),\lim_{x\to b^-}f'(x)$均存在, \\

下面证明一个重要引理:若$\displaystyle\lim_{x\to a^+}f'(x)$存在,则必然等于$f'(a)$, \\

事实上,由$Lagrange$中值定理,$\exists\,\xi_x\in(a,x),\;\text{s.t.}\,\dfrac{f(x)-f(a)}{x-a}=f'(\xi_x)$, \\

$x\rightarrow a^+\text{时},\xi_x\rightarrow a^+$,对上式两边取极限即得到上述引理. \\

由上面的引理,可知$g$在$a,b$处均连续,即$g\in C[a,b]$, \\

$\therefore g\text{在}[a,b]$上严格单调递增,这与$g(a)=g(b)$矛盾! \\

综上,$\exists\,\xi\in(a,b),\;\text{s.t.}\,f'(b)-f'(a)=f''(\xi)(b-a)$. \\

\textbf{注意}:本题并未要求$f'\in C[a,b]$,本题的意义在于,对于导函数来说,$Lagrange$定理中在闭区间上连续的要求不是必要的.

\begin{ti}
【题7】设$f\in C^1[a,b],\int_{a}^{b}f(x)\,dx=\int_{a}^{b}xf(x)\,dx=0,$求证:$\exists\,x_1,x_2\in(a,b),x_1\neq x_2,\;\text{s.t.}\,f(x_1)=f(x_2)=0$
\end{ti}

\textbf{证明}:若$f(x)\text{在}(a,b)$上无零点,由介值定理,$f(x)>0\text{或}f(x)<0$恒成立, \\

这与$\displaystyle\int_{a}^{b}f(x)\,dx=0$矛盾,$\quad\therefore\exists\,c\in(a,b),\;\text{s.t.}\,f(c)=0$, \\

若$f(x)$在$(a,b)$中只有这一个零点, \\

如果$f(x)$在$c$两边同号,则与$\displaystyle\int_{a}^{b}f(x)\,dx=0$矛盾, \\

$\therefore\,f(x)\text{在}c$两边异号, \\

不妨设$x\in(a,c)\text{时},f(x)<0,\;x\in(c,b)\text{时},f(x)>0$, \\

则$(x-c)f(x)>0$对$\forall\,x\in(a,b)$恒成立,$\;\therefore\displaystyle\int_{a}^{b}(x-c)f(x)\,dx>0$, \\

而$\displaystyle\int_{a}^{b}(x-c)f(x)\,dx=\int_{a}^{b}xf(x)\,dx-c\int_{a}^{b}f(x)\,dx=0$,矛盾! \\

所以$f$至少两个零点,即:$\exists\,x_1,x_2\in(a,b),x_1\neq x_2,\;\text{s.t.}\,f(x_1)=f(x_2)=0$.

\begin{ti}
【题8】设$f$在$[a,b]$上二阶可微,$f'(a)=f'(b)=0$,证明:$\exists\,\xi\in(a,b),\;\text{s.t.}\,\left\lvert f''(\xi)\right\rvert\geq\dfrac{4}{(b-a)^2}\left\lvert f(b)-f(a)\right\rvert $
\end{ti}

\textbf{证明}:将$f(x)$分别在$a$处和$b$处展开, \\

$f(x)=f(a)+f'(a)(x-a)+\dfrac{f''(\xi_1)}{2}(x-a)^2=f(a)+\dfrac{f''(\xi_1)}{2}(x-a)^2$, \\

$f(x)=f(b)+f'(b)(x-b)+\dfrac{f''(\xi_2)}{2}(x-b)^2=f(b)+\dfrac{f''(\xi_2)}{2}(x-b)^2$, \\

代入$x=\dfrac{a+b}{2},f(\dfrac{a+b}{2})=f(a)+\dfrac{f''(\xi_1)}{2}\cdot\dfrac{(b-a)^2}{4}=f(b)+\dfrac{f''(\xi_2)}{2}\cdot\dfrac{(b-a)^2}{4}$, \\

则$\left\lvert f(b)-f(a)\right\rvert\cdot\dfrac{4}{(b-a)^2}=\dfrac{1}{2}\left\lvert f''(\xi_2)-f''(\xi_1)\right\rvert$ $\leq\dfrac{1}{2}\left(\left\lvert f''(\xi_1)\right\rvert+\left\lvert f''(\xi_2)\right\rvert\right)$, \\

取$\xi\in\{\xi_1,\xi_2\},\text{使得}\left\lvert f''(\xi)\right\rvert=\max\{\left\lvert f''(\xi_1)\right\rvert,\left\lvert f''(\xi_2) \right\rvert\}$, \\

则$\left\lvert f''(\xi)\right\rvert\geq\dfrac{4}{(b-a)^2}\left\lvert f(b)-f(a)\right\rvert$,得证.

\begin{ti}
【题9】设$f$在$(a,b)$上可微,且在某点$\xi\in(a,b)$处有$f''(\xi)>0$,证明:$\exists\,x_1,x_2\in(a,b),\;\text{s.t.}\,\dfrac{f(x_2)-f(x_1)}{x_2-x_1}=f'(\xi)$
\end{ti}

\textbf{证明}:令$F(x)=f(x)-f'(\xi)x,\;\text{则}F'(\xi)=f'(\xi)-f'(\xi)=0,F''(\xi)=f''(\xi)>0,$ \\

由二阶导数的定义,$\therefore\exists\,\delta,\;\text{s.t.}\,x\in[\xi-\delta,\xi)\text{时},F'(x)<0;\;x\in(\xi,\xi+\delta]\text{时},F'(x)>0,$ \\

$\therefore\,x\in[\xi-\delta,\xi)\text{时},F(x)\text{严格单调递减},x\in(\xi,\xi+\delta]\text{时},F(x)\text{严格单调递增},$ \\

设$m=\min\{F(\xi-\delta),F(\xi+\delta)\},\text{则}F(\xi)<m,$ \\

取$\eta\in(F(\xi),m),\text{则}\exists\,x_1\in(\xi-\delta,\xi),x_2\in(\xi,\xi+\delta),\;\text{s.t.}\,F(x_1)=F(x_2)=\eta,$ \\

$\therefore\,0=\dfrac{F(x_2)-F(x_1)}{x_2-x_1}=\dfrac{f(x_2)-f(x_1)}{x_2-x_1}-f'(\xi),$ \\

$\therefore\,\dfrac{f(x_2)-f(x_1)}{x_2-x_1}=f'(\xi)$,这样的$x_1,x_2$即为所求. \\

\textbf{注}:本题虽然简单,但是可以与$Lagrange$中值定理对比来看,事实上,如果没有其他条件,那么任意的$\xi\in(a,b)$,不一定存在$x_1,x_2$,使得$\dfrac{f(x_2)-f(x_1)}{x_2-x_1}=f'(\xi)$,
读者可以自行思考存在$x_1,x_2$使得$\dfrac{f(x_2)-f(x_1)}{x_2-x_1}=f'(\xi)$的充要条件是什么.

\begin{ti}
【题10】设$f$在$[a,b]$上可微,$f'(a)=f'(b)$,证明:$\exists\,\xi\in(a,b),\;\text{s.t.}\,f'(\xi)=\dfrac{f(\xi)-f(a)}{\xi-a}$
\end{ti}

\textbf{证明}:令$F(x)=f(x)-f(a)-f'(a)(x-a),\text{则}F(a)=0,F'(a)=F'(b)=0,$ \\

\begin{minipage}[t]{0.58\textwidth}
考虑本题的几何含义, \\

定义$H(x)=\begin{cases}
    \dfrac{F(x)}{x-a}&\quad (a<x\leq b) \\
    0&\quad (x=a) \\
\end{cases}$ \\[3pt]

易得$H(a^+)=H(a)$, \\

$\therefore H\text{在}[a,b]\text{上连续}$, $H\text{在}(a,b)\text{上可导}$, \\
\end{minipage}\hfill
\begin{minipage}[t]{0.38\textwidth}
\centering
\resizebox{0.96\linewidth}{!}{%
\begin{tikzpicture}[x=0.90cm,y=1.05cm,baseline=(current bounding box.north)]
    \draw[bookteal!80,line width=0.9pt]
        (0.35,0.45) .. controls (0.85,1.55) and (1.45,2.25) .. (2.15,1.75)
        .. controls (2.75,1.30) and (3.80,1.5253) .. (4.20,1.65)
        .. controls (4.60,1.7747) and (5.15,1.82) .. (5.70,1.55);
    \draw[bookgray!85,line width=0.8pt]
        (-0.10,0.3097) -- (5.70,2.1175);
    \fill[bookblue] (0.35,0.45) circle (1.5pt);
    \fill[bookblue] (4.20,1.65) circle (1.5pt);
    \node[below left,font=\scriptsize] at (0.35,0.45) {$(a,f(a))$};
    \node[below right,font=\scriptsize] at (4.20,1.65) {$(x,f(x))$};
\end{tikzpicture}%
}

\vspace{2pt}
{\centering\small\color{bookgray}图4-2\par}
\end{minipage}
\par


且$x\in(a,b]\text{时},\;H'(x)=\dfrac{(x-a)f'(x)-f(x)+f(a)}{(x-a)^2},$ \\

$\therefore\,H(b)=\dfrac{F(b)}{b-a},\;H'(b)=\dfrac{(b-a)f'(b)-f(b)+f(a)}{(b-a)^2}=-\dfrac{F(b)}{(b-a)^2}=-\dfrac{H(b)}{b-a}.$ \\

\romannum{1} $\quad H(b)>0,\text{则}H'(b)<0,\text{则}\exists\,c\in(a,b),\;\text{s.t.}\,H(c)>H(b),$ \\

又$H(c)>H(b)>0=H(a),\;\therefore\,H$在区间内部取到最大值,\\

即$\exists\,\xi\in(a,b),\;\text{s.t.}\,H'(\xi)=0.$ \\

\romannum{2} $\quad H(b)<0,\text{则}H'(b)>0,\text{则}\exists\,c\in(a,b),\;\text{s.t.}\,H(c)<H(b),$ \\

又$H(c)<H(b)<0=H(a),\;\therefore\,H$在区间内部取到最小值, \\

即$\exists\,\xi\in(a,b),\;\text{s.t.}\,H'(\xi)=0.$ \\

\romannum{3} $\quad H(b)=0$,则由$Rolle$定理,$\exists\,\xi\in(a,b),\;\text{s.t.}\,H'(\xi)=0.$ \\

综上,$\exists\,\xi\in(a,b),\;\text{s.t.}\,H'(\xi)=0$,即$\exists\,\xi\in(a,b),\;\text{s.t.}\,f'(\xi)=\dfrac{f(\xi)-f(a)}{\xi-a}.$ 

\begin{ti}
【题11】设$f$在$[0,+\infty)$上可微,$f(0)=0,\text{且}\exists\,c>0,\;\text{s.t.}\,\left\lvert f'(x)\right\rvert\leq c\left\lvert f(x)\right\rvert,\forall\,x\in[0,+\infty),\;$
证明:$f(x)\equiv 0$
\end{ti}

\textbf{证明}:易得$\left\lvert f\right\rvert\in C[0,\dfrac{1}{2c}]$,由连续函数最值定理,设$\left\lvert f(x_0)\right\rvert=\displaystyle\max_{x\in[0,\frac{1}{2c}]}\left\lvert f(x)\right\rvert$, \\

$\therefore\left\lvert f(x_0)\right\rvert=\left\lvert f(x_0)-f(0)\right\rvert=\left\lvert f'(\xi)\right\rvert x_0\leq c\left\lvert f(\xi)\right\rvert x_0$ \\

$\leq\dfrac{1}{2}\left\lvert f(\xi)\right\rvert\leq\dfrac{1}{2}\left\lvert f(x_0)\right\rvert\Longrightarrow \left\lvert f(x_0)\right\rvert=0,$ $\;\therefore\,f(x)\equiv 0,\;\forall\,x\in[0,\dfrac{1}{2c}]$, \\

假设$x\in[\dfrac{k-1}{2c},\dfrac{k}{2c}]\text{时},f(x)\equiv 0,\text{记}\left\lvert f(x_1)\right\rvert=\displaystyle\max_{x\in[\frac{k}{2c},\frac{k+1}{2c}]}\left\lvert f(x)\right\rvert,$ \\

$\therefore\left\lvert f(x_1)\right\rvert=\left\lvert f(x_1)-f(\dfrac{k}{2c})\right\rvert=\left\lvert f'(\xi_1)\right\rvert(x_1-\dfrac{k}{2c})\leq c\left\lvert f(\xi_1)\right\rvert\cdot\dfrac{1}{2c}$ \\

$\leq\dfrac{1}{2}\left\lvert f(x_1)\right\rvert\Longrightarrow\left\lvert f(x_1)\right\rvert=0,\;$ 即$f(x)\equiv 0,\;\forall\,x\in[\dfrac{k}{2c},\dfrac{k+1}{2c}],$ \\

由数学归纳法,$f(x)\equiv 0,\;\forall\,x\in[\dfrac{n}{2c},\dfrac{n+1}{2c}],\;\text{对}\forall\,n\in\mathbb{N}\text{恒成立}$, \\

即$f(x)\equiv 0,\;\forall\,x\in\displaystyle\bigcup_{n=0}^{\infty}[\dfrac{n}{2c},\dfrac{n+1}{2c}]$,\\

即$f(x)\equiv 0,\;\forall\,x\in[0,+\infty)$,证毕. \\

\textbf{另解}:由题意可得$f(x)=\displaystyle\int_{0}^{x}f'(t)\,dt,$ \\

$\displaystyle\therefore\left\lvert f(x)\right\rvert=\left\lvert \int_{0}^{x}f'(t)\,dt\right\rvert\leq\int_{0}^{x}\left\lvert f'(t)\right\rvert \,dt\leq c\int_{0}^{x}\left\lvert f(t)\right\rvert\,dt $, \\

$\forall\,x_0\in(0,+\infty),\;f\in C[0,x_0],\text{记}M=\displaystyle\max_{t\in[0,x_0]}\left\lvert f(t)\right\rvert,$ \\

则$\displaystyle\left\lvert f(x)\right\rvert\leq c\int_{0}^{x}\left\lvert f(t)\right\rvert \,dt\leq c\cdot Mx,\;\forall\,x\in[0,x_0],$ \\

再将此式带回,可得$\displaystyle\left\lvert f(x)\right\rvert\leq c\int_{0}^{x}\left\lvert f(t)\right\rvert \,dt\leq c^2\int_{0}^{x}Mt\,dt=\dfrac{Mc^2}{2}x^2,$ \\

以此归纳地做下去,可得$x\in[0,x_0]\text{时},\left\lvert f(x)\right\rvert\leq\dfrac{Mc^nx^n}{n!},\text{对}\forall\,n\in\mathbb{N}^+$恒成立, \\

$\therefore\,x\in[0,x_0]\text{时},\left\lvert f(x)\right\rvert\leq\displaystyle\lim_{n\to\infty}\dfrac{Mc^nx^n}{n!}=0,\text{即}f(x)=0,$ \\

由$x_0$的任意性可知,$f\equiv 0$,证毕.



\BookEnd

\ifdefined\BookMaster
\def\BookEnd{}
\else
\documentclass[a4paper,12pt,fleqn,openany,twoside,fontset=windows]{ctexbook}

\usepackage[fleqn]{amsmath}
\usepackage{amssymb,amsthm,mathrsfs}
\usepackage{geometry}
\usepackage[dvipsnames]{xcolor}
\usepackage{graphicx}
\usepackage{tikz}
\usetikzlibrary{calc}
\usepackage[most]{tcolorbox}
\usepackage{fancyhdr}
\usepackage{titlesec}
\usepackage{tocloft}
\usepackage{enumitem}
\usepackage{microtype}
\usepackage{pifont}
\usepackage{hyperref}
\usepackage{romannum}
\usepackage{mathtools}

\definecolor{bookblue}{HTML}{264653}
\definecolor{bookteal}{HTML}{4F7C82}
\definecolor{booklight}{HTML}{EDF4F3}
\definecolor{bookgray}{HTML}{5E686B}

\geometry{
    inner=2.7cm,
    outer=2.5cm,
    top=2.7cm,
    bottom=2.7cm,
    headheight=18pt,
    headsep=18pt,
    footskip=25pt
}
\setlength{\mathindent}{0pt}
\setlength{\parindent}{5pt}
\setlength{\parskip}{3pt plus 1pt minus 1pt}
\linespread{1.25}
\allowdisplaybreaks[3]
\AtBeginDocument{%
    \setlength{\parindent}{5pt}
    \setlength{\abovedisplayskip}{8pt plus 2pt minus 2pt}
    \setlength{\belowdisplayskip}{8pt plus 2pt minus 2pt}
    \setlength{\abovedisplayshortskip}{5pt plus 1pt minus 1pt}
    \setlength{\belowdisplayshortskip}{5pt plus 1pt minus 1pt}
}

\hypersetup{
    unicode=true,
    colorlinks=true,
    linkcolor=bookblue,
    citecolor=bookblue,
    urlcolor=bookteal,
    bookmarksopen=true,
    pdfauthor={Chen},
    pdftitle={数学分析习题整理}
}

\renewcommand{\chaptermark}[1]{%
    \edef\@tmp{第\thechapter 章\quad #1}%
    \expandafter\markboth\expandafter{\@tmp}{}%
}
\fancypagestyle{main}{%
    \fancyhf{}
    \fancyhead[LE]{\small\color{bookgray}\nouppercase{\leftmark}}
    \fancyhead[RO]{\small\color{bookgray}数学分析习题整理}
    \fancyfoot[C]{\color{bookblue}\thepage}
    \fancyfoot[R]{\small\bfseries Chen\;\; github.com/mathlover-1/math-analysis}
    \renewcommand{\headrulewidth}{0.4pt}
    \renewcommand{\footrulewidth}{0pt}
}
\fancypagestyle{plain}{%
    \fancyhf{}
    \fancyfoot[C]{\color{bookblue}\thepage}
    \fancyfoot[R]{\small\bfseries Chen\;\; github.com/mathlover-1/math-analysis}
    \renewcommand{\headrulewidth}{0pt}
    \renewcommand{\footrulewidth}{0pt}
}
\pagestyle{main}

\titleformat{\chapter}[display]
  {\normalfont\sffamily\bfseries\color{bookblue}}
  {\raggedleft\fontsize{46}{50}\selectfont\thechapter}
  {-18pt}
  {\titlerule[1.2pt]\vspace{10pt}\filright\Huge}
\titlespacing*{\chapter}{0pt}{-15pt}{36pt}

\renewcommand{\contentsname}{目\quad 录}
\renewcommand{\cfttoctitlefont}{\hfill\sffamily\bfseries\Huge\color{bookblue}}
\renewcommand{\cftaftertoctitle}{\hfill}
\setlength{\cftbeforetoctitleskip}{5pt}
\setlength{\cftaftertoctitleskip}{28pt}
\setlength{\cftbeforechapskip}{12pt}
\renewcommand{\cftchapfont}{\sffamily\bfseries\large\color{bookblue}}
\renewcommand{\cftchappagefont}{\sffamily\bfseries\large\color{bookblue}}
\renewcommand{\cftchapleader}{\color{bookteal}\cftdotfill{2.5}}

\newtcolorbox{ti}{
    enhanced,
    breakable,
    colback=booklight!55!white,
    colframe=bookteal!35!white,
    boxrule=0.45pt,
    boxsep=0pt,
    left=10pt,
    right=10pt,
    top=9pt,
    bottom=9pt,
    before skip=14pt,
    after skip=14pt,
    arc=2pt,
    outer arc=2pt,
    before upper={\parindent=0pt}
}

\newenvironment{booknotebody}{%
    \begin{center}
    \begin{minipage}{0.84\textwidth}
    \songti\zihao{-4}\linespread{1.45}\selectfont
    \setlength{\parindent}{2\ccwd}
    \setlength{\parskip}{0.75em}
}{%
    \end{minipage}
    \end{center}
}

\newenvironment{booknotelongbody}{%
    \begingroup
    \songti\zihao{-4}\linespread{1.45}\selectfont
    \setlength{\parindent}{2\ccwd}
    \setlength{\parskip}{0.75em}
    \begin{list}{}{%
        \setlength{\leftmargin}{0.08\textwidth}
        \setlength{\rightmargin}{0.08\textwidth}
        \setlength{\topsep}{0pt}
        \setlength{\partopsep}{0pt}
        \setlength{\parsep}{0pt}
        \setlength{\itemsep}{0pt}
    }
    \item\relax
}{%
    \end{list}
    \endgroup
}

\fancypagestyle{plainnowm}{%
    \fancyhf{}
    \fancyfoot[C]{\color{bookblue}\thepage}
    \renewcommand{\headrulewidth}{0pt}
    \renewcommand{\footrulewidth}{0pt}
}

\newcommand{\booknotetitle}[2][plain]{%
    \clearpage
    \phantomsection
    \addcontentsline{toc}{chapter}{#2}
    \markboth{#2}{}
    \thispagestyle{#1}
    \vspace*{1.2cm}
    \begin{center}
        {\sffamily\heiti\bfseries\fontsize{26}{34}\selectfont\color{bookblue}#2\par}
        \vspace{0.8em}
        {\color{bookteal}\rule{4.5cm}{0.8pt}\par}
    \end{center}
    \vspace{0.6cm}
}

\renewcommand{\proofname}{证明}
\renewcommand{\qedsymbol}{\color{bookteal}\ensuremath{\blacksquare}}

\title{数学分析习题整理}
\author{Chen}
\date{}

\makeatletter
% ==================== 旧版封面/封底设计备份 ====================
% 2026-08 重设计,旧代码用 \iffalse...\fi 封存,新版见下方。
\iffalse
\renewcommand{\maketitle}{%
    \begin{titlepage}
        \thispagestyle{empty}
        \setcounter{page}{0}
        \noindent\hspace*{-\parindent}\begin{tikzpicture}[remember picture,overlay]
            \fill[bookblue] (current page.north west) rectangle ([yshift=-4.2cm]current page.north east);
            \fill[booklight] ([yshift=-4.2cm]current page.north west) rectangle (current page.south east);
            \draw[bookteal!55,line width=0.7pt]
                ([xshift=2.4cm,yshift=4.3cm]current page.south west)
                .. controls +(2.5cm,1.8cm) and +(-2.8cm,-1.3cm) ..
                ([xshift=-2.4cm,yshift=7.2cm]current page.south east);
            \draw[bookteal!35,line width=0.7pt]
                ([xshift=2.5cm,yshift=6.2cm]current page.south west)
                .. controls +(3.0cm,-1.2cm) and +(-3.0cm,1.6cm) ..
                ([xshift=-2.3cm,yshift=3.8cm]current page.south east);
        \end{tikzpicture}
        \vspace*{5.6cm}
        \noindent\color{bookblue}\rule{\textwidth}{1.4pt}\par
        \vspace{8pt}
        {\centering\sffamily\bfseries\fontsize{32}{42}\selectfont\color{bookblue}\@title\par}
        \vspace{8pt}
        \noindent\color{bookteal}\rule{\textwidth}{0.6pt}\par
        \vspace{2.6cm}
        {\centering\Large\color{bookgray}\@author\par}
        \vfill
        {\centering\color{bookteal}\large
            $\displaystyle \lim_{n\to\infty} a_n$\qquad
            $\displaystyle \int_a^b f(x)\,\mathrm{d}x$\qquad
            $\displaystyle \nabla f(x,y)$\par
        }
        \vspace*{1.5cm}
    \end{titlepage}
}

\newcommand{\makebackcover}{%
    \begingroup
    \let\cleardoublepage\clearpage
    \begin{titlepage}
        \thispagestyle{empty}
        \noindent\hspace*{-\parindent}\begin{tikzpicture}[remember picture,overlay]
            \fill[bookblue] (current page.north west) rectangle (current page.south east);
            \fill[booklight,rounded corners=7pt]
                ([xshift=2.2cm,yshift=-3.0cm]current page.north west)
                rectangle
                ([xshift=-2.2cm,yshift=3.0cm]current page.south east);
            \fill[bookteal!80]
                (current page.north west) --
                ([xshift=5.3cm]current page.north west) --
                ([xshift=2.1cm,yshift=-4.4cm]current page.north west) --
                ([yshift=-6.2cm]current page.north west) -- cycle;
            \fill[bookteal!55]
                (current page.south east) --
                ([xshift=-5.8cm]current page.south east) --
                ([xshift=-2.4cm,yshift=4.2cm]current page.south east) --
                ([yshift=6.0cm]current page.south east) -- cycle;
            \foreach \radius in {0.75,1.15,1.55,1.95} {
                \draw[bookteal!55,line width=0.55pt]
                    ([xshift=-2.9cm,yshift=-4.4cm]current page.north east)
                    circle[radius=\radius cm];
            }
            \foreach \radius in {0.65,1.05,1.45} {
                \draw[booklight!65,line width=0.55pt]
                    ([xshift=2.4cm,yshift=3.6cm]current page.south west)
                    circle[radius=\radius cm];
            }
            \fill[bookteal!16,rounded corners=3pt]
                ([xshift=3.5cm,yshift=-7.1cm]current page.north west)
                rectangle
                ([xshift=-5.0cm,yshift=7.2cm]current page.south east);
            \fill[bookteal!30,rounded corners=3pt]
                ([xshift=4.4cm,yshift=-8.0cm]current page.north west)
                rectangle
                ([xshift=-4.1cm,yshift=8.1cm]current page.south east);
            \begin{scope}[shift={([xshift=-0.75cm]current page.center)}]
                \clip[rounded corners=3pt] (-5.25,-4.75) rectangle (5.25,4.75);
                \foreach \v in {-3.6,-3.0,...,3.6} {
                    \draw[bookteal!48,line width=0.45pt,opacity=0.72]
                        plot[domain=-3.6:3.6,samples=45,smooth,variable=\u]
                        ({\u+0.55*\v},{0.28*\u-0.35*\v+0.11*(\u*\u-\v*\v)});
                }
                \foreach \u in {-3.6,-3.0,...,3.6} {
                    \draw[bookblue!42,line width=0.42pt,opacity=0.62]
                        plot[domain=-3.6:3.6,samples=45,smooth,variable=\v]
                        ({\u+0.55*\v},{0.28*\u-0.35*\v+0.11*(\u*\u-\v*\v)});
                }
                \draw[bookteal!70,line width=0.75pt,opacity=0.82]
                    plot[domain=-3.6:3.6,samples=55,smooth,variable=\u]
                    ({\u},{0.28*\u+0.11*\u*\u});
                \draw[bookblue!60,line width=0.7pt,opacity=0.75]
                    plot[domain=-3.6:3.6,samples=55,smooth,variable=\v]
                    ({0.55*\v},{-0.35*\v-0.11*\v*\v});
            \end{scope}
            \draw[bookteal!45,line width=0.65pt]
                ([xshift=3.0cm,yshift=-8.2cm]current page.north west)
                .. controls +(2.2cm,2.0cm) and +(-2.0cm,-2.1cm) ..
                ([xshift=-3.0cm,yshift=9.0cm]current page.south east);
            \draw[bookteal!25,line width=0.55pt]
                ([xshift=3.1cm,yshift=8.0cm]current page.south west)
                .. controls +(2.7cm,-1.7cm) and +(-2.6cm,1.8cm) ..
                ([xshift=-3.1cm,yshift=-8.7cm]current page.north east);
            \node[anchor=west,text=bookteal!72,scale=0.78] at
                ([xshift=2.8cm,yshift=-4.0cm]current page.north west) {%
                $\displaystyle e^{\mathrm{i}\pi}+1=0$};
            \node[anchor=east,text=bookblue!62,scale=0.67] at
                ([xshift=-2.7cm,yshift=4.0cm]current page.south east) {%
                $\displaystyle
                f(x)=\sum_{n=0}^{\infty}\frac{f^{(n)}(a)}{n!}(x-a)^n$};
        \end{tikzpicture}
        \mbox{}
    \end{titlepage}
    \endgroup
}
\fi

% ==================== 新版封面(2026-08 重设计) ====================
% 设计要点:顶部深绿纵向渐变+波浪下缘+两层半透明过渡波(解决深浅绿生硬过渡);
% 深色区螺旋线/同心圆母题;标题+英文副题+作者居中;
% 中部 Fourier 部分和逼近方波曲线;底部浅色波纹(无文字)。
\renewcommand{\maketitle}{%
    \begin{titlepage}
        \thispagestyle{empty}
        \setcounter{page}{0}
        \noindent\hspace*{-\parindent}\begin{tikzpicture}[remember picture,overlay]
            % ---------- 底色 ----------
            \fill[booklight] (current page.north west) rectangle (current page.south east);

            % ---------- 顶部深色区:纵向渐变 + 波浪下缘 ----------
            \shade[top color=bookblue, bottom color=bookteal!55!bookblue]
                (current page.north west) -- (current page.north east) --
                ([yshift=-8.6cm]current page.north east)
                .. controls ([xshift=-5.2cm,yshift=-10.6cm]current page.north east)
                         and ([xshift=6.0cm,yshift=-8.0cm]current page.north west) ..
                ([yshift=-10.0cm]current page.north west) -- cycle;

            % ---------- 过渡层 1:半透明波浪 ----------
            \fill[bookteal!60!bookblue,opacity=0.42]
                ([yshift=-8.6cm]current page.north east)
                .. controls ([xshift=-5.2cm,yshift=-10.6cm]current page.north east)
                         and ([xshift=6.0cm,yshift=-8.0cm]current page.north west) ..
                ([yshift=-10.0cm]current page.north west) --
                ([yshift=-12.6cm]current page.north west)
                .. controls ([xshift=6.5cm,yshift=-11.0cm]current page.north west)
                         and ([xshift=-5.5cm,yshift=-13.2cm]current page.north east) ..
                ([yshift=-11.4cm]current page.north east) -- cycle;

            % ---------- 过渡层 2:更浅的波浪 ----------
            \fill[bookteal!35,opacity=0.35]
                ([yshift=-10.6cm]current page.north east)
                .. controls ([xshift=-6.0cm,yshift=-12.4cm]current page.north east)
                         and ([xshift=5.0cm,yshift=-9.8cm]current page.north west) ..
                ([yshift=-11.8cm]current page.north west) --
                ([yshift=-13.8cm]current page.north west)
                .. controls ([xshift=5.5cm,yshift=-12.6cm]current page.north west)
                         and ([xshift=-6.5cm,yshift=-14.6cm]current page.north east) ..
                ([yshift=-12.8cm]current page.north east) -- cycle;

            % ---------- 深色区装饰:右上螺旋线 ----------
            \begin{scope}[shift={([xshift=-3.4cm,yshift=-4.3cm]current page.north east)}]
                \draw[booklight!45!bookteal,opacity=0.55,line width=0.6pt]
                    plot[domain=0:1440,samples=220,smooth,variable=\t]
                    ({0.0028*\t*cos(\t)},{0.0028*\t*sin(\t)});
            \end{scope}

            % ---------- 深色区装饰:左下同心圆 ----------
            \foreach \r in {0.7,1.25,1.8,2.35,2.9}{
                \draw[booklight!35!bookteal,opacity=0.5,line width=0.5pt]
                    ([xshift=0.4cm,yshift=-7.9cm]current page.north west) circle[radius=\r cm];
            }

            % ---------- 主标题 ----------
            \node[anchor=center,text=booklight] at
                ([yshift=-4.5cm]current page.north)
                {\sffamily\bfseries\fontsize{40}{46}\selectfont \@title};

            % ---------- 分隔短线 ----------
            \draw[booklight!65,line width=1.2pt]
                ([xshift=-2.6cm,yshift=-5.75cm]current page.north) --
                ([xshift=2.6cm,yshift=-5.75cm]current page.north);

            % ---------- 英文副题 ----------
            \node[anchor=center,text=booklight!70!bookteal] at
                ([yshift=-6.55cm]current page.north)
                {\itshape\large Mathematical Analysis \textperiodcentered\
                 A Collection of Exercises};

            % ---------- 作者(紧随标题下方) ----------
            \node[anchor=center,text=booklight!88] at
                ([yshift=-8.2cm]current page.north)
                {\sffamily\large \@author\quad 整理};

            % ---------- 中部:Fourier 部分和逼近方波 ----------
            \begin{scope}[shift={([xshift=-5.2cm,yshift=-18.1cm]current page.north)}]
                % 坐标轴
                \draw[bookgray!55,line width=0.5pt,->] (-0.5,0) -- (10.9,0);
                \draw[bookgray!55,line width=0.5pt,->] (0,-1.75) -- (0,1.95);
                % 目标方波(虚线)
                \draw[bookgray!75,line width=0.6pt,dashed]
                    (0,0.92) -- (5,0.92) (5,-0.92) -- (10,-0.92);
                % N=1
                \draw[bookteal!45,line width=0.7pt]
                    plot[domain=0:360,samples=90,smooth,variable=\x]
                    ({\x/36},{1.18*sin(\x)});
                % N=3
                \draw[bookteal!80,line width=0.8pt]
                    plot[domain=0:360,samples=120,smooth,variable=\x]
                    ({\x/36},{1.18*(sin(\x)+sin(3*\x)/3)});
                % N=7
                \draw[bookblue,line width=0.95pt]
                    plot[domain=0:360,samples=150,smooth,variable=\x]
                    ({\x/36},{1.18*(sin(\x)+sin(3*\x)/3+sin(5*\x)/5+sin(7*\x)/7)});
                % 刻度标注
                \node[anchor=north east,text=bookgray!80,scale=0.62] at (-0.12,0) {$0$};
                \node[anchor=north,text=bookgray!80,scale=0.62] at (5,0) {$\pi$};
                \node[anchor=north,text=bookgray!80,scale=0.62] at (10,0) {$2\pi$};
            \end{scope}
            % 图注
            \node[anchor=center,text=bookgray,scale=0.8] at
                ([yshift=-20.5cm]current page.north)
                {$S_N(x)=\displaystyle\sum_{k=1}^{N}\frac{\sin\bigl((2k-1)x\bigr)}{2k-1}$
                 \ 逐步逼近方波};

            % ---------- 底部浅色波纹(无文字,与顶部波浪呼应) ----------
            \fill[bookteal!22,opacity=0.5]
                ([yshift=2.4cm]current page.south west)
                .. controls ([xshift=6.0cm,yshift=3.9cm]current page.south west)
                         and ([xshift=-6.5cm,yshift=1.0cm]current page.south east) ..
                ([yshift=2.9cm]current page.south east) --
                (current page.south east) -- (current page.south west) -- cycle;
            \fill[bookteal!40!bookblue,opacity=0.28]
                ([yshift=1.3cm]current page.south west)
                .. controls ([xshift=5.5cm,yshift=2.7cm]current page.south west)
                         and ([xshift=-5.5cm,yshift=0.3cm]current page.south east) ..
                ([yshift=1.9cm]current page.south east) --
                (current page.south east) -- (current page.south west) -- cycle;
        \end{tikzpicture}
        \mbox{}
    \end{titlepage}
}

% ==================== 新版封底(纯背景,2026-08 重设计) ====================
% 设计要点:对角渐变底+角部半透明叠层;中下部外摆线(玫瑰形)为视觉主体;
% 底部两层半透明白色波浪;无任何文字。
\newcommand{\makebackcover}{%
    \begingroup
    \let\cleardoublepage\clearpage
    \begin{titlepage}
        \thispagestyle{empty}
        \noindent\hspace*{-\parindent}\begin{tikzpicture}[remember picture,overlay]
            % ---------- 底色:对角渐变 ----------
            \shade[top color=bookblue, bottom color=bookteal!48!bookblue,
                   shading angle=-30]
                (current page.north west) rectangle (current page.south east);

            % ---------- 角部半透明叠层 ----------
            \fill[white,opacity=0.05]
                ([xshift=2.0cm,yshift=1.0cm]current page.north east) circle[radius=7.2cm];
            \fill[white,opacity=0.06]
                ([xshift=-0.5cm,yshift=-1.5cm]current page.north east) circle[radius=3.6cm];

            % ---------- 左下同心圆 ----------
            \foreach \r in {1.0,1.75,2.5,3.25,4.0,4.75}{
                \draw[booklight!40,opacity=0.55,line width=0.55pt]
                    ([xshift=-0.6cm,yshift=-0.4cm]current page.south west) circle[radius=\r cm];
            }
            % ---------- 右下螺旋线 ----------
            \begin{scope}[shift={([xshift=-2.9cm,yshift=3.2cm]current page.south east)}]
                \draw[booklight!45!bookteal,opacity=0.5,line width=0.55pt]
                    plot[domain=0:1260,samples=200,smooth,variable=\t]
                    ({0.0036*\t*cos(\t)},{0.0036*\t*sin(\t)});
            \end{scope}
            % ---------- 左上斜线组 ----------
            \foreach \i in {0,1,2}{
                \draw[bookteal!60!booklight,opacity=0.5,line width=0.6pt]
                    ([xshift={1.6cm+\i*0.45cm}]current page.north west) --
                    ([yshift={-1.6cm-\i*0.45cm}]current page.north west);
            }

            % ---------- 中下部视觉主体:外摆线(玫瑰形) ----------
            \begin{scope}[shift={([yshift=-16.6cm]current page.north)}]
                % 外圈
                \draw[booklight!35,opacity=0.6,line width=0.6pt]
                    (0,0) circle[radius=4.9cm];
                \draw[booklight!22,opacity=0.5,line width=0.5pt]
                    (0,0) circle[radius=5.25cm];
                % 主玫瑰线
                \draw[booklight!55!bookteal,opacity=0.8,line width=0.85pt]
                    plot[domain=0:1080,samples=400,smooth,variable=\t]
                    ({0.335*(8*cos(\t)-5*cos(2.6667*\t))},
                     {0.335*(8*sin(\t)-5*sin(2.6667*\t))});
                % 内层玫瑰线
                \draw[booklight!38!bookteal,opacity=0.6,line width=0.6pt]
                    plot[domain=0:1080,samples=400,smooth,variable=\t]
                    ({0.21*(8*cos(\t)-5*cos(2.6667*\t))},
                     {0.21*(8*sin(\t)-5*sin(2.6667*\t))});
            \end{scope}

            % ---------- 底部半透明波浪 ----------
            \fill[white,opacity=0.06]
                ([yshift=4.6cm]current page.south west)
                .. controls ([xshift=6.5cm,yshift=6.4cm]current page.south west)
                         and ([xshift=-6.0cm,yshift=2.6cm]current page.south east) ..
                ([yshift=5.2cm]current page.south east) --
                (current page.south east) -- (current page.south west) -- cycle;
            \fill[booklight,opacity=0.08]
                ([yshift=2.6cm]current page.south west)
                .. controls ([xshift=5.5cm,yshift=4.3cm]current page.south west)
                         and ([xshift=-6.5cm,yshift=1.0cm]current page.south east) ..
                ([yshift=3.3cm]current page.south east) --
                (current page.south east) -- (current page.south west) -- cycle;
        \end{tikzpicture}
        \mbox{}
    \end{titlepage}
    \endgroup
}
\makeatother

% 页脚右下方常态化水印(署名 + 仓库地址,加粗黑色),已在 main/plain 分页样式中定义,所有页面均生效.

\begin{document}
\mainmatter
\pagestyle{main}
\setcounter{chapter}{4}
\def\BookEnd{\end{document}}
\fi

\chapter{多元微分学与多元积分学}
\begin{ti}
【题1】若$f(x,y)$是在$\mathbb{R}^2\backslash\{0\}$上连续的$n$次齐次函数(注:$n$次齐次函数是指满足$f(t\mathbf{x})=t^nf(\mathbf{x})$的函数),
证明:$\left\lvert f(x,y)\right\rvert\leq C(x^2+y^2)^{\frac{n}{2}}$,其中$C$为正常数
\end{ti}

\textbf{证明}:考虑集合$S=\left\{(x,y)\,|\;x^2+y^2=1\right\}$,则$S$为有界闭集, \\

又$f$为连续函数,$\therefore\left\lvert f\right\rvert $也是连续函数,$\therefore\exists\,C>0,\,(x,y)\in S\text{时},\left\lvert f(x,y)\right\rvert\leq C$, \\

对$\forall\,(x,y)\in\mathbb{R}^2\backslash\{0\},\text{令}r=\sqrt{x^2+y^2},$ 则$f(x,y)=f(r\cdot\dfrac{x}{r},r\cdot\dfrac{y}{r})=r^nf(\dfrac{x}{r},\dfrac{y}{r}),$ \\

又$\because(\dfrac{x}{r},\dfrac{y}{r})\in S,\quad\therefore\left\lvert f(\dfrac{x}{r},\dfrac{y}{r})\right\rvert\leq C$, \\

$\therefore\left\lvert f(x,y)\right\rvert\leq C(x^2+y^2)^{\frac{n}{2}}$.

\begin{ti}
【题2】如果$F(x,y)$在矩形$D=\left\{(x,y)\in\mathbb{R}^2\,|\;\left\lvert x-x_0\right\rvert<a,\left\lvert y-y_0\right\rvert<b\right\}$上连续,
$F(x_0,y_0)=0,$且$F(x,y)$对每一个$x\in(x_0-a,x_0+a)$关于$y$严格单调,求证:在$(x_0,y_0)$某个邻域内,由$F(x,y)=0$可确定唯一的函数$y=f(x)$,且$f(x)$在$x_0$的一个邻域里连续
\end{ti}

\textbf{证明}:不妨设关于$\displaystyle y$严格单调递增.取$\displaystyle(x_0,y_1),(x_0,y_2)\in D$,且$\displaystyle y_1>y_0>y_2$, \\

则$\displaystyle F(x_0,y_1)>0,F(x_0,y_2)<0$.取$\displaystyle0<\alpha<a$使$\displaystyle S=[x_0-\alpha,x_0+\alpha]\times[y_2,y_1]\subset D$, \\

因为$\displaystyle F$在闭矩形$\displaystyle S$上一致连续,令$\displaystyle\varepsilon_1=\dfrac{\min\{|F(x_0,y_1)|,|F(x_0,y_2)|\}}{2}$, \\

可取$\displaystyle\delta_0>0$,使$\displaystyle u,v\in S$且$\displaystyle\left\lVert u-v\right\rVert <\delta_0$时,$\displaystyle|F(u)-F(v)|<\varepsilon_1$, \\

令$\displaystyle\delta=\min\{\delta_0,\alpha\}$,对每个$\displaystyle x\in(x_0-\delta,x_0+\delta)$,有$\displaystyle F(x,y_1)>0,F(x,y_2)<0$, \\

故由连续性和严格单调性,存在唯一的$\displaystyle f(x)\in(y_2,y_1)$使$\displaystyle F(x,f(x))=0$,且$\displaystyle f(x_0)=y_0$, \\

下证$\displaystyle f$连续,任取$\displaystyle x^*\in(x_0-\delta,x_0+\delta)$及$\displaystyle\varepsilon>0$, \\

取$\displaystyle 0<\rho<\min\{\varepsilon,f(x^*)-y_2,y_1-f(x^*)\}$, \\

则由严格单调性,$\displaystyle F(x^*,f(x^*)-\rho)<0<F(x^*,f(x^*)+\rho)$, \\

由$\displaystyle F$的连续性,存在$\displaystyle\eta_0>0$,使$\displaystyle|x-x^*|<\eta_0$时上述两端仍分别为负和正, \\

令$\displaystyle\eta=\min\{\eta_0,\delta-|x^*-x_0|\}>0$,当$\displaystyle|x-x^*|<\eta$时, \\

唯一零点$\displaystyle f(x)$位于$\displaystyle(f(x_*)-\rho,f(x_*)+\rho)$内,故$\displaystyle|f(x)-f(x_*)|<\rho<\varepsilon$, \\

所以$\displaystyle f$在$\displaystyle(x_0-\delta,x_0+\delta)$上连续,命题得证.

\begin{ti}
【题3】$D=\left\{u\in\mathbb{R}^n\,|\;\left\lVert u\right\rVert\leq\frac{1}{2}\right\},$ $f\colon D\to\mathbb{R}\text{为}C^1$映射, \\
$\left\lVert \nabla f(0)\right\rVert=1,\text{且}\left\lVert \nabla f(u)-\nabla f(v)\right\rVert\leq\left\lVert u-v\right\rVert $, \\
证明:存在唯一的点$\xi$,使得$f(\xi)=\displaystyle\max_{D}f$
\end{ti}

\textbf{证明}:$\because D$为有界闭集,$f$为连续映射,从而$f$在$D$上能取到最大值, \\

当$u\in int D\text{时},\left\lVert \nabla f(u)-\nabla f(0)\right\rVert\leq\left\lVert u\right\rVert<\dfrac{1}{2}\Rightarrow\dfrac{1}{2}<\left\lVert \nabla f(u)\right\rVert<\dfrac{3}{2},$ \\

特别地,$\left\lVert \nabla f(u)\right\rVert\neq 0,$从而$f$的最大值只能在$\partial D$上取到, \\

令$F(u)=f(u)-\dfrac{1}{2}\left\lVert u\right\rVert^2,\;\therefore\,\nabla F(u)=\nabla f(u)-u$,同理$F$可在$D$上取到最大值, \\

当$u\in int D\text{时},\left\lVert \nabla F(u)\right\rVert=\left\lVert \nabla f(u)-u\right\rVert\geq\left\lVert \nabla f(u)\right\rVert-\left\lVert u\right\rVert>0$, \\

$\therefore\,\nabla F(u)\neq 0$,即$F$的最大值也只能在$\partial D$上取到, \\

又$\displaystyle\max_{D}F=\max_{\partial D}F=\max_{\partial D}f-\dfrac{1}{8}=\max_{D}f-\dfrac{1}{8}$, \\

$\therefore F\text{与}f$有完全相同的最大值点,因此只需要证明$F$的最大值点唯一即可, \\

对任意$\displaystyle u,v\in D$,由条件与$Cauchy$不等式, \\

$\displaystyle(\nabla F(u)-\nabla F(v))\cdot(u-v)$ \\

$\displaystyle=(\nabla f(u)-\nabla f(v))\cdot(u-v)-\left\lVert u-v\right\rVert^2$ \\

$\displaystyle\leq\left\lVert\nabla f(u)-\nabla f(v)\right\rVert\left\lVert u-v\right\rVert-\left\lVert u-v\right\rVert^2\leq0,$ \\

令$\displaystyle w(t)=(1-t)u+tv,\;\psi(t)=F(w(t))$,若$\displaystyle0\leq s<t\leq1$, \\

则$\displaystyle(\nabla F(w(t))-\nabla F(w(s)))\cdot(w(t)-w(s))\leq0,$ \\

而$\displaystyle w(t)-w(s)=(t-s)(v-u)$,故$\displaystyle\psi'(t)\leq\psi'(s)$, \\

因此$\displaystyle\psi'$单调不增,$\displaystyle\psi$为凹函数,从而$\displaystyle F$在$\displaystyle D$上为凹函数, \\

假设$\displaystyle u^0,v^0$是$\displaystyle F$的两个不同最大值点,记$\displaystyle F(u^0)=F(v^0)=M$, \\

对任意$\displaystyle t\in(0,1)$,由凹性与$\displaystyle M$的最大性, \\

$\displaystyle M\geq F((1-t)u^0+tv^0)\geq(1-t)F(u^0)+tF(v^0)=M,$ \\

故线段开段上每一点都是$\displaystyle F$的最大值点, \\

由于闭球$\displaystyle D$严格凸,$\displaystyle(1-t)u^0+tv^0\in\operatorname{int}D$, \\

由$Fermat$必要条件可知该点满足$\displaystyle\nabla F=0$,与前面已证的$\displaystyle\nabla F\neq0$矛盾, \\

从而$F$只能有唯一的最大值点,即存在唯一的点$\xi$,使得$f(\xi)=\displaystyle\max_{D}f$,得证.

\begin{ti}
【题4】二元函数$f(x,y)$,定义域为$I=[a,b]\times[c,d]$,$f$关于每个变量单调递增,证明:$f\in R(I)$
\end{ti}

\textbf{证明}:取等距分划$\displaystyle a=x_0<\cdots<x_n=b,\quad c=y_0<\cdots<y_n=d$,并令 \\

$\displaystyle R_{ij}=[x_{i-1},x_i]\times[y_{j-1},y_j],\quad\Delta x=\dfrac{b-a}{n},\quad\Delta y=\dfrac{d-c}{n}.$ \\

由单调性,$\displaystyle f$在$\displaystyle R_{ij}$上的振幅为$\displaystyle\omega_{ij}=f(x_i,y_j)-f(x_{i-1},y_{j-1})$ \\

$\displaystyle=[f(x_i,y_j)-f(x_{i-1},y_j)]+[f(x_{i-1},y_j)-f(x_{i-1},y_{j-1})]$, \\

因此沿$\displaystyle i,j$分别求和,可得$\displaystyle S(f,P_n)-s(f,P_n)=\sum_{i,j}\omega_{ij}\Delta x\Delta y$ \\

$\displaystyle=\Delta x\Delta y\left\{\sum_{j=1}^n[f(b,y_j)-f(a,y_j)]+\sum_{i=1}^n[f(x_{i-1},d)-f(x_{i-1},c)]\right\}$ \\

$\displaystyle\leq\frac{2(b-a)(d-c)}{n}[f(b,d)-f(a,c)]\to0.$ \\

故对任意$\displaystyle\varepsilon>0$,可取充分大的$\displaystyle n$,使$\displaystyle S(f,P_n)-s(f,P_n)<\varepsilon$,从而$\displaystyle f\in R(I)$,得证.

\begin{ti}
【题5】$\gamma(t):[a,b]\to \mathbb{R}^2,\gamma(t)=(x(t),y(t))$,且 $x(t)\in C^\alpha[a,b],y(t)\in C^\beta[a,b],0<\alpha,\beta<1, \alpha+\beta>1$, 
证明:$\gamma([a,b])$ 是 $\mathbb{R}^2$ 中零面积集.(注:$x(t)\in C^\alpha[a,b]$指的是$\displaystyle|x(t_1)-x(t_2)|\leq C_x\left\lvert t_1-t_2\right\rvert^{\alpha}$)
\end{ti}

\textbf{证明}:将 $[a,b]$ $n$等分,$a=t_0<t_1<\dots<t_n=b, \Delta t_i=\dfrac{b-a}{n}$, \\

 设$x(t)$在$[t_{i-1},t_i]$上最大值,最小值分别为$\alpha_i, \beta_i$, \\
 
 $\phantom{\text{设}}y(t)$在$[t_{i-1},t_i]$上最大值,最小值分别为$\mu_i, \nu _i$, \\
 
 令$I_i=[\beta_i,\alpha_i]\times[\nu_i,\mu_i]$ ,则 $\displaystyle\gamma([a,b])\subset\bigcup_{i=1}^n I_i$, 取$C=\max\{C_x,C_y\}$, \\
 
 $\because \alpha_i - \beta_i \leq C|t_i - t_{i-1}|^\alpha = \dfrac{C(b-a)^\alpha}{n^\alpha}$, \\

 $\phantom{\because}\mu_i - \nu_i \leq C|t_i - t_{i-1}|^\beta = \dfrac{C(b-a)^\beta}{n^\beta}$, \\
 
 $\displaystyle\therefore \sum_{i=1}^n \sigma(I_i) = \sum_{i=1}^n \dfrac{C^2(b-a)^{\alpha+\beta}}{n^{\alpha+\beta}} = \dfrac{C^2(b-a)^{\alpha+\beta}}{n^{\alpha+\beta-1}}$, \\
 
 $\displaystyle\because \alpha+\beta>1, \quad\therefore \lim_{n\to\infty} \dfrac{C^2(b-a)^{\alpha+\beta}}{n^{\alpha+\beta-1}}=0$, \\

$\displaystyle\therefore \forall\,\varepsilon>0,\;\exists N,\;\text{s.t.}\,\forall\,n\geq N,\;\gamma([a,b])\subset\bigcup_{i=1}^n I_i$,且$\displaystyle\sum_{i=1}^n \sigma(I_i)<\varepsilon$, \\

从而 $\gamma([a,b])$是$\mathbb{R}^2$中零面积集,得证.

\begin{ti}
【题6】设 $f\in C^2(\overline{B_1(0)})$,在 $B_1(0)$ 上 $C^2$,设 $f$ 在 $\partial B_1(0)$ 上最小值为 $a$,
$f(0)=b<a$,考虑映射 $\phi:B_1(0)\to\mathbb{R}^n,\phi(x)=\nabla f(x)$,$\Sigma=\{x\in B_1(0)|\nabla^2 f(x)$半正定$\}$,
证明:$\phi(\Sigma)$ 中包含一个半径为 $a-b$ 的开球
\end{ti}

\textbf{证明}:$\displaystyle\because f(0)=b<a=\min_{\partial B_1(0)} f$,且 $\overline{B_1(0)}$ 是有界闭集, $f\in C^2(\overline{B_1(0)})$, \\

$\therefore\,f$在$\overline{B_1(0)}$ 上有最小值,且最小值点在 $B_1(0)$ 取到, \\

即 $\exists\,x^*\in B_1(0),\;\text{s.t.}\, \nabla f(x^*)=0$,且 $\nabla^2 f(x^*)\geq0,\quad\therefore x^*\in\Sigma$, \\

$\because\,0=\phi(x^*),\quad\therefore 0\in\phi(\Sigma)$, \\

对$\forall v\in B_{a-b}(0)$,考虑函数$g(x)=f(x)-v\cdot x,\;\text{则}g\in C^0(\overline{B_1(0)})$,且$g(0)=f(0)=b$, \\

$x\in\partial B_1(0)$ 时, $g(x)\geq a - v\cdot x \geq a - \left\lVert v\right\rVert\left\lVert x\right\rVert   > a - (a-b)=b$, \\

从而$g(x)$在$\overline{B_1(0)}$上的最小值在$B_1(0)$取到, \\

即$\exists\,y^0\in B_1(0)$,$\;\text{s.t.}\,\nabla g(y^0)=0$且$\nabla^2 g(y^0)\geq0$, \\

又$\nabla g(y^0)=\nabla f(y^0)-v$, $\nabla^2 g(y^0)=\nabla^2 f(y^0)\geq0$, \\

$\therefore\,y^0\in\Sigma$$\quad\therefore\,v=\nabla f(y^0)\in\phi(\Sigma)$, \\

从而$B_{a-b}(0)\subset\phi(\Sigma)$,证毕.

\begin{ti}
【题7】设$D\subset\mathbb{R}^2$为平面区域,$u\in C^2(D)$,满足$\dfrac{\partial^2 u}{\partial x^2}+\dfrac{\partial^2 u}{\partial y^2}=0$,
证明:对任何$\overline{B_r(a)}\subset D$,都有$\displaystyle\int_{\partial B_r(a)}\dfrac{\partial u}{\partial \vec{n}}ds=0$,
其中$\vec{n}$为$\partial B_r(a)$的单位外法向量,$\dfrac{\partial u}{\partial \vec{n} }$为方向导数.
\end{ti}

\textbf{证明}:设$a=(a_1,a_2)$,则$\partial B_r(a)$:$\begin{cases}x=a_1+r\cos\theta\\y=a_2+r\sin\theta\end{cases}\quad\theta\in[0,2\pi]$, \\

$\therefore\,\vec{n}=(\cos\theta,\sin\theta)$, $ds=rd\theta$, \\

又$u\in C^2(D)$,$\quad\therefore\dfrac{\partial u}{\partial \vec{n} }=\nabla u\cdot\vec{n}=\dfrac{\partial u}{\partial x}\cos\theta+\dfrac{\partial u}{\partial y}\sin\theta$,\\

$\therefore\displaystyle\int_{\partial B_r(a)}\dfrac{\partial u}{\partial \vec{n} }ds$
$=\displaystyle\int_0^{2\pi}[\dfrac{\partial u}{\partial x}(a_1+r\cos\theta,a_2+r\sin\theta)\cdot\cos\theta+\dfrac{\partial u}{\partial y}(a_1+r\cos\theta,a_2+r\sin\theta)\sin\theta] rd\theta$ \\

$=\displaystyle\int_0^{2\pi}[\dfrac{\partial u}{\partial x}\cdot(r\cos\theta)+\dfrac{\partial u}{\partial y}\cdot(r\sin\theta)]d\theta$
$=\displaystyle\int_0^{2\pi}[\dfrac{\partial u}{\partial x}\cdot y'(\theta)-\dfrac{\partial u}{\partial y}\cdot x'(\theta)]d\theta$ \\

$=\displaystyle\oint_{\partial B_r(a)}\dfrac{\partial u}{\partial x}dy-\dfrac{\partial u}{\partial y}dx$, \\

设$P(x,y)=-\dfrac{\partial u}{\partial y}$,$\;Q(x,y)=\dfrac{\partial u}{\partial x}$,$\quad\therefore\,P,Q\in C^1(D)$, $\dfrac{\partial Q}{\partial x}=\dfrac{\partial^2 u}{\partial x^2}$,$\;\dfrac{\partial P}{\partial y}=-\dfrac{\partial^2 u}{\partial y^2}$, \\

由$Green$公式,$\displaystyle\oint_{\partial B_r(a)}\dfrac{\partial u}{\partial x}dy-\dfrac{\partial u}{\partial y}dx=\iint_{B_r(a)}(\dfrac{\partial^2 u}{\partial x^2}+\dfrac{\partial^2 u}{\partial y^2})dxdy=0$, \\

$\therefore\displaystyle\int_{\partial B_r(a)}\dfrac{\partial u}{\partial \vec{n} }ds=0$,得证.

\begin{ti}
【题8】$\Omega$为有界$C^1$区域,$\displaystyle u\in C^2(\overline\Omega)$,$\displaystyle\Delta u=u$,且$\displaystyle u|_{\partial\Omega}\equiv 0$,
证明:$u\equiv0$
\end{ti}

\textbf{微分学视角}:假设$u$不恒为$0$,存在$x^0\in int\Omega$,$\;\text{s.t.}\,u(x^0)>0$, \\

且$\displaystyle u(x^0)=\max_{\Omega}u$,则$\nabla u(x^0)=0$,且$\nabla^2 u(x^0)\leq 0$, \\

由于$\Delta u = \dfrac{\partial^2 u}{\partial x^2} + \dfrac{\partial^2 u}{\partial y^2} = \text{tr}(\nabla^2 u)$, \\

$\therefore\,\Delta u(x^0)\leq 0$,与$\Delta u(x^0)=u(x^0)$矛盾!  从而$u \equiv 0$,得证. \\

\textbf{积分学视角}:考虑积分$\displaystyle\int_\Omega u^2 dxdy$,则$\displaystyle\int_\Omega u^2 dxdy = \int_\Omega u\Delta u dxdy$, \\

$\because\,\text{div}(u \nabla u) = \text{div}((u\dfrac{\partial u}{\partial x},u\dfrac{\partial u}{\partial y})) = \dfrac{\partial(u\dfrac{\partial u}{\partial x})}{\partial x}+\dfrac{\partial (u\dfrac{\partial u}{\partial y})}{\partial y}$ \\

$=u\dfrac{\partial^2 u}{\partial x^2} + (\dfrac{\partial u}{\partial x})^2 + u\dfrac{\partial^2 u}{\partial y^2} + (\dfrac{\partial u}{\partial y})^2$ $=u\Delta u+||\nabla u||^2$, \\

由$Gauss$公式,$\displaystyle\int_{\Omega} u\Delta u dxdy = \int_\Omega \text{div}(u \nabla u) dxdy - \int_\Omega ||\nabla u||^2 dxdy$ \\

$\displaystyle=\int_{\partial\Omega} u\nabla u \cdot \vec{n} ds - \int_\Omega ||\nabla u||^2 dxdy$, \\

$\displaystyle=-\int_\Omega ||\nabla u||^2 dxdy \leq 0$,而$\int_\Omega u^2 dxdy \geq 0$, \\

从而$\displaystyle\int_\Omega u^2 dxdy = 0 \Rightarrow u \equiv 0$,得证.

\begin{ti}
【题9】计算$\displaystyle\oint_\Gamma\dfrac{xdy - ydx}{ax^2 + 2cxy + by^2}$,
其中$\Gamma:x^2 + y^2 =1$,方向为逆时针且$ab > c^2$
\end{ti}

\textbf{证明}:分母$\displaystyle ax^2+2cxy+by^2=(x,y)M\begin{pmatrix}x\\y\end{pmatrix}$, \\

其中$\displaystyle M=\begin{pmatrix}a&c\\c&b\end{pmatrix}$为实对称矩阵,且$\displaystyle\det M=ab-c^2>0$, \\

可选$\displaystyle P\in SO(2)$,使$\displaystyle P^TMP=\begin{pmatrix}\lambda_1&0\\0&\lambda_2\end{pmatrix}$, \\

令$\displaystyle\begin{pmatrix}x\\y\end{pmatrix}=P\begin{pmatrix}u\\v\end{pmatrix}$,则单位圆及其逆时针方向保持不变,且$\displaystyle x\,dy-y\,dx=u\,dv-v\,du$, \\

由$\displaystyle ab>c^2$可知$\displaystyle a,b$同号,两个特征值也与$\displaystyle a$同号, \\

令$\displaystyle s=\operatorname{sgn}(a)$,则可写$\displaystyle\lambda_i=s\mu_i$,其中$\displaystyle\mu_i>0$,且$\displaystyle\mu_1\mu_2=\lambda_1\lambda_2=ab-c^2$, \\

取$\displaystyle u=\cos\theta,v=\sin\theta$,则$\displaystyle\oint_\Gamma\dfrac{x\,dy-y\,dx}{ax^2+2cxy+by^2}$ \\

$\displaystyle=s\int_0^{2\pi}\dfrac{d\theta}{\mu_1\cos^2\theta+\mu_2\sin^2\theta}$ $\displaystyle=4s\int_0^{\frac{\pi}{2}}\dfrac{d\theta}{\mu_1\cos^2\theta+\mu_2\sin^2\theta}$, \\

令$\displaystyle t=\tan\theta$,则$\displaystyle4s\int_0^{\frac{\pi}{2}}\dfrac{d\theta}{\mu_1\cos^2\theta+\mu_2\sin^2\theta}=4s\int_0^{+\infty}\dfrac{dt}{\mu_1+\mu_2t^2}=\dfrac{2\pi s}{\sqrt{\mu_1\mu_2}},$ \\

故:$\quad\displaystyle\boxed{\oint_\Gamma\dfrac{x\,dy-y\,dx}{ax^2+2cxy+by^2}=\dfrac{2\pi\operatorname{sgn}(a)}{\sqrt{ab-c^2}}}.$

\begin{ti}
【题10】$\Omega\subset \mathbb{R}^3$为有$C^1$边界的有界区域,$\displaystyle u\in C^2(\overline\Omega,\mathbb{R})$,
$\Delta u = 0$,证明:对任意内点$p \in \Omega$,有:
$\displaystyle u(p)=\dfrac{1}{4\pi}\int_{\partial\Omega}\left(\dfrac{\boldsymbol{r}\cdot\boldsymbol{n}}{r^3} u + \dfrac{1}{r}\dfrac{\partial u}{\partial \boldsymbol{n}} \right)d\sigma$,
其中$p = (x_0, y_0, z_0)$,$\boldsymbol{r} = (x - x_0, y - y_0, z - z_0)$,$\boldsymbol{n}$为$\partial\Omega$单位外法向量
\end{ti}

\textbf{证明}:固定$\displaystyle\varepsilon_0>0$使$\displaystyle\overline{B_{\varepsilon_0}(p)}\subset\Omega$,并取$\displaystyle0<\varepsilon<\varepsilon_0$.因为$\displaystyle\dfrac{\partial u}{\partial \boldsymbol{n}}=\nabla u\cdot\boldsymbol{n}$,令$\displaystyle\boldsymbol{X}=\dfrac{\boldsymbol{r}}{r^3}u+\dfrac{1}{r}\nabla u$,考虑$\displaystyle\Omega\setminus B_\varepsilon(p)$, \\

记$\displaystyle I_1 =\int_{\partial\Omega}(\dfrac{\boldsymbol{r} \cdot \boldsymbol{n}}{r^3} u +\dfrac{1}{r}\dfrac{\partial u}{\partial \boldsymbol{n}}) d\sigma$,
$\displaystyle I_2 = \int_{\partial B_\varepsilon(p)} (\dfrac{\boldsymbol{r} \cdot \boldsymbol{n}}{r^3} u +\dfrac{1}{r}\dfrac{\partial u}{\partial \boldsymbol{n}}) d\sigma$, \\

则由$Gauss$公式,$\displaystyle I_1 + I_2 = \int_{\Omega\backslash B_{\varepsilon}(p)}\text{div}\boldsymbol{X} \, dxdydz$,记$\boldsymbol{X} = (P, Q, R)$, \\

$\therefore\dfrac{\partial P}{\partial x} = \dfrac{1}{r}\dfrac{\partial^2 u}{\partial x^2} + (-\dfrac{1}{r^2})\dfrac{x-x_0}{r}\dfrac{\partial u}{\partial x}+\dfrac{x-x_0}{r^3}\dfrac{\partial u}{\partial x} + u\cdot\dfrac{r^2 - 3(x - x_0)^2}{r^5}$ \\

$=\dfrac{1}{r}\dfrac{\partial^2 u}{\partial x^2} + u\cdot\dfrac{r^2 - 3(x - x_0)^2}{r^5}$, \\

同理$\dfrac{\partial Q}{\partial y} = \dfrac{1}{r} \dfrac{\partial^2 u}{\partial y^2} + u\cdot\dfrac{r^2 - 3(y - y_0)^2}{r^5}$,
$\dfrac{\partial R}{\partial z} = \dfrac{1}{r}\dfrac{\partial^2 u}{\partial z^2} + u\cdot\dfrac{r^2 - 3(z - z_0)^2}{r^5}$, \\

$\therefore\,\text{div} \boldsymbol{X} = \dfrac{1}{r}\Delta u + u\cdot\dfrac{3r^2 - 3r^2}{r^5} = 0$, 
$\quad\therefore I_1 + I_2 = 0$, \\

在$\partial B_\varepsilon(p)$上,诱导定向的$\boldsymbol{n} = -\dfrac{\boldsymbol{r}}{r}$,且$r = \varepsilon$, \\

$\displaystyle\therefore\,I_2 = -\int_{\partial B_\varepsilon(p)}(\dfrac{u}{\varepsilon^2}+\dfrac{\nabla u}{\varepsilon}\cdot\dfrac{\boldsymbol{r}}{r} \,) d\sigma$, \\

令$\displaystyle C=\max_{\overline{B_{\varepsilon_0}(p)}}\left\lVert \nabla u\right\rVert <\infty$.对每个$\displaystyle0<\varepsilon<\varepsilon_0$,有 \\

$\displaystyle\left|\int_{\partial B_\varepsilon(p)}\frac{\nabla u}{\varepsilon}\cdot\frac{\boldsymbol{r}}{r}\,d\sigma\right|\leq\frac{C}{\varepsilon}\,4\pi\varepsilon^2=4\pi C\varepsilon\to0,$ \\

又由$\displaystyle u$在$\displaystyle p$处连续,有$\displaystyle\sup_{|x-p|=\varepsilon}|u(x)-u(p)|\to0$, \\

从而$\displaystyle\left|\frac{1}{\varepsilon^2}\int_{\partial B_\varepsilon(p)}u\,d\sigma-4\pi u(p)\right|\leq4\pi\sup_{|x-p|=\varepsilon}|u(x)-u(p)|\to0,$ \\

$\displaystyle\therefore\lim_{\varepsilon\to0}I_2=-4\pi u(p)$,又$\displaystyle I_1+I_2=0$对每个$\displaystyle0<\varepsilon<\varepsilon_0$成立, \\

$\therefore\,I_1 = 4\pi u(p)$,即$\displaystyle u(p) = \dfrac{1}{4\pi}\int_{\partial \Omega} \left( \dfrac{\boldsymbol{r} \cdot \boldsymbol{n}}{r^3} u + \dfrac{1}{r} \dfrac{\partial u}{\partial \boldsymbol{n}} \right) d\sigma$.



\BookEnd

\backmatter
\booknotetitle{参考资料}
\begin{booknotebody}

\begin{enumerate}[leftmargin=2\ccwd,itemsep=0.55em,topsep=0.4em]
    \item 梅加强.\textit{数学分析}[M].第2版.北京:高等教育出版社,2020.ISBN 978-7-04-053344-6.
    \item 谢惠民,恽自求,易法槐,钱定边.\textit{数学分析习题课讲义:上册}[M].第2版.北京:高等教育出版社,2018.ISBN 978-7-04-049851-6.
    \item 谢惠民,恽自求,易法槐,钱定边.\textit{数学分析习题课讲义:下册}[M].第2版.北京:高等教育出版社,2019.ISBN 978-7-04-051152-9.
\end{enumerate}


\end{booknotebody}



\makebackcover

\end{document}
